% Fonctionnement d'un moteur à combustion interne — PCT 3ème — Cours signé OnBuch+
% Programme officiel MINESEC. Compiler avec : tectonic N19-fonctionnement-moteur-combustion-interne.tex
\newcommand{\DOCMATIERE}{PCT}
\newcommand{\DOCNIVEAU}{3ème}
\newcommand{\DOCMODULE}{Module 4 : Technologie et mécanique}
\newcommand{\DOCLECON}{N19}
\newcommand{\DOCTITRE}{Fonctionnement d'un moteur à combustion interne}
% OnBuch+ — préambule « cours signés » ENRICHI (compilé avec tectonic / XeLaTeX)
% Système visuel « Pop » d'OnBuch+ : fond crème (jamais blanc), bordures encre
% épaisses, ombres « dures » décalées, titres Archivo Black en capitales.
% Version MATHÉMATIQUES : graphes (pgfplots/tikz), boîtes pédagogiques
% (définition, propriété, méthode, attention, le savais-tu, exercice résolu,
% exercice, TP, bilan), page de garde et signature OnBuch+ ancrée en bas.
%
% Macros attendues dans main.tex AVANT \input{preamble} :
%   \DOCMATIERE  \DOCNIVEAU  \DOCMODULE  \DOCLECON (numéro)  \DOCTITRE
\documentclass[11pt]{article}

\usepackage{fontspec}
\usepackage[french]{babel}
\usepackage[a4paper,margin=2cm,top=3.4cm,bottom=2.8cm,headheight=1.4cm,footskip=1.3cm]{geometry}
\usepackage{amsmath,amssymb,amsfonts}
\usepackage{mathtools} % \xmapsto, \coloneqq... souvent utilisés par les modèles
\usepackage{cancel} % simplifications barrées (bilans de forces)
% \abs{z} (module, valeur absolue) et \norm{u} (norme), utilisés par les modèles
\providecommand{\abs}{}\let\abs\relax\DeclarePairedDelimiter{\abs}{\lvert}{\rvert}
\providecommand{\norm}{}\let\norm\relax\DeclarePairedDelimiter{\norm}{\lVert}{\rVert}
\usepackage[table]{xcolor} % [table] : fournit aussi \rowcolors (alternance de lignes)
\usepackage{pagecolor}
\usepackage{tikz}
\usetikzlibrary{arrows.meta,positioning,calc,shapes.geometric,shapes.misc,shapes.arrows,decorations.pathmorphing,decorations.pathreplacing,decorations.markings,patterns,shadows,mindmap,trees,fit,backgrounds,matrix,intersections,angles,quotes}
\usepackage{pgfplots}
\usepackage[version=4]{mhchem} % équations nucléaires \ce{^{235}_{92}U}
\usepackage[european,siunitx]{circuitikz} % schémas électriques
\pgfplotsset{compat=1.18}
\usepgfplotslibrary{fillbetween,groupplots} % aires entre courbes (calcul intégral)
\usepackage{enumitem}
\usepackage{fancyhdr}
% [most] charge toutes les bibliothèques tcolorbox (listings, minted, raster,
% documentation...) et gonfle inutilement la mémoire TeX sur un document long
% (~300 boîtes) : on ne charge que ce qui est réellement utilisé.
\usepackage[skins,breakable]{tcolorbox}
\usepackage{graphicx}
\usepackage{colortbl}
\usepackage{booktabs}
\usepackage{tabularx}
\usepackage{multirow}
\usepackage{diagbox} % case d'en-tête diagonale des tableaux à double entrée
% Colonnes extensibles centrée / gauche / droite (C, L, R), souvent utilisées par les modèles
\newcolumntype{C}{>{\centering\arraybackslash}X}
\newcolumntype{L}{>{\raggedright\arraybackslash}X}
\newcolumntype{R}{>{\raggedleft\arraybackslash}X}
\usepackage{array}
\usepackage{multicol}
\usepackage{siunitx}
% parse-numbers=false : accepte aussi \SI{2,5 \times 10^{-3}}{...}
\sisetup{locale=FR,output-decimal-marker={,},group-minimum-digits=4,parse-numbers=false}
\DeclareSIUnit\ohm{\ensuremath{\symup{\Omega}}} % U+2126 absent de la police
\DeclareSIUnit\division{div} % réglages d'oscilloscope (ms/div, V/div)
\DeclareSIUnit\tr{tr} % tours (vitesse de rotation, ex. \SI{1500}{\tr\per\minute})
\DeclareSIUnit\molar{M} % concentration molaire (ex. \SI{3}{\molar}, \SI{1}{\micro\molar})
\DeclareSIUnit\an{an} % année (le modèle confond parfois avec \year, un compteur TeX) — ex. \SI{5}{\kilo\meter\per\an}
\DeclareSIUnit\FCFA{FCFA} % franc CFA (ex. \SI{500}{\FCFA\per\kilo\gram})
\DeclareSIUnit\degreeBrix{\textdegree Brix} % teneur en sucre d'un sirop/jus (réfractométrie)
\DeclareSIUnit\month{mois} % le modèle utilise parfois \month, un compteur TeX, comme unité de durée
\DeclareSIUnit\microliter{\micro\liter} % le modèle écrit parfois \microliter en un seul token
\DeclareSIUnit\million{millions} % dénombrement (ex. \SI{15}{\million\per\mL} spermatozoïdes)
\usepackage{qrcode}
% \degree : commande fréquemment utilisée par les modèles pour un angle ou une
% température en dehors de \SI{}{} (non fournie par les paquets chargés).
\providecommand{\degree}{^{\circ}}
% \qed (fin de démonstration), utilisé par les modèles sans amsthm
\providecommand{\qed}{\hfill$\square$}
% \barre : double barre des valeurs interdites dans un tableau de variations (tkz-tab)
\providecommand{\barre}{\ensuremath{\|}}
% environnement proof (amsthm non chargé) utilisé par les modèles
\ifdefined\proof\else\newenvironment{proof}[1][Démonstration]{\par\noindent\textit{#1.}\ }{\qed\par}\fi
\usepackage{lastpage}
\usepackage{hyperref}
\hypersetup{hidelinks}

% ============================================================
% Palette « Pop » OnBuch+ — mêmes hex que la classe P (plus_ui.dart)
% ============================================================
\definecolor{poporange}{HTML}{F59321}
\definecolor{popdark}{HTML}{E07A0C}
\definecolor{popcream}{HTML}{FFF6E8}
\definecolor{popink}{HTML}{1C1714}
\definecolor{poppurple}{HTML}{7A5AE0}
\definecolor{popgreen}{HTML}{2E9E6A}
\definecolor{poppink}{HTML}{E0457B}
\definecolor{popblue}{HTML}{2F7DE1}
\definecolor{popgold}{HTML}{C98A12}
\definecolor{popmuted}{HTML}{8A7F76}
\definecolor{popline}{HTML}{EADFD2}
\definecolor{popgray}{HTML}{8A7F76}
\colorlet{popgrey}{popgray}
% Alias tolérants (le modèle nomme parfois les couleurs différemment)
\colorlet{popwhite}{white}
\colorlet{popblack}{popink}
\colorlet{popred}{poppink}
\colorlet{popyellow}{popgold}
% Teintes claires (fonds de tableaux/encadrés) — le modèle nomme parfois
% les couleurs avec un suffixe L (ex. popblueL) sans les avoir définies.
\colorlet{poporangeL}{poporange!20}
\colorlet{popdarkL}{popdark!20}
\colorlet{poppurpleL}{poppurple!20}
\colorlet{popgreenL}{popgreen!20}
\colorlet{poppinkL}{poppink!20}
\colorlet{popblueL}{popblue!20}
\colorlet{popgoldL}{popgold!20}
\colorlet{popredL}{popred!20}
\colorlet{popgrayL}{popgray!20}
% Teintes claires pour fonds de tableaux / zones
\colorlet{popblueL}{popblue!10!white}
\colorlet{poporangeL}{poporange!14!white}
\colorlet{popgreenL}{popgreen!12!white}
\colorlet{poppurpleL}{poppurple!12!white}
\colorlet{poppinkL}{poppink!12!white}
\colorlet{popgoldL}{popgold!14!white}

\pagecolor{popcream}

% ============================================================
% Typographie
% ============================================================
\defaultfontfeatures{Ligatures=TeX}
\setmainfont{PlusJakartaSans-Medium.ttf}[Path=../../../fonts/,BoldFont=PlusJakartaSans-Bold.ttf]
% Maths en sans-serif (Fira Math) avec chiffres et lettres droites de Plus
% Jakarta, pour que les formules se fondent dans le texte.
\usepackage[math-style=ISO,bold-style=ISO]{unicode-math}
\setmathfont{FiraMath-Regular.otf}[Path=../../../fonts/]
\setmathfont{PlusJakartaSans-Medium.ttf}[Path=../../../fonts/,range={up/num,up/{Latin,latin},\mathup}]
\usepackage{icomma} % virgule décimale sans espace en mode math
% Caractères mathématiques Unicode absents des polices de texte (Plus Jakarta,
% Archivo Black) : rendus par leur équivalent mathématique, en texte comme en
% formule — sinon ils s'affichent en carré vide (ex. « Arithmétique dans ℤ »).
\usepackage{newunicodechar}
\newunicodechar{ℝ}{\ensuremath{\mathbb{R}}}
\newunicodechar{ℤ}{\ensuremath{\mathbb{Z}}}
\newunicodechar{ℂ}{\ensuremath{\mathbb{C}}}
\newunicodechar{ℕ}{\ensuremath{\mathbb{N}}}
\newunicodechar{ℚ}{\ensuremath{\mathbb{Q}}}
\newunicodechar{𝔻}{\ensuremath{\mathbb{D}}}
\newunicodechar{α}{\ensuremath{\alpha}}
\newunicodechar{β}{\ensuremath{\beta}}
\newunicodechar{γ}{\ensuremath{\gamma}}
\newunicodechar{δ}{\ensuremath{\delta}}
\newunicodechar{ε}{\ensuremath{\varepsilon}}
\newunicodechar{θ}{\ensuremath{\theta}}
\newunicodechar{λ}{\ensuremath{\lambda}}
\newunicodechar{μ}{\ensuremath{\mu}}
\newunicodechar{σ}{\ensuremath{\sigma}}
\newunicodechar{τ}{\ensuremath{\tau}}
\newunicodechar{φ}{\ensuremath{\varphi}}
\newunicodechar{ψ}{\ensuremath{\psi}}
\newunicodechar{ω}{\ensuremath{\omega}}
\newunicodechar{Σ}{\ensuremath{\Sigma}}
% majuscules produites par \MakeUppercase dans les titres (α-aminés -> Α)
\newunicodechar{Α}{\ensuremath{\alpha}}
\newunicodechar{Β}{\ensuremath{\beta}}
\newunicodechar{Γ}{\ensuremath{\Gamma}}
\newunicodechar{Δ}{\ensuremath{\Delta}}
\newunicodechar{Ε}{\ensuremath{\varepsilon}}
\newunicodechar{Θ}{\ensuremath{\Theta}}
\newunicodechar{Λ}{\ensuremath{\Lambda}}
\newunicodechar{Μ}{\ensuremath{\mu}}
\newunicodechar{Τ}{\ensuremath{\tau}}
\newunicodechar{Φ}{\ensuremath{\Phi}}
\newunicodechar{Ψ}{\ensuremath{\Psi}}
\newunicodechar{Ω}{\ensuremath{\Omega}}
\newunicodechar{↦}{\ensuremath{\mapsto}}
\newunicodechar{∈}{\ensuremath{\in}}
\newunicodechar{∉}{\ensuremath{\notin}}
\newunicodechar{∧}{\ensuremath{\wedge}}
\newunicodechar{⇔}{\ensuremath{\Leftrightarrow}}
\newunicodechar{⟺}{\ensuremath{\Longleftrightarrow}}
\newunicodechar{⇒}{\ensuremath{\Rightarrow}}
\newunicodechar{⟹}{\ensuremath{\Longrightarrow}}
\newunicodechar{∘}{\ensuremath{\circ}}
\newunicodechar{∪}{\ensuremath{\cup}}
\newunicodechar{∩}{\ensuremath{\cap}}
\newunicodechar{≡}{\ensuremath{\equiv}}
\newunicodechar{⊂}{\ensuremath{\subset}}
\newunicodechar{⊥}{\ensuremath{\perp}}
\newunicodechar{∃}{\ensuremath{\exists}}
\newunicodechar{∀}{\ensuremath{\forall}}
\newunicodechar{∛}{\ensuremath{\sqrt[3]{\,}}}
\newunicodechar{∜}{\ensuremath{\sqrt[4]{\,}}}
\newunicodechar{⃗}{\ensuremath{{}^{\to}}}
\newunicodechar{ⁿ}{\ifmmode^{n}\else\textsuperscript{\ensuremath{n}}\fi}
\newunicodechar{ˣ}{\ifmmode^{x}\else\textsuperscript{\ensuremath{x}}\fi}
\newunicodechar{ᵇ}{\ifmmode^{b}\else\textsuperscript{\ensuremath{b}}\fi}
\newunicodechar{ʳ}{\ifmmode^{r}\else\textsuperscript{\ensuremath{r}}\fi}
\newunicodechar{ᵘ}{\ifmmode^{u}\else\textsuperscript{\ensuremath{u}}\fi}
\newunicodechar{ʸ}{\ifmmode^{y}\else\textsuperscript{\ensuremath{y}}\fi}
\newunicodechar{ᵃ}{\ifmmode^{a}\else\textsuperscript{\ensuremath{a}}\fi}
\newunicodechar{ᵉ}{\ifmmode^{e}\else\textsuperscript{\ensuremath{e}}\fi}
\newunicodechar{ᵏ}{\ifmmode^{k}\else\textsuperscript{\ensuremath{k}}\fi}
\newunicodechar{ᵖ}{\ifmmode^{p}\else\textsuperscript{\ensuremath{p}}\fi}
\newunicodechar{ᴺ}{\ifmmode^{N}\else\textsuperscript{\ensuremath{N}}\fi}
\newunicodechar{ᶜ}{\ifmmode^{c}\else\textsuperscript{\ensuremath{c}}\fi}
\newunicodechar{ʰ}{\ifmmode^{h}\else\textsuperscript{\ensuremath{h}}\fi}
\newunicodechar{ᵠ}{\ifmmode^{q}\else\textsuperscript{\ensuremath{q}}\fi}
\newunicodechar{ₙ}{\ifmmode_{n}\else\textsubscript{\ensuremath{n}}\fi}
\newunicodechar{ₐ}{\ifmmode_{a}\else\textsubscript{\ensuremath{a}}\fi}
\newunicodechar{ᵢ}{\ifmmode_{i}\else\textsubscript{\ensuremath{i}}\fi}
\newunicodechar{ᵣ}{\ifmmode_{r}\else\textsubscript{\ensuremath{r}}\fi}
\newunicodechar{ₖ}{\ifmmode_{k}\else\textsubscript{\ensuremath{k}}\fi}
\newunicodechar{ᵦ}{\ifmmode_{\beta}\else\textsubscript{\ensuremath{\beta}}\fi}
\newunicodechar{ₓ}{\ifmmode_{x}\else\textsubscript{\ensuremath{x}}\fi}
\newfontfamily\popdisplay{ArchivoBlack-Regular.ttf}[Path=../../../fonts/]
\newfontfamily\poplogo{SpaceGrotesk-Bold.ttf}[Path=../../../fonts/]
\newfontfamily\popsemi{PlusJakartaSans-SemiBold.ttf}[Path=../../../fonts/]
\newfontfamily\popxbold{PlusJakartaSans-ExtraBold.ttf}[Path=../../../fonts/]
\newfontfamily\popxboldls{PlusJakartaSans-ExtraBold.ttf}[Path=../../../fonts/,LetterSpace=3.0]

\setlist[enumerate]{leftmargin=1.7em,itemsep=5pt,topsep=4pt}
\setlist[itemize]{leftmargin=1.7em,itemsep=4pt,topsep=4pt,label={\color{poporange}\textbullet}}
\setlength{\parindent}{0pt}
\setlength{\parskip}{5pt}
\renewcommand{\arraystretch}{1.25}

% pgfplots : style Pop par défaut
\pgfplotsset{
  every axis/.append style={axis line style={popink,line width=1pt},tick style={popink},
    grid style={popline,line width=0.5pt},label style={font=\small},tick label style={font=\footnotesize}},
  popcurve/.style={poporange,line width=1.8pt,no markers,smooth},
}
% Même style disponible dans un \draw TikZ simple (hors pgfplots)
\tikzset{popcurve/.style={poporange,line width=1.8pt}}
\tikzset{popgrid/.style={popline,very thin}}
\colorlet{popcurve}{poporange} % le nom du style est parfois utilisé comme couleur

% ============================================================
% Pill / squiggle
% ============================================================
\newcommand{\poppill}[3]{%
  \tikz[baseline=-0.55ex]\node[draw=popink,line width=1.2pt,rounded corners=2.6pt,%
    fill=#2,inner xsep=6.5pt,inner ysep=3pt]{%
    \popxboldls\fontsize{7}{7}\selectfont\color{#3}\MakeUppercase{#1}};%
}
\newcommand{\popsquiggle}[1]{%
  \tikz[baseline=-0.4ex]{
    \draw[#1,line width=2.6pt,line cap=round]
      plot [smooth] coordinates {(0,0.09) (0.34,0.01) (0.68,0.21) (1.02,0.01) (1.36,0.10)};
  }%
}
% Mot-clé surligné (marqueur orange sous le texte)
\newcommand{\cle}[1]{{\popxbold\color{popdark}#1}}

% ============================================================
% En-tête / pied
% ============================================================
\pagestyle{fancy}
\fancyhf{}
\renewcommand{\headrulewidth}{0pt}
\renewcommand{\footrulewidth}{0pt}
\lhead{\raisebox{-0.15ex}{\poplogo\fontsize{13}{13}\selectfont\color{popink}OnBuch\color{poporange}+}}
\rhead{\poppill{\DOCMATIERE\ \textbullet\ \DOCNIVEAU\ \textbullet\ Leçon \DOCLECON}{white}{popink}}
\lfoot{\popxbold\fontsize{7.6}{7.6}\selectfont\color{popmuted}\MakeUppercase{Cours signé OnBuch+}\ \textbullet\ {\popsemi\DOCTITRE}}
\rfoot{\tikz[baseline=-0.6ex]\node[rounded corners=8.5pt,fill=popink,minimum height=17pt,minimum width=17pt,inner xsep=5pt,inner sep=0]{\popxbold\fontsize{8}{8}\selectfont\color{popcream}\thepage\,/\,\pageref*{LastPage}};}

% ============================================================
% Titres
% ============================================================
\newcommand{\chaptitre}[2]{%
  \begin{tcolorbox}[enhanced,colback=poporange,colframe=popink,boxrule=2.6pt,arc=11pt,
      drop shadow=popink,left=15pt,right=15pt,top=12pt,bottom=11pt,before skip=3pt,after skip=18pt]
    \poppill{#2}{popink}{popcream}\par\vspace{9pt}
    {\raggedright\hyphenpenalty=10000\exhyphenpenalty=10000\popdisplay\fontsize{22}{24}\selectfont\color{popcream}\MakeUppercase{#1}\par}
  \end{tcolorbox}
}
% Section numérotée : pastille orange avec le numéro + titre Archivo
\newcounter{popsec}
\newcommand{\coursec}[1]{%
  \stepcounter{popsec}\setcounter{popsub}{0}%
  \par\vspace{16pt}\noindent
  \begin{minipage}{\linewidth}
  \tikz[baseline=-0.7ex]\node[draw=popink,line width=1.6pt,fill=poporange,rounded corners=4pt,minimum size=22pt,inner sep=0]{\popdisplay\fontsize{12}{12}\selectfont\color{popcream}\thepopsec};%
  \hspace{8pt}{\popdisplay\fontsize{15}{17}\selectfont\color{popink}\MakeUppercase{#1}}\par
  \vspace{4pt}\noindent{\color{poporange}\rule{36pt}{3.6pt}}
  \end{minipage}\par\nopagebreak\vspace{8pt}
}
\newcounter{popsub}
\newcommand{\courssub}[1]{%
  \stepcounter{popsub}%
  \par\vspace{9pt}\noindent{\popxbold\fontsize{11.5}{13}\selectfont\color{popink}{\color{poporange}\thepopsec.\thepopsub}\hspace{6pt}#1}\par\nopagebreak\vspace{4pt}
}

% ============================================================
% Boîtes pédagogiques — même recette : fond blanc, bordure épaisse colorée,
% ombre dure encre, badge plein. Titre optionnel [#1].
% ============================================================
\tcbset{popbase/.style={breakable,enhanced,colback=white,boxrule=2.3pt,arc=9pt,
  shadow={3pt}{-3pt}{0pt}{popink},left=14pt,right=14pt,top=11pt,bottom=11pt,before skip=11pt,after skip=15pt,
  fontupper=\popsemi\fontsize{10.6}{14.5}\selectfont\color{popink},
  fontlower=\popsemi\fontsize{10.6}{14.5}\selectfont\color{popink}}}
% Coche dessinée (le glyphe ✓ n'existe pas dans les polices du document)
\newcommand{\popcheck}[1]{\tikz[baseline=-0.5ex]\draw[#1,line width=1.6pt,line cap=round,line join=round](-0.13,0.0)--(-0.04,-0.09)--(0.13,0.1);}
\providecommand{\coche}{\popcheck{popgreen}}
\providecommand{\resultat}[1]{\par\smallskip\noindent\fcolorbox{popgreen}{popgreenL}{\parbox{\dimexpr\linewidth-2\fboxsep-2\fboxrule\relax}{\textbf{Résultat.} #1}}\par\smallskip}
\newcommand{\popboxhead}[3]{\poppill{#1}{#2}{white}\ifx\relax#3\relax\else\hspace{6pt}{\popxbold\fontsize{10.6}{12}\selectfont\color{popink}#3}\fi\par\vspace{7pt}}

\newtcolorbox{aretenir}[1][]{popbase,colframe=popgreen,before upper={\popboxhead{À retenir}{popgreen}{#1}}}
\newtcolorbox{exemplebox}[1][]{popbase,colframe=poporange,before upper={\popboxhead{Exemple}{poporange}{#1}}}
\newtcolorbox{definition}[1][]{popbase,colframe=popblue,before upper={\popboxhead{Définition}{popblue}{#1}}}
\newtcolorbox{propriete}[1][]{popbase,colframe=poppurple,before upper={\popboxhead{Propriété}{poppurple}{#1}}}
\newtcolorbox{methode}[1][]{popbase,colframe=popgold,before upper={\popboxhead{Méthode}{popgold}{#1}}}
\newtcolorbox{attention}[1][]{popbase,colframe=poppink,before upper={\popboxhead{Attention}{poppink}{#1}}}
\newtcolorbox{experience}[1][]{popbase,colframe=popblue,colback=popblueL,before upper={\popboxhead{Expérience}{popblue}{#1}}}
% « Le savais-tu ? » : carte violette pleine (contexte, culture, Cameroun)
\newtcolorbox{savaistu}[1][]{popbase,colback=poppurple,colframe=popink,
  fontupper=\popsemi\fontsize{10.4}{14.2}\selectfont\color{white},
  before upper={\poppill{Le savais-tu\,?}{white}{poppurple}\ifx\relax#1\relax\else\hspace{6pt}{\popxbold\color{white}#1}\fi\par\vspace{7pt}}}
% Exercice résolu : énoncé en haut, \tcblower puis la solution sur fond crème
\newcounter{popexo}\newcounter{popexr}
\newtcolorbox{exoresolu}[1][]{popbase,colframe=poporange,colbacklower=popcream,
  segmentation style={popink,line width=1.2pt,dashed},
  before upper={\stepcounter{popexr}\popboxhead{Exercice résolu \thepopexr}{poporange}{#1}},
  before lower={\poppill{Solution}{popink}{popcream}\par\vspace{6pt}}}
% Exercice (non résolu en place) — la correction est dans la partie Corrigés
\newtcolorbox{exercice}[1][]{popbase,colframe=popink,
  before upper={\stepcounter{popexo}\popboxhead{Exercice \thepopexo}{popink}{#1}}}
% Corrigé
\newtcolorbox{corrige}[1][]{popbase,colframe=popgreen,colback=popgreenL,
  before upper={\popboxhead{Corrigé}{popgreen}{#1}}}
% Objectifs / prérequis / bilan
\newtcolorbox{objectifs}{popbase,colframe=popink,colback=poporangeL,
  before upper={\popboxhead{Objectifs de la leçon}{popink}{}}}
\newtcolorbox{prerequis}{popbase,colframe=popink,colback=popblueL,
  before upper={\popboxhead{Prérequis}{popblue}{}}}
\newtcolorbox{bilan}{popbase,colframe=popink,colback=white,boxrule=2.8pt,
  before upper={\popboxhead{Fiche bilan}{poporange}{L'essentiel à connaître pour l'examen}}}
% Figure encadrée avec légende
\newtcolorbox{popfigure}[1][]{enhanced,breakable=false,colback=white,colframe=popline,boxrule=1.4pt,arc=8pt,
  left=10pt,right=10pt,top=10pt,bottom=8pt,before skip=10pt,after skip=12pt,halign=center,
  lower separated=false,halign lower=center,
  fontlower=\popsemi\fontsize{9}{11.5}\selectfont\color{popmuted},
  #1}
\newcommand{\legende}[1]{\par\vspace{6pt}{\popsemi\fontsize{9}{11.5}\selectfont\color{popmuted}#1}}

% ============================================================
% Signature OnBuch+
% ============================================================
\newcommand{\poplogoinline}[1]{{\poplogo\fontsize{#1}{#1}\selectfont\color{popink}OnBuch\color{poporange}+}}
\newcommand{\signatureonbuch}{%
  \begin{tcolorbox}[enhanced,colback=popink,colframe=popink,boxrule=2.2pt,arc=10pt,
      drop shadow=poporange,left=14pt,right=14pt,top=10pt,bottom=10pt,before skip=0pt,after skip=0pt]
    \tikz[baseline=-0.7ex]\node[circle,fill=popgreen,draw=popcream,line width=1.3pt,minimum size=22pt,inner sep=0]{\popcheck{white}};%
    \hspace{9pt}\begin{minipage}[c]{0.62\linewidth}
      {\popxbold\fontsize{12}{13}\selectfont\color{popcream}Signé\ \textbullet\ {\poplogo OnBuch\color{poporange}+}}\par\vspace{2pt}
      {\raggedright\popsemi\fontsize{8}{10.5}\selectfont\color{popline}\MakeUppercase{Équipe pédagogique OnBuch+}\\\MakeUppercase{Conforme au programme officiel MINESEC}\par}
    \end{minipage}\hfill
    {\poplogo\fontsize{22}{22}\selectfont\color{popcream}OnBuch\color{poporange}+}
  \end{tcolorbox}%
}

% ============================================================
% Page de garde
% #1 titre · #2 matière · #3 niveau · #4 module · #5 n° leçon · #6 durée ·
% #7 description / objectifs (texte ou itemize)
% ============================================================
\newcommand{\pagegarde}[7]{%
  \thispagestyle{empty}
  \newgeometry{a4paper,margin=1.7cm,top=1.6cm,bottom=1.4cm}%
  \begin{tikzpicture}[remember picture,overlay]
    \fill[poporange!18] ([xshift=-3.2cm,yshift=-3.6cm]current page.north east) circle (4.6cm);
    \fill[poppurple!14] ([xshift=2.4cm,yshift=7.6cm]current page.south west) circle (3.8cm);
  \end{tikzpicture}%
  {\poplogo\fontsize{30}{30}\selectfont\color{popink}OnBuch\color{poporange}+}\hfill
  \raisebox{0.6ex}{\poppill{Cours signé \textbullet\ Premium}{popink}{popcream}}\par
  \vspace{1pt}\popsquiggle{poppurple}\par
  \vspace{8pt}
  \begin{tcolorbox}[enhanced,colback=poporange,colframe=popink,boxrule=2.8pt,arc=13pt,
      drop shadow=popink,left=18pt,right=18pt,top=12pt,bottom=13pt,after skip=10pt]
    \poppill{#2 \textbullet\ #3}{popink}{popcream}\hspace{5pt}\poppill{#4}{white}{popink}\par\vspace{8pt}
    {\popdisplay\fontsize{14}{14}\selectfont\color{popink}LEÇON #5}\par\vspace{4pt}
    {\raggedright\hyphenpenalty=10000\exhyphenpenalty=10000\popdisplay\fontsize{25}{27}\selectfont\color{popcream}\MakeUppercase{#1}\par}
    \vspace{8pt}
    {\popxbold\fontsize{9.5}{9.5}\selectfont\color{popink}\MakeUppercase{Durée conseillée : #6}}
  \end{tcolorbox}
  \begin{tcolorbox}[enhanced,colback=white,colframe=popink,boxrule=2.2pt,arc=10pt,
      drop shadow=popink,left=16pt,right=16pt,top=10pt,bottom=10pt,after skip=8pt]
    \poppill{Dans cette leçon}{poppurple}{white}\par\vspace{5pt}
    \popsemi\fontsize{9.3}{12.6}\selectfont\color{popink}#7
  \end{tcolorbox}
  \noindent\begin{minipage}[t]{0.487\linewidth}
    \begin{tcolorbox}[enhanced,colback=popgreen,colframe=popink,boxrule=2.2pt,arc=10pt,
        drop shadow=popink,left=11pt,right=11pt,top=10pt,bottom=10pt,height=3.1cm,valign=center]
      \tcbox[colback=white,colframe=popink,boxrule=1.3pt,arc=5pt,boxsep=0pt,nobeforeafter,
        left=3pt,right=3pt,top=3pt,bottom=3pt,valign=center]{\qrcode[height=1.9cm]{https://play.google.com/store/apps/details?id=com.luvvix.onbuch&hl=fr}}%
      \hspace{8pt}\begin{minipage}[c]{0.5\linewidth}
        {\popdisplay\fontsize{11}{12.5}\selectfont\color{white}\MakeUppercase{Télécharge OnBuch}}\par\vspace{4pt}
        {\raggedright\popsemi\fontsize{8.4}{11}\selectfont\color{white}Gratuit \textbullet\ Google Play\\Tuteur IA, annales, concours, résultats\par}
      \end{minipage}
    \end{tcolorbox}
  \end{minipage}\hfill
  \begin{minipage}[t]{0.487\linewidth}
    \begin{tcolorbox}[enhanced,colback=poppurple,colframe=popink,boxrule=2.2pt,arc=10pt,
        drop shadow=popink,left=11pt,right=11pt,top=10pt,bottom=10pt,height=3.1cm,valign=center]
      \tikz[baseline=-0.7ex]\node[circle,fill=white,draw=popink,line width=1.3pt,minimum size=1.6cm,inner sep=0]{\popdisplay\fontsize{24}{24}\selectfont\color{poppurple}+};%
      \hspace{8pt}\begin{minipage}[c]{0.6\linewidth}
        {\popdisplay\fontsize{11}{12.5}\selectfont\color{white}\MakeUppercase{Passe à OnBuch+}}\par\vspace{4pt}
        {\raggedright\popsemi\fontsize{8.4}{11}\selectfont\color{white}Tous les cours signés, Tuteur IA illimité, fiches PDF — onglet OnBuch+ de l'app\par}
      \end{minipage}
    \end{tcolorbox}
  \end{minipage}
  \vspace{8pt}
  \popcontenu
  \vfill
  \signatureonbuch
  \restoregeometry
  \newpage
}

% Rangée « Ce cours contient » (page de garde)
\newcommand{\poptile}[3]{% #1 couleur #2 gros texte #3 légende
  \begin{tcolorbox}[enhanced,colback=#1,colframe=popink,boxrule=2pt,arc=9pt,shadow={2.5pt}{-2.5pt}{0pt}{popink},
      left=6pt,right=6pt,top=9pt,bottom=9pt,halign=center,valign=center,height=1.75cm,width=\linewidth,nobeforeafter]
    {\popdisplay\fontsize{12.5}{13}\selectfont\color{white}\MakeUppercase{#2}}\par\vspace{3pt}
    {\centering\popsemi\fontsize{7.8}{9.5}\selectfont\color{white}#3\par}
  \end{tcolorbox}}
\newcommand{\popcontenu}{%
  \par\noindent{\popxboldls\fontsize{7.5}{7.5}\selectfont\color{popmuted}CE COURS SIGNÉ CONTIENT}\par\vspace{7pt}
  \noindent\begin{minipage}[t]{0.235\linewidth}\poptile{popblue}{Cours}{complet, illustré, pas à pas}\end{minipage}\hfill
  \begin{minipage}[t]{0.235\linewidth}\poptile{poporange}{Méthodes}{et exercices résolus}\end{minipage}\hfill
  \begin{minipage}[t]{0.235\linewidth}\poptile{poppink}{Exercices}{type examen + corrigés}\end{minipage}\hfill
  \begin{minipage}[t]{0.235\linewidth}\poptile{popgreen}{Fiche bilan}{l'essentiel à retenir}\end{minipage}\par}

% Sommaire
\newtcolorbox{sommaire}{enhanced,colback=white,colframe=popink,boxrule=2.2pt,arc=10pt,
  drop shadow=popink,left=18pt,right=18pt,top=13pt,bottom=13pt,before skip=6pt,after skip=20pt,
  before upper={\poppill{Sommaire}{poppurple}{white}\par\vspace{9pt}\popsemi\fontsize{10.6}{14.5}\selectfont\color{popink}}}

% Signature de fin de cours — poussée tout en bas de la dernière page
\newcommand{\findecours}{%
  \par\vfill\nopagebreak
  \begin{center}\popsquiggle{poporange}\end{center}\vspace{4pt}
  \signatureonbuch
}
\AtBeginDocument{\def\llbracket{\mathopen{[\mkern-3.5mu[}}\def\rrbracket{\mathclose{]\mkern-3.5mu]}}} % intervalles d'entiers (glyphe ⟦ absent de la police)
% Glyphes absents de Fira Math : redéfinis à partir de glyphes disponibles
\AtBeginDocument{%
  \def\setminus{\mathbin{\backslash}}\let\smallsetminus\setminus
  \def\vdots{\mathord{\vbox{\baselineskip3.2pt\lineskiplimit0pt\kern4pt\hbox{.}\hbox{.}\hbox{.}}}}%
  \def\ddots{\mathinner{\mkern1mu\raise6.5pt\hbox{.}\mkern2mu\raise3.6pt\hbox{.}\mkern2mu\raise0.7pt\hbox{.}\mkern1mu}}%
  \def\triangle{\mathord{\scalebox{0.9}{$\symup{\Delta}$}}}%
  \def\nabla{\mathord{\raisebox{\depth}{\scalebox{1}[-1]{$\symup{\Delta}$}}}}%
  \def\checkmark{\popcheck{popdark}}%
  \def\star{\mathbin{\ast}}\def\bigstar{\mathord{\ast}}%
  \def\diamond{\mathbin{\scalebox{0.6}{$\Diamond$}}}%
  \def\bigcup{\mathop{\mathchoice{\raisebox{-0.3em}{\scalebox{1.7}{$\cup$}}}{\scalebox{1.25}{$\cup$}}{\cup}{\cup}}\displaylimits}%
  \def\bigcap{\mathop{\mathchoice{\raisebox{-0.3em}{\scalebox{1.7}{$\cap$}}}{\scalebox{1.25}{$\cap$}}{\cap}{\cap}}\displaylimits}%
  \def\lesseqqgtr{\mathrel{\lessgtr}}\def\bigoplus{\mathop{\oplus}\displaylimits}%
}
\providecommand{\ssi}{si et seulement si} % abréviation fréquente
\AtBeginDocument{\providecommand{\wideparen}{\overparen}} % arc de cercle

%% FIGFIX-BEGIN v3 — correctifs globaux des figures (docs/onbuch-generation/tools/figfix)
\usepackage{adjustbox}\usepackage{etoolbox}
% toute figure TikZ/pgfplots plus large que la ligne est réduite à la largeur disponible
\BeforeBeginEnvironment{tikzpicture}{\begin{adjustbox}{max width=\linewidth}}
\AfterEndEnvironment{tikzpicture}{\end{adjustbox}}
% légendes pgfplots lisibles par-dessus les courbes
\pgfplotsset{every axis/.append style={legend style={fill=white,fill opacity=0.92,text opacity=1,draw=black!15,inner sep=3pt}}}
% étiquettes d'arêtes (syntaxe edge["..."]) avec fond blanc
\tikzset{every edge quotes/.append style={fill=white,inner sep=1.2pt}}
% bulles de cartes mentales : texte à la ligne dans la bulle, bulle un peu plus grande
\tikzset{every concept/.append style={text width=2.0cm,align=center,minimum size=2.3cm,inner sep=2pt}}
%% FIGFIX-END
\begin{document}

\pagegarde{Fonctionnement d'un moteur à combustion interne}{PCT}{3ème}{Module 4 : Technologie et mécanique}{N19}{2 séances (1 h chacune)}{Cette leçon présente les éléments constitutifs d'un moteur à combustion interne et leurs rôles, en détaillant le cycle à quatre temps (admission, compression, détente, échappement). Elle établit ensuite la différence entre moteur à essence et moteur diesel, et sensibilise aux consignes de sécurité liées à l'usage des carburants et des moteurs.
\vspace{4pt}
\begin{multicols}{2}\small
\begin{itemize}
\item À la fin de cette leçon, je sais définir un moteur à combustion interne et citer ses éléments essentiels.
\item Je sais identifier sur un schéma le cylindre, le piston, la bielle, le vilebrequin, les soupapes et la bougie (ou l'injecteur) et donner le rôle de chacun.
\item Je sais décrire les quatre temps du cycle moteur dans l'ordre et expliquer ce qui se passe à chaque temps.
\item Je sais interpréter un schéma du cycle à quatre temps et repérer le temps moteur.
\item Je sais distinguer un moteur à essence d'un moteur diesel (allumage, carburant, organes) et justifier la différence.
\item Je sais expliquer comment le mouvement alternatif du piston est transformé en mouvement de rotation.
\end{itemize}\end{multicols}}

\begin{sommaire}
\begin{multicols}{2}
\begin{enumerate}
\item Le moteur à combustion interne : définition et principe
\item Les éléments d'un moteur et leurs rôles
\item Le cycle à quatre temps
\item Moteur à essence et moteur diesel : différences
\item Sécurité, entretien et environnement
\item Activité d'intégration
\item Méthodes et exercices résolus
\item Exercices
\item Corrigés des exercices
\item Fiche bilan
\end{enumerate}
\end{multicols}
\end{sommaire}

\begin{objectifs}
\begin{itemize}
\item À la fin de cette leçon, je sais définir un moteur à combustion interne et citer ses éléments essentiels.
\item Je sais identifier sur un schéma le cylindre, le piston, la bielle, le vilebrequin, les soupapes et la bougie (ou l'injecteur) et donner le rôle de chacun.
\item Je sais décrire les quatre temps du cycle moteur dans l'ordre et expliquer ce qui se passe à chaque temps.
\item Je sais interpréter un schéma du cycle à quatre temps et repérer le temps moteur.
\item Je sais distinguer un moteur à essence d'un moteur diesel (allumage, carburant, organes) et justifier la différence.
\item Je sais expliquer comment le mouvement alternatif du piston est transformé en mouvement de rotation.
\item Je sais appliquer les consignes de sécurité liées aux carburants et aux gaz d'échappement dans la vie courante.
\item Je sais réaliser un calcul simple lié au fonctionnement du moteur (nombre de tours, cylindrée).
\end{itemize}
\end{objectifs}

\begin{prerequis}
\begin{itemize}
\item Notion d'énergie et de transformation d'énergie (énergie chimique, thermique, mécanique)
\item Combustion d'un combustible : combustible, comburant (dioxygène), énergie d'activation
\item Notions de mouvement : translation et rotation (mécanique, classe de 4ème)
\item Notion de pression d'un gaz et effet de la température sur un gaz
\item Lecture et interprétation de schémas techniques simples
\end{itemize}
\end{prerequis}

\newpage

\coursec{Le moteur à combustion interne : définition et principe}

À Douala comme à Yaoundé, les motos-taxis « benskin » et les taxis jaunes fonctionnent toute la journée grâce à un moteur à combustion interne. Pourtant, certains véhicules roulent à l'essence, d'autres au gasoil (diesel), et le chauffeur sait qu'il ne faut jamais se tromper de carburant à la station Total. Pourquoi le moteur « tourne-t-il » par saccades régulières ? Que se passe-t-il exactement à l'intérieur du cylindre quand on tourne la clé ? Cette leçon ouvre le capot pour comprendre comment l'explosion d'un carburant devient le mouvement de nos véhicules.

\textbf{Problématique :} comment une simple étincelle et un peu d'essence peuvent-ils mettre en mouvement une moto de plus de \SI{150}{kg} ? Pour répondre à cette question, nous allons d'abord définir ce qu'est un moteur, puis suivre pas à pas le chemin de l'énergie, du réservoir jusqu'à la roue.

\courssub{Qu'est-ce qu'un moteur ?}

Observe autour de toi : le ventilateur de la maison tourne grâce à l'électricité d'ENEO, le moulin à maïs du quartier tourne grâce au gasoil, et la pirogue motorisée du Wouri avance grâce à l'essence. Tous ces appareils ont un point commun : ils possèdent un \cle{moteur}, c'est-à-dire un dispositif qui produit du mouvement.

Mais d'où vient ce mouvement ? Le ventilateur ne tourne pas « tout seul » : il reçoit de l'énergie électrique. Le moteur de la moto ne tourne pas tout seul non plus : il reçoit de l'essence. Dans les deux cas, le moteur reçoit une certaine forme d'énergie et la \emph{transforme} en mouvement.

\begin{definition}[Moteur]
Un \cle{moteur} est un dispositif qui transforme une forme d'énergie (électrique, chimique, thermique\ldots) en \cle{énergie mécanique}, c'est-à-dire en mouvement.
\end{definition}

\begin{exemplebox}[Des moteurs dans la vie courante]
\begin{itemize}
\item Le \textbf{moteur électrique} du ventilateur transforme l'énergie électrique en énergie mécanique (rotation des pales).
\item Le \textbf{moteur à essence} d'une moto-taxi transforme l'énergie chimique de l'essence en énergie mécanique (rotation de la roue arrière).
\item Le \textbf{moteur diesel} d'un camion de transport transforme l'énergie chimique du gasoil en énergie mécanique.
\end{itemize}
\end{exemplebox}

\begin{attention}[Le moteur ne « crée » pas d'énergie]
Erreur fréquente : croire que le moteur \emph{fabrique} de l'énergie. C'est faux ! Un moteur ne fait que \textbf{transformer} une énergie qu'il reçoit. De plus, la transformation n'est jamais parfaite : une grande partie de l'énergie du carburant est \emph{perdue sous forme de chaleur} (c'est pour cela qu'un moteur devient brûlant et qu'il faut un système de refroidissement). On dit que le moteur a un \cle{rendement} inférieur à 1 : seule une fraction de l'énergie du carburant devient du mouvement utile.
\end{attention}

\courssub{Moteur à combustion interne : une définition à bien comprendre}

Parmi tous les moteurs, intéressons-nous à ceux qui brûlent un \cle{carburant} (essence, gasoil, pétrole lampant). La \cle{combustion} est une réaction chimique entre un combustible (le carburant) et un comburant (le dioxygène de l'air), qui dégage beaucoup de chaleur. Mais où cette combustion a-t-elle lieu ?

\begin{itemize}
\item Dans la vieille \textbf{machine à vapeur} des locomotives d'autrefois, on brûlait du bois ou du charbon \emph{à l'extérieur} de la machine, sous une chaudière, pour chauffer de l'eau : la combustion était \textbf{externe}.
\item Dans le moteur de ta moto, l'essence brûle \emph{à l'intérieur} du moteur lui-même, dans une chambre appelée \cle{cylindre} : la combustion est \textbf{interne}.
\end{itemize}

\begin{definition}[Moteur à combustion interne]
Un \cle{moteur à combustion interne} est une machine qui transforme l'énergie chimique d'un carburant en énergie mécanique, la combustion du carburant se déroulant \textbf{à l'intérieur} du moteur, dans le cylindre.
\end{definition}

\begin{aretenir}[Les deux idées clés]
\begin{enumerate}
\item Le carburant brûle \textbf{dans} le moteur (d'où le mot « interne »).
\item La chaleur dégagée par la combustion pousse une pièce mobile, le \cle{piston}, qui produit le mouvement.
\end{enumerate}
\end{aretenir}

\courssub{Expérience : l'explosion qui pousse}

Comment être sûr que la combustion d'un mélange air-essence peut produire un mouvement ? Réalisons une expérience simple, mais spectaculaire.

\begin{experience}[La boîte de conserve qui saute]
\textbf{Matériel :} une boîte de conserve vide percée d'un petit trou sur le côté, quelques gouttes d'essence, une bougie allumée ou une allumette longue, un couvercle léger posé sur la boîte (ou une seringue dont le piston est libre).

\textbf{Protocole :}
\begin{enumerate}
\item Verser quelques gouttes d'essence dans la boîte et agiter pour que les vapeurs se mélangent à l'air.
\item Poser le couvercle léger sur l'ouverture.
\item Approcher la flamme du petit trou latéral.
\end{enumerate}

\textbf{Observations :} on entend une détonation et le couvercle est violemment projeté vers le haut. (Avec la seringue, le piston est repoussé brusquement.)

\textbf{Interprétation :} le mélange air-vapeurs d'essence s'enflamme \emph{à l'intérieur} de la boîte. La combustion dégage une grande quantité de chaleur : les gaz se dilatent brusquement et la pression augmente très fort dans le récipient fermé. Cette poussée des gaz chauds est capable de déplacer le couvercle, c'est-à-dire de produire un \textbf{mouvement}.

\textbf{Conclusion :} la combustion d'un mélange air-carburant dans une enceinte fermée transforme l'énergie chimique en énergie mécanique. C'est exactement ce qui se passe dans le cylindre d'un moteur : la boîte joue le rôle du cylindre, et le couvercle celui du piston.
\end{experience}

\begin{attention}[Sécurité]
Cette expérience est réalisée par le professeur, avec de très petites quantités d'essence, loin de tout bidon de carburant. L'essence et ses vapeurs sont extrêmement inflammables : c'est pour cette raison qu'il est interdit de fumer ou de téléphoner près des pompes à la station-service.
\end{attention}

\courssub{Les conditions de la combustion : le triangle du feu}

L'expérience précédente échoue si l'on retire l'un des trois ingrédients. Sans essence, rien ne brûle ; sans air, la combustion s'arrête ; sans flamme, le mélange ne s'enflamme pas tout seul (à température ordinaire). La combustion exige donc la réunion de trois conditions, que l'on représente par le \cle{triangle du feu}.

\begin{definition}[Triangle du feu]
La combustion n'est possible que si trois éléments sont réunis :
\begin{enumerate}
\item un \cle{combustible} : la substance qui brûle (essence, gasoil, bois, gaz domestique\ldots) ;
\item un \cle{comburant} : le dioxygène de l'air, qui entretient la combustion ;
\item une \cle{source de chaleur} (ou énergie d'activation) : étincelle ou flamme qui déclenche la réaction.
\end{enumerate}
Si l'un des trois manque, la combustion est impossible ou s'arrête : c'est le principe des extincteurs.
\end{definition}

\begin{exemplebox}[Le triangle du feu dans le moteur]
Dans le moteur à essence, les trois éléments sont réunis à chaque instant : l'essence (combustible) est injectée, l'air (comburant) est aspiré, et la \textbf{bougie} produit l'étincelle au bon moment. Si la bougie est défectueuse, le moteur ne démarre pas : c'est la panne classique que le mécanicien du quartier vérifie en premier !
\end{exemplebox}

\courssub{La chaîne énergétique du moteur}

Suivons maintenant le chemin de l'énergie, du réservoir d'essence jusqu'à la roue de la moto. Trois étapes se succèdent :

\begin{enumerate}
\item L'essence possède de l'\textbf{énergie chimique}, stockée dans ses molécules (hydrocarbures).
\item La combustion libère cette énergie sous forme d'\textbf{énergie thermique} (chaleur) : les gaz brûlés, très chauds, se dilatent et leur pression augmente brutalement.
\item La poussée des gaz déplace le piston : l'énergie devient \textbf{énergie mécanique}, d'abord un mouvement de va-et-vient du piston, transformé ensuite en mouvement de rotation de l'arbre moteur, puis des roues.
\end{enumerate}

\begin{popfigure}
\begin{center}
\begin{tikzpicture}[
bloc/.style={rectangle, rounded corners=3pt, draw=popdark, thick, fill=popblueL, minimum height=1.1cm, text width=3.1cm, align=center, font=\small},
fleche/.style={-{Stealth[length=3mm]}, very thick, poporange}]
\node[bloc, align=center] (chim) at (0,0){\textbf{Énergie chimique}\\ du carburant (essence)};
\node[bloc, fill=poporangeL, align=center] (therm) at (5.2,0){\textbf{Énergie thermique}\\ (combustion dans le cylindre)};
\node[bloc, fill=popgreenL, align=center] (meca) at (10.4,0){\textbf{Énergie mécanique}\\ (mouvement du piston et des roues)};
\draw[fleche] (chim) -- (therm) node[midway, above, font=\small\itshape, popdark]{combustion};
\draw[fleche] (therm) -- (meca) node[midway, above, font=\small\itshape, popdark]{poussée des gaz};
\node[below=0.35cm of therm, font=\small\itshape, popmuted] {une partie de la chaleur est perdue (échappement, refroidissement)};
\end{tikzpicture}
\end{center}
\legende{Chaîne énergétique d'un moteur à combustion interne : l'énergie chimique du carburant devient énergie thermique, puis énergie mécanique. Une partie de l'énergie est perdue en chaleur.}
\end{popfigure}

\begin{popfigure}
\begin{center}
\begin{tikzpicture}[scale=1]
% cylindre
\draw[thick, popdark] (0,0) -- (0,4.2);
\draw[thick, popdark] (3,0) -- (3,4.2);
\draw[thick, popdark] (0,4.2) -- (3,4.2);
% bougie
\draw[thick, popdark, fill=popgoldL] (1.35,4.2) rectangle (1.65,4.7);
\draw[poporange, thick] (1.5,4.2) -- (1.5,3.95);
\node[font=\small, popdark] at (4.6,4.45) {bougie (étincelle)};
\draw[-{Stealth}, popmuted] (3.4,4.4) -- (1.7,4.45);
% gaz brûlés
\fill[poporange!35] (0.06,2.6) rectangle (2.94,4.14);
\foreach \x/\y in {0.5/3.0, 1.2/3.6, 2.0/3.1, 2.5/3.7, 0.9/2.8, 1.7/3.9, 2.4/2.9, 0.4/3.8}
  \fill[poporange] (\x,\y) circle (0.05);
\node[font=\small, popdark] at (1.5,3.3) {gaz brûlés};
\node[font=\small, popdark] at (1.5,2.85) {très chauds};
% piston
\draw[thick, popdark, fill=popblueL] (0.06,1.7) rectangle (2.94,2.6);
\node[font=\small\bfseries, popdark] at (1.5,2.15) {PISTON};
% force
\draw[-{Stealth[length=4mm]}, very thick, poppink] (1.5,1.55) -- (1.5,0.45);
\node[font=\small, poppink, right] at (1.6,1.0) {poussée des gaz};
% bielle
\draw[thick, popdark] (1.5,1.7) -- (1.9,0.3);
\draw[thick, popdark, fill=popgreenL] (1.9,0.3) circle (0.35);
\node[font=\small, popdark] at (4.4,0.3) {arbre moteur (rotation)};
\draw[-{Stealth}, popmuted] (3.3,0.35) -- (2.3,0.32);
% légende cylindre
\node[font=\small, popdark] at (-1.6,3.5) {cylindre};
\draw[-{Stealth}, popmuted] (-0.9,3.45) -- (-0.05,3.45);
\end{tikzpicture}
\end{center}
\legende{Principe du moteur à combustion interne : l'étincelle de la bougie enflamme le mélange air-essence ; les gaz brûlés se dilatent et poussent le piston vers le bas ; ce mouvement est transmis à l'arbre moteur qui entraîne les roues.}
\end{popfigure}

\begin{methode}[Analyser le fonctionnement d'une machine]
Pour analyser le fonctionnement de n'importe quelle machine (moteur, groupe électrogène, ventilateur, pompe), procède toujours en trois étapes :
\begin{enumerate}
\item Identifier la \textbf{forme d'énergie d'entrée} (ce que la machine reçoit : carburant, électricité\ldots).
\item Décrire les \textbf{transformations} successives de cette énergie à l'intérieur de la machine.
\item Identifier la \textbf{forme d'énergie utile en sortie} (ce que la machine produit : mouvement, électricité\ldots), sans oublier les pertes (chaleur, frottements).
\end{enumerate}
\end{methode}

\begin{exoresolu}[Chaîne énergétique d'une moto-taxi]
\textbf{Énoncé.} Un « benskin » roule à l'essence. 1) Quelle forme d'énergie le moteur reçoit-il ? 2) Quelle forme d'énergie produit-il pour faire avancer la moto ? 3) Pourquoi le moteur est-il chaud après un long trajet ?
\tcblower
\textbf{Solution détaillée.}
\begin{enumerate}
\item Le moteur reçoit de l'\textbf{énergie chimique}, stockée dans l'essence.
\item Il produit de l'\textbf{énergie mécanique} : la combustion chauffe les gaz, qui poussent le piston ; le mouvement du piston devient la rotation de la roue.
\item La transformation n'est pas totale : une partie de l'énergie thermique de la combustion n'est pas convertie en mouvement et chauffe le moteur. C'est une \textbf{perte} inévitable, d'où la nécessité du refroidissement.
\end{enumerate}
\end{exoresolu}

\courssub{Des applications partout au Cameroun}

Les moteurs à combustion interne sont partout autour de nous :

\begin{itemize}
\item les \textbf{motos} et \textbf{voitures} (moteurs à essence ou diesel) ;
\item les \textbf{groupes électrogènes}, qui produisent de l'électricité pendant les coupures ;
\item les \textbf{pirogues motorisées} sur le Wouri, la Sanaga ou le fleuve Logone ;
\item les \textbf{pompes à eau} et les \textbf{moulins à maïs} des villages ;
\item les \textbf{tondeuses}, \textbf{motoculteurs} et engins de chantier.
\end{itemize}

\begin{savaistu}[Le groupe électrogène, un moteur caché]
Pendant les délestages à Yaoundé ou à Douala, les boutiques et les familles allument le groupe électrogène. Sais-tu ce qu'il contient ? Un \textbf{petit moteur à combustion interne} (à essence ou à gasoil) couplé à un \textbf{alternateur}. Le moteur transforme l'énergie chimique du carburant en énergie mécanique (rotation), puis l'alternateur transforme cette rotation en énergie électrique. La chaîne énergétique complète est donc : énergie chimique $\rightarrow$ énergie thermique $\rightarrow$ énergie mécanique $\rightarrow$ énergie électrique !
\end{savaistu}

\begin{aretenir}[L'essentiel de la section]
\begin{itemize}
\item Un \textbf{moteur} transforme une forme d'énergie en \textbf{énergie mécanique} (mouvement).
\item Dans un \textbf{moteur à combustion interne}, le carburant brûle \emph{à l'intérieur} du cylindre (contrairement à la machine à vapeur, à combustion externe).
\item La combustion exige le \textbf{triangle du feu} : combustible + comburant (air) + source de chaleur (étincelle).
\item Chaîne énergétique : énergie \textbf{chimique} $\rightarrow$ énergie \textbf{thermique} $\rightarrow$ énergie \textbf{mécanique}, avec des pertes en chaleur.
\end{itemize}
\end{aretenir}

\begin{exercice}[Pour t'entraîner]
1) Recopie et complète : « Un moteur à combustion interne transforme l'énergie \ldots\ du carburant en énergie \ldots\ ; la combustion a lieu \ldots\ du moteur. »
2) Cite les trois éléments du triangle du feu et indique celui qui manque lorsqu'on étouffe un feu avec une couverture.
3) Un chauffeur de taxi affirme : « Mon moteur crée de l'énergie grâce à l'essence. » Corrige cette phrase en une phrase scientifiquement exacte.
\end{exercice}

\begin{corrige}[Exercice : Pour t'entraîner]
1) « Un moteur à combustion interne transforme l'énergie \textbf{chimique} du carburant en énergie \textbf{mécanique} ; la combustion a lieu \textbf{à l'intérieur} du moteur. »
2) Combustible, comburant (dioxygène de l'air), source de chaleur. La couverture prive le feu de \textbf{comburant} : la combustion s'arrête.
3) Phrase corrigée : « Mon moteur \textbf{transforme} l'énergie chimique de l'essence en énergie mécanique ; une partie de cette énergie est perdue sous forme de chaleur. » Un moteur ne crée jamais d'énergie.
\end{corrige}

\coursec{Les éléments d'un moteur et leurs rôles}

Dans la section précédente, nous avons découvert qu'un moteur à combustion interne transforme l'énergie chimique d'un carburant (essence ou gasoil) en énergie mécanique, grâce à une combustion qui se produit \emph{à l'intérieur} de la machine. Mais comment cette explosion contrôlée parvient-elle à faire tourner les roues de la moto-taxi qui t'amène au collège ? Pour le comprendre, il faut ouvrir le capot et observer les pièces qui travaillent ensemble. C'est exactement ce que fait le mécanicien du quartier lorsqu'il démonte le moteur d'une moto : il retrouve toujours les mêmes organes essentiels. Découvrons-les un à un.

\courssub{Le cylindre et le piston : le c\oe ur du moteur}

\textbf{Observation préalable.} Lorsque tu gonfles un pneu avec une pompe à vélo, tu pousses une pièce mobile (le piston de la pompe) dans un tube (le corps de la pompe) : l'air est comprimé et chassé vers le pneu. Le moteur fonctionne sur le même principe, mais inversé : ce sont les gaz brûlés qui poussent le piston.

\begin{definition}[Le cylindre]
Le \cle{cylindre} est une cavité cylindrique creusée dans le bloc moteur. C'est la \textbf{chambre dans laquelle se déplace le piston} et \textbf{où se produit la combustion} du mélange air-carburant. Sa partie supérieure, fermée par la \textbf{culasse}, constitue la chambre de combustion proprement dite.
\end{definition}

\begin{definition}[Le piston]
Le \cle{piston} est une pièce mobile, de forme cylindrique, qui \textbf{coulisse à l'intérieur du cylindre}. Il joue un double rôle :
\begin{itemize}
\item il \textbf{comprime} le mélange air-carburant lorsqu'il monte ;
\item il \textbf{reçoit la poussée} des gaz brûlés lors de la combustion et la transmet à la bielle lorsqu'il descend.
\end{itemize}
Des segments (anneaux métalliques) assurent l'étanchéité entre le piston et la paroi du cylindre.
\end{definition}

\begin{attention}[Étanchéité du piston]
Le piston n'est pas collé au cylindre : il doit coulisser librement, mais sans laisser fuir les gaz. Ce sont les \textbf{segments} qui jouent le rôle de joint. Un piston usé laisse passer les gaz brûlés : le moteur perd de la puissance et fume bleu (il brûle de l'huile). C'est une panne classique des vieilles motos.
\end{attention}

\textbf{Repères utiles : PMH, PMB et cylindrée.} Au cours de son mouvement de va-et-vient, le piston atteint deux positions extrêmes :

\begin{definition}[Point mort haut et point mort bas]
\begin{itemize}
\item Le \cle{point mort haut} (\textbf{PMH}) est la position la plus haute atteinte par le piston dans le cylindre : le volume au-dessus du piston est alors \emph{minimal}.
\item Le \cle{point mort bas} (\textbf{PMB}) est la position la plus basse atteinte par le piston : le volume au-dessus du piston est alors \emph{maximal}.
\end{itemize}
On dit « point mort » car, à ces deux positions, le piston s'arrête un instant avant de repartir en sens inverse.
\end{definition}

\begin{definition}[La cylindrée]
La \cle{cylindrée} est le \textbf{volume balayé par le piston entre le point mort haut et le point mort bas}. Elle s'exprime en centimètres cubes (cm$^3$). Pour un moteur à plusieurs cylindres, la cylindrée totale est la somme des cylindrées de chaque cylindre.
\end{definition}

\begin{exemplebox}[Des cylindrées de la vie courante]
\begin{itemize}
\item Une moto « 100 » ou « 125 » possède une cylindrée d'environ \SI{100}{cm^3} ou \SI{125}{cm^3} : c'est le volume balayé par son unique piston.
\item Une voiture annoncée « 1\,300 cm$^3$ » possède le plus souvent \textbf{4 cylindres} de $\dfrac{1300}{4} = \SI{325}{cm^3}$ chacun.
\item Plus la cylindrée est grande, plus le moteur peut brûler de carburant à chaque cycle, et plus il est puissant (mais aussi gourmand en carburant !).
\end{itemize}
\end{exemplebox}

\courssub{Les soupapes : les portes de la chambre de combustion}

\textbf{Intuition.} Pour que la combustion se renouvelle sans cesse, il faut faire entrer le mélange frais (air + carburant) et évacuer les gaz brûlés. Il faut donc des « portes » qui s'ouvrent et se ferment au bon moment : ce sont les soupapes.

\begin{definition}[Les soupapes]
Les \cle{soupapes} sont des organes en forme de champignon, placés dans la culasse, qui \textbf{s'ouvrent et se ferment} pour commander les échanges gazeux :
\begin{itemize}
\item la \textbf{soupape d'admission} s'ouvre pour laisser \emph{entrer} le mélange air-carburant (moteur à essence) ou l'air seul (moteur diesel) ;
\item la \textbf{soupape d'échappement} s'ouvre pour laisser \emph{sortir} les gaz brûlés vers le pot d'échappement.
\end{itemize}
Pendant la compression et la combustion, les deux soupapes sont fermées : la chambre de combustion est alors parfaitement close. Leur ouverture est commandée par l'\textbf{arbre à cames}, synchronisé avec le vilebrequin.
\end{definition}

\courssub{La bougie et l'injecteur : deux organes à ne pas confondre}

\begin{definition}[La bougie d'allumage]
Dans un \textbf{moteur à essence}, la \cle{bougie d'allumage} est une pièce vissée dans la culasse, dont l'extrémité plonge dans la chambre de combustion. Elle produit une \textbf{étincelle électrique} entre ses deux électrodes : c'est cette étincelle qui \textbf{déclenche la combustion} du mélange air-essence comprimé.
\end{definition}

\begin{definition}[L'injecteur]
Dans un \textbf{moteur diesel}, il n'y a pas de bougie d'allumage. L'\cle{injecteur} est un organe qui \textbf{pulvérise le gasoil} (en fines gouttelettes, comme un vaporisateur de parfum) dans l'air fortement comprimé et donc très chaud. Le gasoil s'enflamme alors \emph{spontanément}, sans étincelle.
\end{definition}

\begin{attention}[Bougie et injecteur : le piège classique du BEPC]
Ne confonds jamais ces deux organes :
\begin{itemize}
\item la \textbf{bougie} sert à \emph{allumer} le mélange : elle existe dans le moteur \textbf{à essence} ;
\item l'\textbf{injecteur} sert à \emph{amener le carburant} : il caractérise le moteur \textbf{diesel}.
\end{itemize}
Dire « la bougie pulvérise l'essence » ou « l'injecteur produit l'étincelle » est une erreur grave. Retiens : \textbf{bougie = étincelle (essence)} ; \textbf{injecteur = pulvérisation du gasoil (diesel)}.
\end{attention}

\courssub{Le carburateur et le système d'injection : la préparation du mélange}

\begin{definition}[Carburateur et système d'injection]
Dans un \textbf{moteur à essence}, le mélange air-carburant doit être préparé avant d'entrer dans le cylindre :
\begin{itemize}
\item le \cle{carburateur} (moteurs anciens, motos) est un dispositif qui \textbf{mélange l'air et l'essence} dans de bonnes proportions, par aspiration ;
\item le \cle{système d'injection} (véhicules modernes) \textbf{pulvérise l'essence} directement dans l'air admis, de façon électronique et précise.
\end{itemize}
Dans les deux cas, le rôle est le même : \textbf{préparer le mélange air-carburant} avant son admission dans le cylindre.
\end{definition}

\courssub{La bielle et le vilebrequin : du va-et-vient à la rotation}

\textbf{Le problème à résoudre.} Le piston se déplace en ligne droite, en va-et-vient (on dit : mouvement \textbf{rectiligne alternatif}). Or les roues du véhicule doivent \emph{tourner}. Comment transformer un aller-retour en rotation ? Observe le pédalier d'un vélo : tes jambes montent et descendent, mais la pédale et le plateau tournent. Le moteur utilise exactement le même artifice.

\begin{definition}[La bielle]
La \cle{bielle} est une tige rigide qui \textbf{relie le piston au vilebrequin}. Elle est articulée à ses deux extrémités : sur le piston (par l'axe de piston) et sur le maneton du vilebrequin. Elle transmet la poussée du piston tout en accompagnant la rotation.
\end{definition}

\begin{definition}[Le vilebrequin]
Le \cle{vilebrequin} (ou \textbf{arbre moteur}) est un arbre coudé qui reçoit la poussée transmise par la bielle. Il \textbf{transforme le mouvement alternatif (aller-retour) du piston en mouvement de rotation continu}, rotation ensuite transmise aux roues par l'embrayage, la boîte de vitesses et la transmission (chaîne pour la moto).
\end{definition}

\begin{propriete}[Transformation du mouvement (admise)]
Le système \textbf{bielle-vilebrequin} (appelé aussi système bielle-manivelle) transforme le \textbf{mouvement rectiligne alternatif} du piston en \textbf{mouvement de rotation continu} du vilebrequin. Inversement, au démarrage, c'est la rotation du démarreur qui provoque le va-et-vient des pistons. Cette propriété est admise : sa justification géométrique n'est pas exigible au BEPC, mais tu dois savoir la citer et l'illustrer sur un schéma.
\end{propriete}

\begin{experience}[Observer la transformation du mouvement]
\textbf{Matériel :} un vélo retourné, posé sur la selle et le guidon.\\
\textbf{Protocole :}
\begin{enumerate}
\item Fais tourner lentement une pédale à la main et observe le mouvement de ton genou imaginaire : la manivelle tourne, la jambe monterait et descendrait.
\item Repère la manivelle (elle joue le rôle du vilebrequin) et la jambe (elle joue le rôle de la bielle et du piston).
\end{enumerate}
\textbf{Observations :} le pied décrit un cercle, mais la jambe effectue un mouvement de montée-descente.\\
\textbf{Interprétation :} la manivelle transforme le va-et-vient de la jambe en rotation du pédalier. Dans le moteur, c'est le couple bielle-vilebrequin qui réalise la même transformation, mais dans le sens moteur : la poussée des gaz sur le piston devient la rotation du vilebrequin.
\end{experience}

\courssub{Schémas de synthèse}

\begin{popfigure}
\begin{center}
\begin{tikzpicture}[scale=0.9, every node/.style={font=\small}]
% Cylindre (bloc moteur)
\draw[fill=popblueL, thick] (-2.2,-1) rectangle (2.2,3.4);
% Intérieur du cylindre
\draw[fill=white, thick] (-1.1,-1) rectangle (1.1,3.4);
% Piston
\draw[fill=poporange, thick] (-1.1,0.2) rectangle (1.1,1.2);
\node at (0,0.7) {\textbf{Piston}};
% Culasse
\draw[fill=popmuted!40, thick] (-2.2,3.4) rectangle (2.2,4.2);
% Soupapes
\draw[fill=popgreen, thick] (-1.4,3.0) -- (-0.7,3.0) -- (-0.85,3.4) -- (-1.25,3.4) -- cycle;
\draw[fill=popgreen, thick] (0.7,3.0) -- (1.4,3.0) -- (1.25,3.4) -- (0.85,3.4) -- cycle;
\draw[thick] (-1.05,3.4) -- (-1.05,4.2);
\draw[thick] (1.05,3.4) -- (1.05,4.2);
% Bougie
\draw[fill=poppink, thick] (-0.15,3.4) rectangle (0.15,4.6);
\draw[thick] (-0.08,3.4) -- (-0.08,3.1) (0.08,3.4) -- (0.08,3.15);
% Bielle
\draw[line width=2.5pt, popdark] (0,0.7) -- (0.9,-2.1);
% Vilebrequin (manivelle)
\draw[line width=2.5pt, popdark] (0,-2.8) -- (0.9,-2.1);
\draw[fill=popdark] (0,-2.8) circle (0.12);
\draw[fill=popdark] (0.9,-2.1) circle (0.1);
\draw[dashed, popmuted] (0,-2.8) circle (1.06);
% Flèche rotation
\draw[-{Stealth}, thick, poppurple] (1.5,-2.2) arc (40:-40:1.1);
% Flèche va-et-vient piston
\draw[{Stealth}-{Stealth}, thick, poppurple] (1.7,0.2) -- (1.7,1.2);
% Légendes fléchées
\node[right, align=left] at (2.6,4.4) {\textbf{1. Bougie}};
\draw[-{Stealth}, popink] (2.6,4.3) -- (0.2,4.2);
\node[right, align=left] at (2.6,3.5) {\textbf{2. Soupapes}};
\draw[-{Stealth}, popink] (2.6,3.4) -- (1.35,3.2);
\node[left, align=right] at (-2.6,3.8) {\textbf{3. Cylindre}};
\draw[-{Stealth}, popink] (-2.6,3.7) -- (-2.2,2.6);
\node[left, align=right] at (-2.6,0.7) {\textbf{4. Piston}};
\draw[-{Stealth}, popink] (-2.6,0.7) -- (-1.15,0.7);
\node[right, align=left] at (2.6,-0.8) {\textbf{5. Bielle}};
\draw[-{Stealth}, popink] (2.6,-0.85) -- (0.6,-0.8);
\node[right, align=left] at (2.6,-2.8) {\textbf{6. Vilebrequin}};
\draw[-{Stealth}, popink] (2.6,-2.8) -- (0.15,-2.8);
\end{tikzpicture}
\end{center}
\legende{Coupe schématique d'un cylindre de moteur à essence : 1. bougie ; 2. soupapes (admission et échappement) ; 3. cylindre ; 4. piston ; 5. bielle ; 6. vilebrequin. La flèche violette verticale indique le va-et-vient du piston ; la flèche courbe indique la rotation du vilebrequin.}
\end{popfigure}

\begin{popfigure}
\begin{center}
\begin{tikzpicture}[scale=1.0, every node/.style={font=\small}]
% Cylindre
\draw[fill=popblueL, thick] (-0.8,1.2) rectangle (0.8,3.6);
\draw[fill=white, thick] (-0.55,1.2) rectangle (0.55,3.6);
% Piston
\draw[fill=poporange, thick] (-0.55,1.6) rectangle (0.55,2.4);
\node at (0,2.0) {\textbf{P}};
% Axe piston
\draw[fill=popdark] (0,1.6) circle (0.08);
% Bielle
\draw[line width=2.5pt, popdark] (0,1.6) -- (0.75,-0.35);
% Manivelle
\draw[line width=2.5pt, popdark] (0,-1.0) -- (0.75,-0.35);
\draw[fill=popdark] (0,-1.0) circle (0.1);
\draw[fill=popdark] (0.75,-0.35) circle (0.08);
% Trajectoire du maneton
\draw[dashed, popmuted] (0,-1.0) circle (0.75);
% Flèche rotation
\draw[-{Stealth}, very thick, poppurple] (1.15,-0.45) arc (35:-35:0.85);
\node[right, poppurple] at (1.3,-1.0) {rotation};
% Flèche alternatif
\draw[{Stealth}-{Stealth}, very thick, poppurple] (1.3,1.5) -- (1.3,2.5);
\node[right, poppurple, align=left] at (1.45,2.0) {va-et-vient};
% Labels
\node[left, align=right] at (-1.0,2.0) {piston};
\draw[-{Stealth}, popink] (-1.0,2.0) -- (-0.6,2.0);
\node[left, align=right] at (-1.0,0.6) {bielle};
\draw[-{Stealth}, popink] (-1.0,0.6) -- (0.25,0.55);
\node[below] at (0,-2.0) {vilebrequin (manivelle)};
\draw[-{Stealth}, popink] (0,-1.9) -- (0,-1.15);
\end{tikzpicture}
\end{center}
\legende{Le système bielle-manivelle : le va-et-vient rectiligne du piston (P) entraîne, par l'intermédiaire de la bielle, la rotation continue du vilebrequin. Le maneton (point de fixation de la bielle sur le vilebrequin) décrit un cercle en pointillés.}
\end{popfigure}

\begin{methode}[Identifier les organes sur un schéma de moteur]
Face à un schéma de moteur au BEPC, procède toujours dans le même ordre :
\begin{enumerate}
\item \textbf{Suis le trajet des gaz} : repère l'entrée (admission, côté carburateur ou injecteur), la chambre de combustion (haut du cylindre), puis la sortie (échappement).
\item \textbf{Repère les pièces fixes} : le cylindre (le tube) et la culasse (le couvercle qui porte soupapes et bougie ou injecteur).
\item \textbf{Repère les pièces mobiles} de haut en bas : le piston (dans le cylindre), la bielle (la tige), le vilebrequin (l'arbre coudé en bas).
\item \textbf{Identifie le type de moteur} : présence d'une bougie $\Rightarrow$ moteur à essence ; présence d'un injecteur $\Rightarrow$ moteur diesel.
\end{enumerate}
\end{methode}

\begin{exoresolu}[Cylindrée d'une voiture]
\textbf{Énoncé.} Une voiture possède un moteur de 4 cylindres identiques. Sa cylindrée totale affichée est de \SI{1300}{cm^3}. (1) Rappeler la définition de la cylindrée. (2) Calculer la cylindrée d'un cylindre. (3) Citer l'organe qui transforme le mouvement du piston en rotation.
\tcblower
\textbf{Solution détaillée.}
\begin{enumerate}
\item La cylindrée est le \textbf{volume balayé par le piston entre le point mort haut (PMH) et le point mort bas (PMB)} ; pour un moteur à plusieurs cylindres, on additionne les volumes balayés par chaque piston.
\item Les 4 cylindres étant identiques :
\[ C_{\text{1 cylindre}} = \dfrac{C_{\text{totale}}}{4} = \dfrac{1300}{4} = \SI{325}{cm^3}. \]
\item C'est le \textbf{vilebrequin}, entraîné par la \textbf{bielle}, qui transforme le mouvement rectiligne alternatif du piston en mouvement de rotation continu.
\end{enumerate}
\end{exoresolu}

\begin{savaistu}[Pourquoi plusieurs cylindres ?]
Un moteur à un seul cylindre (moto, groupe électrogène) produit des à-coups : une seule poussée tous les deux tours de vilebrequin. En associant 3, 4, 6 cylindres ou plus dont les pistons poussent à tour de rôle, on obtient une rotation plus régulière et plus puissante. C'est pourquoi les grosses voitures et les camions qui montent la côte de Djebelé ou les pentes de l'Ouest ont de gros moteurs à plusieurs cylindres. Le vilebrequin possède alors autant de manetons (coudes) que de cylindres.
\end{savaistu}

\begin{aretenir}[L'essentiel de la section]
\begin{itemize}
\item Le \textbf{cylindre} est la chambre où coulisse le piston et où se produit la combustion.
\item Le \textbf{piston} comprime le mélange et reçoit la poussée des gaz brûlés ; ses positions extrêmes sont le \textbf{PMH} et le \textbf{PMB}.
\item La \textbf{cylindrée} est le volume balayé par le piston entre le PMH et le PMB ; elle s'exprime en cm$^3$.
\item Les \textbf{soupapes} (admission, échappement) commandent l'entrée du mélange et la sortie des gaz brûlés.
\item La \textbf{bougie} produit l'étincelle (moteur à essence) ; l'\textbf{injecteur} pulvérise le gasoil (moteur diesel).
\item Le \textbf{carburateur} ou le \textbf{système d'injection} prépare le mélange air-essence.
\item La \textbf{bielle} relie le piston au \textbf{vilebrequin}, qui transforme le mouvement alternatif en rotation.
\end{itemize}
\end{aretenir}

\coursec{Le cycle à quatre temps}

Dans la section précédente, nous avons démonté (en imagination !) le moteur de la moto-taxi : nous avons identifié le \cle{cylindre}, le \cle{piston}, les \cle{soupapes}, la \cle{bougie}, la \cle{bielle} et le \cle{vilebrequin}. Nous savons donc maintenant \emph{de quoi} le moteur est fait. Reste la grande question : \emph{comment} toutes ces pièces travaillent-elles ensemble pour transformer l'essence en mouvement ? La réponse tient en deux mots : le \cle{cycle à quatre temps}. C'est le cœur de cette leçon, et une question classique du BEPC.

\courssub{Principe général du cycle}

\begin{experience}[Observation : le rythme d'une moto au ralenti]
\textbf{Matériel :} une moto-taxi au ralenti, une montre ou un chronomètre.

\textbf{Protocole :} approche-toi (sans toucher !) d'une moto dont le moteur tourne au ralenti. Écoute attentivement le bruit du pot d'échappement.

\textbf{Observations :} on entend une succession régulière de « poum-poum-poum ». Le bruit n'est pas continu : il est fait de coups réguliers, séparés par de courts silences.

\textbf{Interprétation :} chaque « poum » correspond à une \cle{explosion} (une combustion brutale) dans le cylindre, suivie de l'évacuation des gaz brûlés. Le moteur ne travaille donc pas en continu : il fonctionne par \emph{coups successifs}, c'est-à-dire par \cle{cycles} qui se répètent tant qu'il y a du carburant et de l'air.

\textbf{Conclusion :} le fonctionnement d'un moteur à combustion interne est \emph{cyclique} : la même suite d'opérations se répète indéfiniment.
\end{experience}

\begin{definition}[Cycle à quatre temps]
Le \cle{cycle à quatre temps} est la suite de quatre opérations — \cle{admission}, \cle{compression}, \cle{détente} (ou explosion), \cle{échappement} — au cours desquelles le piston effectue \textbf{quatre courses} (deux descentes et deux montées), ce qui correspond à \textbf{deux tours de vilebrequin}. Ce cycle se répète tant que le moteur tourne.
\end{definition}

\begin{aretenir}[L'ordre des temps à mémoriser]
L'ordre des quatre temps est immuable :
\begin{center}
\textbf{1. Admission \quad $\rightarrow$ \quad 2. Compression \quad $\rightarrow$ \quad 3. Détente \quad $\rightarrow$ \quad 4. Échappement}
\end{center}
Astuce de mémorisation : « \emph{À Coton Douala, l'Essence} » (Admission, Compression, Détente, Échappement). Attention : la détente s'appelle aussi \emph{explosion} ou \emph{temps moteur}.
\end{aretenir}

\courssub{Premier temps : l'admission}

Imagine que tu tires la seringue d'une pompe à vélo : en tirant le piston, tu crées de la place et l'air entre. L'admission fonctionne exactement ainsi.

\begin{itemize}
\item La \cle{soupape d'admission} \textbf{s'ouvre} ; la soupape d'échappement reste fermée.
\item Le \textbf{piston descend} du point mort haut (PMH) vers le point mort bas (PMB) : c'est la première course.
\item En descendant, le piston aspire dans le cylindre le \cle{mélange air-essence} préparé par le carburateur (ou l'injection). Dans un moteur diesel, c'est de \textbf{l'air seul} qui entre (le gasoil sera injecté plus tard).
\end{itemize}

Le vilebrequin a effectué un \textbf{demi-tour}. Ce temps ne produit aucun travail : c'est le vilebrequin, grâce à son inertie, qui entraîne le piston.

\courssub{Deuxième temps : la compression}

\begin{itemize}
\item Les \textbf{deux soupapes sont fermées} : le cylindre est hermétiquement clos.
\item Le \textbf{piston remonte} du PMB vers le PMH : deuxième course.
\item Le mélange air-essence (ou l'air seul en diesel), prisonnier, est comprimé dans un volume de plus en plus petit : sa \cle{pression} et sa \cle{température} augmentent fortement.
\end{itemize}

\begin{experience}[Pourquoi la compression chauffe-t-elle le mélange ?]
\textbf{Matériel :} une pompe à vélo.

\textbf{Protocole :} bouche la sortie de la pompe avec un doigt, puis pousse vigoureusement le piston. Touche ensuite le corps de la pompe.

\textbf{Observation :} la pompe s'échauffe nettement.

\textbf{Interprétation :} comprimer un gaz augmente sa température. Dans le cylindre, le mélange comprimé atteint plusieurs centaines de degrés : il est prêt à brûler instantanément. C'est aussi ce principe qui permet au moteur diesel de s'allumer sans bougie (auto-allumage par compression).
\end{experience}

Ce temps consomme du travail (il faut pousser le piston) : c'est encore le vilebrequin qui le fournit. Le vilebrequin a maintenant effectué \textbf{un tour complet}.

\courssub{Troisième temps : la détente, le temps moteur}

C'est le moment décisif, le « poum » entendu au ralenti.

\begin{itemize}
\item Les \textbf{deux soupapes restent fermées}.
\item Au moment où le piston arrive en haut (PMH) :
  \begin{itemize}
  \item Dans un \textbf{moteur à essence} : la \cle{bougie} produit une \textbf{étincelle}. Le mélange comprimé et chaud \textbf{brûle brutalement} : c'est l'\cle{explosion}.
  \item Dans un \textbf{moteur diesel} : l'\cle{injecteur} pulvérise le gasoil qui s'enflamme spontanément au contact de l'air très chaud (auto-allumage).
  \end{itemize}
\item Les gaz brûlés, très chauds et à très haute pression, \textbf{poussent violemment le piston vers le bas} : troisième course.
\item Par l'intermédiaire de la bielle, cette poussée fait tourner le vilebrequin : l'énergie chimique du carburant est transformée en \cle{travail mécanique}.
\end{itemize}

\begin{definition}[Temps moteur]
Le \cle{temps moteur} (ou détente, ou explosion) est le temps pendant lequel la combustion brutale du mélange pousse le piston vers le bas et produit le \cle{travail mécanique} qui fait tourner le vilebrequin. C'est le \textbf{seul} des quatre temps qui fournit du travail.
\end{definition}

\begin{attention}[Piège classique du BEPC]
\textbf{Erreur fréquente :} croire que les quatre temps produisent du travail, ou que « moteur à quatre temps » signifie « quatre temps moteurs ». C'est faux ! Sur les quatre temps, \textbf{un seul est moteur} : la détente. Les trois autres temps (admission, compression, échappement) \emph{consomment} de l'énergie : ils sont entraînés grâce à l'\cle{inertie} du vilebrequin et du \cle{volant moteur}, une masse lourde solidaire du vilebrequin qui stocke l'énergie du temps moteur et la restitue pendant les autres temps. « Quatre temps » désigne le nombre de courses du piston par cycle, rien d'autre.
\end{attention}

\courssub{Quatrième temps : l'échappement}

\begin{itemize}
\item La \cle{soupape d'échappement} \textbf{s'ouvre} ; la soupape d'admission reste fermée.
\item Le \textbf{piston remonte} du PMB vers le PMH : quatrième course.
\item En remontant, il \textbf{chasse les gaz brûlés} hors du cylindre, vers le pot d'échappement.
\end{itemize}

Le vilebrequin achève son \textbf{deuxième tour}. Le cylindre est vidé de ses gaz brûlés : un nouveau cycle peut commencer par une nouvelle admission. Et ainsi de suite, des milliers de fois par minute !

\courssub{Schéma récapitulatif des quatre temps}

\begin{popfigure}
\begin{center}
\begin{tikzpicture}[scale=0.88, every node/.style={transform shape}]
% --- Case 1 : Admission ---
\begin{scope}[shift={(0,0)}]
\draw[thick, popdark] (-1.5,0) rectangle (1.5,3.2);
\draw[thick, popdark] (-1.5,3.2) -- (-1.5,4.2) -- (1.5,4.2) -- (1.5,3.2);
% soupapes
\draw[fill=popgreen] (-0.9,4.2) -- (-0.4,4.2) -- (-0.65,3.8) -- cycle;
\draw[fill=popmuted] (0.9,4.2) -- (0.4,4.2) -- (0.65,3.8) -- cycle;
\draw[-{Stealth}, popgreen, thick] (-0.65,4.9) -- (-0.65,4.3);
% piston bas + flèche descente
\draw[fill=popblueL, thick] (-1.4,0.7) rectangle (1.4,1.3);
\draw[-{Stealth}, popblue, very thick] (0,2.6) -- (0,1.6);
\node[popdark, font=\bfseries] at (0,-0.5) {1. Admission};
\node[popmuted, font=\small] at (0,-1.0) {piston descend};
\end{scope}
% --- Case 2 : Compression ---
\begin{scope}[shift={(4.2,0)}]
\draw[thick, popdark] (-1.5,0) rectangle (1.5,3.2);
\draw[thick, popdark] (-1.5,3.2) -- (-1.5,4.2) -- (1.5,4.2) -- (1.5,3.2);
\draw[fill=popmuted] (-0.9,4.2) -- (-0.4,4.2) -- (-0.65,3.8) -- cycle;
\draw[fill=popmuted] (0.9,4.2) -- (0.4,4.2) -- (0.65,3.8) -- cycle;
\draw[fill=popblueL, thick] (-1.4,2.0) rectangle (1.4,2.6);
\draw[-{Stealth}, popblue, very thick] (0,0.9) -- (0,1.9);
\node[popdark, font=\bfseries] at (0,-0.5) {2. Compression};
\node[popmuted, font=\small] at (0,-1.0) {piston monte};
\end{scope}
% --- Case 3 : Détente ---
\begin{scope}[shift={(8.4,0)}]
\draw[thick, popdark] (-1.5,0) rectangle (1.5,3.2);
\draw[thick, popdark] (-1.5,3.2) -- (-1.5,4.2) -- (1.5,4.2) -- (1.5,3.2);
\draw[fill=popmuted] (-0.9,4.2) -- (-0.4,4.2) -- (-0.65,3.8) -- cycle;
\draw[fill=popmuted] (0.9,4.2) -- (0.4,4.2) -- (0.65,3.8) -- cycle;
% étincelle / injection
\node[popgold, font=\large] at (0,3.9) {$\star$};
\node[popmuted, font=\tiny, align=center] at (0,3.5) {Étincelle (essence)\\\textit{ou injection (diesel)}};
\draw[fill=poporangeL, thick] (-1.4,2.0) rectangle (1.4,2.6);
\draw[-{Stealth}, poporange, very thick] (0,1.9) -- (0,0.9);
\node[popdark, font=\bfseries] at (0,-0.5) {3. Détente (moteur)};
\node[popmuted, font=\small] at (0,-1.0) {piston poussé vers le bas};
\end{scope}
% --- Case 4 : Échappement ---
\begin{scope}[shift={(12.6,0)}]
\draw[thick, popdark] (-1.5,0) rectangle (1.5,3.2);
\draw[thick, popdark] (-1.5,3.2) -- (-1.5,4.2) -- (1.5,4.2) -- (1.5,3.2);
\draw[fill=popmuted] (-0.9,4.2) -- (-0.4,4.2) -- (-0.65,3.8) -- cycle;
\draw[fill=poppink] (0.9,4.2) -- (0.4,4.2) -- (0.65,3.8) -- cycle;
\draw[-{Stealth}, poppink, thick] (0.65,4.3) -- (0.65,4.9);
\draw[fill=popblueL, thick] (-1.4,2.0) rectangle (1.4,2.6);
\draw[-{Stealth}, popblue, very thick] (0,0.9) -- (0,1.9);
\node[popdark, font=\bfseries] at (0,-0.5) {4. Échappement};
\node[popmuted, font=\small] at (0,-1.0) {piston monte};
\end{scope}
\end{tikzpicture}
\end{center}
\legende{Les quatre temps du cycle. En vert : soupape d'admission ouverte ; en rose : soupape d'échappement ouverte ; en gris : soupape fermée. Le rectangle clair représente le piston, la flèche bleue ou orange son mouvement, l'étoile l'étincelle (essence) ou l'injection (diesel).}
\end{popfigure}

\begin{center}
\rowcolors{1}{popcream}{white}
\begin{tabularx}{\linewidth}{|l|X|X|X|}
\hline
\rowcolor{popblueL} \textbf{Temps} & \textbf{Mouvement du piston} & \textbf{Soupapes} & \textbf{Rôle} \\
\hline
1. Admission & Descend (PMH $\rightarrow$ PMB) & Admission \textbf{ouverte}, échappement fermée & Le mélange air-essence (ou l'air seul en diesel) entre dans le cylindre \\
\hline
2. Compression & Monte (PMB $\rightarrow$ PMH) & Les deux \textbf{fermées} & Le mélange est comprimé : pression et température augmentent \\
\hline
3. Détente (temps moteur) & Descend (Poussé PMH $\rightarrow$ PMB) & Les deux \textbf{fermées} & La combustion pousse le piston : seul temps qui \textbf{produit du travail} \\
\hline
4. Échappement & Monte (PMB $\rightarrow$ PMH) & Échappement \textbf{ouverte}, admission fermée & Les gaz brûlés sont chassés hors du cylindre \\
\hline
\end{tabularx}
\end{center}

\courssub{Comment reconnaître un temps sur un schéma ?}

Au BEPC, on te présentera souvent un schéma de cylindre et on te demandera : « \emph{Quel temps est représenté ? Justifie.} » Voici la démarche infaillible.

\begin{methode}[Identifier un temps moteur sur un schéma]
\begin{enumerate}
\item \textbf{Observer les soupapes.} Une soupape ouverte ? Si c'est celle d'admission $\rightarrow$ c'est l'\emph{admission}. Si c'est celle d'échappement $\rightarrow$ c'est l'\emph{échappement}.
\item \textbf{Si les deux soupapes sont fermées}, regarder le sens de déplacement du piston (flèche) : piston qui \emph{monte} $\rightarrow$ \emph{compression} ; piston qui \emph{descend} $\rightarrow$ \emph{détente}.
\item \textbf{Chercher l'indice d'allumage} : étincelle de la bougie (moteur à essence) ou présence de l'injecteur (moteur diesel). Cela confirme le début de la détente (temps moteur).
\item \textbf{Rédiger la justification} en citant les deux indices : état des soupapes \emph{et} sens du mouvement du piston. Un seul indice ne suffit pas toujours !
\end{enumerate}
\end{methode}

\courssub{Lien avec la rotation du vilebrequin : calculs}

Récapitulons le bilan cinématique du cycle :
\begin{itemize}
\item 1 cycle = \textbf{4 courses} du piston = \textbf{2 tours de vilebrequin} ;
\item 1 cycle = \textbf{1 seule explosion} (un seul temps moteur) par cylindre.
\end{itemize}

On en déduit une relation très utile pour les exercices :
\[ \text{nombre d'explosions par minute (par cylindre)} = \frac{\text{vitesse de rotation (tr/min)}}{2} \]

\begin{exoresolu}[Nombre d'explosions dans le moteur d'une moto]
\textbf{Énoncé.} Le moteur (monocylindre) d'une moto-taxi tourne à \num{3000} tr/min.
\begin{enumerate}
\item Combien de temps moteur (explosions) se produisent-ils par minute dans le cylindre ?
\item Combien de cycles complets le piston effectue-t-il en une seconde ?
\end{enumerate}
\tcblower
\textbf{Solution détaillée.}
\begin{enumerate}
\item Un cycle complet (qui contient une seule explosion) correspond à 2 tours de vilebrequin. Donc :
\[ N = \frac{\num{3000}}{2} = \num{1500} \ \text{explosions par minute.} \]
\item Le nombre de cycles par minute est aussi de \num{1500} (un cycle = une explosion). En une seconde :
\[ \frac{\num{1500}}{60} = 25 \ \text{cycles par seconde.} \]
Le piston effectue donc 25 cycles complets, soit 100 courses, chaque seconde !
\end{enumerate}
\textbf{Remarque :} si la moto avait deux cylindres, il y aurait $2 \times \num{1500} = \num{3000}$ explosions par minute au total, mais toujours \num{1500} \emph{par cylindre}.
\end{exoresolu}

\begin{exercice}[À toi de jouer]
Le moteur d'un groupe électrogène tourne à \num{3600} tr/min.
\begin{enumerate}
\item Combien d'explosions se produisent-elles par minute dans son unique cylindre ?
\item Pendant ces \num{3600} tours, combien de courses le piston a-t-il effectuées ?
\end{enumerate}
\end{exercice}

\begin{corrige}[Exercice : groupe électrogène]
\begin{enumerate}
\item Nombre d'explosions par minute : $\dfrac{\num{3600}}{2} = \num{1800}$ explosions par minute.
\item Un tour de vilebrequin correspond à 2 courses du piston (une montée et une descente). Donc : $\num{3600} \times 2 = \num{7200}$ courses par minute.
\end{enumerate}
\end{corrige}

\begin{attention}[Entretien et sécurité liés au cycle]
\begin{itemize}
\item \textbf{Huile moteur :} elle lubrifie le piston dans le cylindre (frottements lors des 4 courses) et les paliers du vilebrequin. Sans huile, le moteur « serre » (grippage) en quelques minutes. Vérifie le niveau \emph{tous les matins} sur terrain plat, moteur froid.
\item \textbf{Bougie (essence) / Injecteur (diesel) :} une bougie encrassée ou un injecteur bouché rate des explosions $\rightarrow$ perte de puissance, consommation excessive, démarrage difficile. Change la bougie selon le carnet d'entretien (souvent tous les \num{10000}--\num{15000} km).
\item \textbf{Filtre à air :} s'il est bouché, l'admission est pauvre $\rightarrow$ mélange trop riche, fumée noire, encrassement. Nettoie-le ou change-le régulièrement (poussière rouge du Sahel ou poussière des routes non bitumées !).
\item \textbf{Refroidissement :} la détente chauffe énormément le cylindre. Le liquide de refroidissement (eau + antigel) ou l'aération (moto) doit évacuer cette chaleur. Ne \textbf{jamais} ouvrir le vase d'expansion d'un moteur chaud : risque de brûlures graves par ébullition retardée.
\item \textbf{Pot d'échappement :} gaz toxiques (monoxyde de carbone CO inodore et mortel). Ne jamais faire tourner un moteur (moto, groupe électrogène) dans un local fermé (garage, chambre).
\end{itemize}
\end{attention}

\begin{savaistu}[Le « tac-tac-tac » de la moto]
Le bruit régulier « tac-tac-tac » ou « poum-poum-poum » d'une moto au ralenti n'est autre que la succession des explosions dans le cylindre, suivies de l'évacuation des gaz par le pot d'échappement. Quand le conducteur accélère, le régime monte (par exemple de \num{1500} à \num{6000} tr/min), les explosions se rapprochent, et le bruit devient un vrombissement presque continu. Certains mécaniciens expérimentés de Douala ou de Yaoundé diagnostiquent une panne rien qu'à l'oreille : un rythme irrégulier trahit une bougie encrassée qui « rate » des explosions, ou une prise d'air parasite qui déséquilibre le mélange !
\end{savaistu}

\begin{aretenir}[Bilan de la section]
\begin{itemize}
\item Le moteur à quatre temps fonctionne par \cle{cycles} : \textbf{admission, compression, détente, échappement}.
\item Un cycle = 4 courses du piston = \textbf{2 tours de vilebrequin} = \textbf{1 explosion} par cylindre.
\item Seul le \textbf{3\up{e} temps (détente)} est moteur ; les trois autres sont entraînés par l'inertie du vilebrequin et du volant moteur.
\item Différence essence/diesel : allumage par \textbf{bougie} (essence) vs \textbf{auto-allumage par compression} + injection (diesel).
\item Pour identifier un temps sur un schéma : regarder l'état des soupapes \emph{et} le sens de déplacement du piston.
\item Nombre d'explosions par minute et par cylindre = vitesse de rotation (tr/min) divisée par 2.
\end{itemize}
\end{aretenir}

\coursec{Moteur à essence et moteur diesel : différences}

Dans la section précédente, nous avons étudié le \cle{cycle à quatre temps} : admission, compression, détente (temps moteur), échappement. Nous avons vu que ce cycle est commun à presque tous les moteurs qui font rouler nos moto-taxis, nos voitures, nos camions et tourner nos groupes électrogènes. Mais à la station-service Total ou Ola Energy, tu as sûrement remarqué deux pompes différentes : l'une pour la \emph{super} (essence), l'autre pour le \emph{gasoil} (diesel). Pourquoi deux carburants ? Parce qu'il existe deux grandes familles de moteurs à combustion interne : le \cle{moteur à essence} et le \cle{moteur diesel}. Ils fonctionnent tous les deux selon le cycle à quatre temps, mais ils diffèrent par un point essentiel : \textbf{la manière dont le mélange s'enflamme}. C'est toute l'originalité de cette section.

\courssub{Le point commun : deux moteurs à quatre temps}

Avant d'opposer les deux moteurs, insistons sur ce qui les rassemble.

\begin{aretenir}[Point commun fondamental]
Le moteur à essence et le moteur diesel sont tous les deux des \cle{moteurs à combustion interne} fonctionnant selon le \cle{cycle à quatre temps} : admission, compression, détente (temps moteur), échappement. Dans les deux cas, l'énergie provient de la combustion d'un carburant dérivé du pétrole, à l'intérieur même du cylindre.
\end{aretenir}

La différence ne se situe donc ni dans le nombre de temps, ni dans le principe général (une combustion pousse le piston). Elle se situe dans \textbf{ce qui est admis dans le cylindre} et dans \textbf{ce qui déclenche la combustion}. Pour le comprendre, menons une petite investigation.

\courssub{Une observation qui pose question : la pompe à vélo}

\begin{experience}[Comprimer un gaz le chauffe-t-il ?]
\textbf{Matériel :} une pompe à vélo ou à ballon, ta main.

\textbf{Protocole :}
\begin{enumerate}
\item Bouche l'extrémité du tuyau de la pompe avec un doigt.
\item Enfonce le piston de la pompe vigoureusement plusieurs fois de suite.
\item Touche le corps métallique de la pompe, près de sa base.
\end{enumerate}

\textbf{Observation :} le corps de la pompe devient nettement chaud, alors qu'aucune flamme ni aucune source de chaleur ne l'a touché.

\textbf{Interprétation :} quand on comprime l'air (on le force à occuper un volume plus petit), on lui communique de l'énergie. Cette énergie se transforme en chaleur : la température de l'air comprimé augmente, et cette chaleur se transmet au métal de la pompe.

\textbf{Conclusion :} comprimer un gaz élève sa température. Plus la compression est forte, plus l'échauffement est important.
\end{experience}

Cette observation toute simple est la clé de toute la différence entre les deux moteurs.

\begin{propriete}[Propriété admise : la compression chauffe les gaz]
Quand on comprime un gaz, sa température augmente. Plus le \cle{taux de compression} est élevé, plus la température atteinte est grande. Dans un moteur diesel, l'air seul est comprimé si fortement qu'il atteint une température suffisante pour enflammer spontanément le gasoil qu'on y injecte. Cette propriété est admise au niveau de la classe de 3ème (sa démonstration n'est pas exigible au BEPC).
\end{propriete}

\courssub{Le moteur à essence : allumage commandé par la bougie}

Dans un moteur à essence, le cylindre n'admet pas de l'air pur : il aspire un \cle{mélange air + essence} déjà préparé. Autrefois, ce mélange était préparé par un organe appelé \cle{carburateur} ; aujourd'hui, sur les véhicules modernes, c'est un système d'\cle{injection} électronique (injection indirecte) qui pulvérise l'essence dans l'air admis, avant l'entrée dans le cylindre.

Ce mélange est ensuite comprimé par le piston, mais \emph{modérément} : si on le comprimait trop, il s'enflammerait tout seul n'importe quand (cliquetis), et le moteur fonctionnerait mal. Il faut donc déclencher la combustion \textbf{au bon moment}, précisément lorsque le piston arrive en haut (point mort haut). C'est le rôle de la \cle{bougie} : au moment voulu, une étincelle électrique jaillit entre ses électrodes et met le feu au mélange. On parle d'\cle{allumage commandé}.

\begin{definition}[Allumage commandé (moteur à essence)]
L'\cle{allumage commandé} est le mode d'inflammation dans lequel la combustion du mélange air--essence est déclenchée, à un instant précis, par l'étincelle électrique produite par la \cle{bougie}. Sans bougie (ou si la bougie est défectueuse), le mélange ne s'enflamme pas et le moteur ne démarre pas.
\end{definition}

\begin{exemplebox}[La panne de bougie du moto-taxi]
Un moto-taxi de quartier refuse de démarrer un matin. Le mécanicien dévisse la bougie, la présente contre le moteur et actionne le démarreur : aucune étincelle ne jaillit. Il remplace la bougie, et le moteur repart aussitôt. La bougie est l'organe d'allumage du moteur à essence : c'est souvent la première pièce qu'on vérifie sur une moto en panne.
\end{exemplebox}

\courssub{Le moteur diesel : allumage par compression}

Le moteur diesel, inventé par l'ingénieur allemand Rudolf Diesel à la fin du XIX\textsuperscript{e} siècle, fonctionne de manière plus astucieuse : il se passe complètement de bougie \emph{d'allumage}.

Voici le raisonnement, étape par étape :

\begin{enumerate}
\item \textbf{Admission :} le cylindre n'aspire que de l'\emph{air seul}, sans aucun carburant.
\item \textbf{Compression :} le piston comprime cet air \emph{très fortement}, beaucoup plus que dans un moteur à essence. D'après la propriété admise plus haut, l'air devient alors très chaud (plusieurs centaines de degrés Celsius).
\item \textbf{Injection et inflammation :} à cet instant précis (point mort haut), un organe appelé \cle{injecteur} pulvérise le gasoil en fines gouttelettes dans cet air brûlant. Au contact de l'air chaud, le gasoil s'enflamme \emph{spontanément}, sans aucune étincelle. C'est l'\cle{auto-inflammation}.
\item \textbf{Détente et échappement :} la combustion pousse le piston (temps moteur), puis les gaz brûlés sont évacués, comme dans tout moteur à quatre temps.
\end{enumerate}

\begin{definition}[Allumage par compression (moteur diesel)]
L'\cle{allumage par compression} (ou auto-inflammation) est le mode d'inflammation dans lequel le carburant (le gasoil), injecté dans de l'air fortement comprimé et donc très chaud, s'enflamme \emph{spontanément}, sans étincelle. Le moteur diesel ne possède donc \textbf{pas de bougie d'allumage} : c'est la compression elle-même qui provoque l'allumage. (Il peut être équipé de \emph{bougies de préchauffage} pour faciliter le démarrage à froid, mais elles ne servent pas à l'allumage en régime normal.)
\end{definition}

\begin{aretenir}[Le taux de compression]
Le \cle{taux de compression} compare le volume du cylindre quand le piston est en bas au volume quand il est en haut. Il est \textbf{plus élevé dans un moteur diesel que dans un moteur à essence}. C'est cette compression plus forte qui permet de chauffer suffisamment l'air pour enflammer le gasoil sans bougie. (Valeur admise à ce niveau : inutile de retenir des chiffres précis.)
\end{aretenir}

\courssub{Schéma comparatif des deux culasses}

La \cle{culasse} est la pièce qui ferme le haut du cylindre. C'est sur elle que se trouve l'organe d'allumage : la bougie pour l'essence, l'injecteur pour le diesel. Observons les deux schémas côte à côte.

\begin{popfigure}
\begin{center}
\begin{tikzpicture}[scale=0.9, every node/.style={font=\small}]
% ===== MOTEUR ESSENCE =====
% cylindre
\draw[thick, popblue] (0,0) -- (0,3.2) (2.4,0) -- (2.4,3.2);
% culasse
\draw[thick, fill=popblueL] (0,3.2) rectangle (2.4,3.9);
% bougie
\draw[thick, fill=popgoldL] (1.0,3.9) rectangle (1.4,4.7);
\draw[thick, popdark] (1.2,3.9) -- (1.2,3.55);
\draw[poporange, thick] (1.05,3.45) -- (1.2,3.55) -- (1.35,3.45);
\node[poporange] at (1.2,3.25) {\scriptsize $\star$};
% piston
\draw[thick, fill=popblue!40] (0.08,0.6) rectangle (2.32,1.3);
\draw[thick] (1.2,0.6) -- (1.2,-0.5);
% flèches mélange
\draw[-{Stealth}, popgreen, thick] (-0.9,3.55) -- (-0.05,3.55);
\node[popgreen, align=center] at (-1.15,2.9) {\scriptsize mélange\\ \scriptsize air+essence};
% labels
\node[font=\bfseries, popblue] at (1.2,5.1) {Moteur à essence};
\draw[-{Stealth}, popdark] (1.45,4.3) -- (2.9,4.3) node[right, align=left] {\scriptsize bougie (étincelle)};
\node at (1.2,0.95) {\scriptsize piston};
\node[align=center] at (1.2,2.2) {\scriptsize mélange\\ \scriptsize comprimé};

% ===== MOTEUR DIESEL =====
\begin{scope}[xshift=7.2cm]
\draw[thick, poporange] (0,0) -- (0,3.2) (2.4,0) -- (2.4,3.2);
\draw[thick, fill=poporangeL] (0,3.2) rectangle (2.4,3.9);
% injecteur
\draw[thick, fill=poppurple!30] (1.0,3.9) rectangle (1.4,4.7);
\draw[thick, popdark] (1.2,3.9) -- (1.2,3.55);
\draw[poppurple, thick] (1.2,3.55) -- (0.9,3.15);
\draw[poppurple, thick] (1.2,3.55) -- (1.2,3.1);
\draw[poppurple, thick] (1.2,3.55) -- (1.5,3.15);
% piston
\draw[thick, fill=poporange!40] (0.08,0.6) rectangle (2.32,1.3);
\draw[thick] (1.2,0.6) -- (1.2,-0.5);
% flèche air seul
\draw[-{Stealth}, popgreen, thick] (-0.9,3.55) -- (-0.05,3.55);
\node[popgreen, align=center] at (-1.0,2.9) {\scriptsize air seul};
% labels
\node[font=\bfseries, poporange] at (1.2,5.1) {Moteur diesel};
\draw[-{Stealth}, popdark] (1.45,4.3) -- (2.9,4.3) node[right, align=left] {\scriptsize injecteur (gasoil)};
\node at (1.2,0.95) {\scriptsize piston};
\node[align=center] at (1.2,2.2) {\scriptsize air très\\ \scriptsize comprimé,\\ \scriptsize très chaud};
\end{scope}
\end{tikzpicture}
\end{center}
\legende{Comparaison des deux culasses. À gauche (essence) : la \textbf{bougie} produit une étincelle qui enflamme le mélange air + essence admis dans le cylindre. À droite (diesel) : l'\textbf{injecteur} pulvérise le gasoil dans l'air seul, fortement comprimé et très chaud ; le gasoil s'enflamme spontanément. C'est l'organe présent sur la culasse qui permet de reconnaître le type de moteur.}
\end{popfigure}

\begin{methode}[Comment distinguer les deux moteurs sur un schéma]
Face à un schéma de moteur au BEPC ou dans un exercice, procède ainsi :
\begin{enumerate}
\item Repère la \cle{culasse}, la pièce qui ferme le haut du cylindre.
\item Cherche l'organe d'allumage fixé sur la culasse.
\item Si tu vois une \textbf{bougie} (avec une étincelle symbolisée) : c'est un \textbf{moteur à essence}, à allumage commandé.
\item Si tu vois un \textbf{injecteur} (avec un jet de carburant pulvérisé) et aucune bougie : c'est un \textbf{moteur diesel}, à allumage par compression.
\item Vérifie en cohérence : mélange air + essence admis $\rightarrow$ essence ; air seul admis puis carburant injecté $\rightarrow$ diesel.
\end{enumerate}
\end{methode}

\courssub{Tableau comparatif des deux moteurs}

Récapitulons l'ensemble des différences dans un tableau à cinq lignes, à savoir reproduire et compléter le jour du BEPC.

\begin{center}
\begin{tabularx}{\linewidth}{|l|X|X|}
\hline
\rowcolor{popblueL} \textbf{Critère} & \textbf{Moteur à essence} & \textbf{Moteur diesel} \\
\hline
\textbf{Carburant} & Essence (super) & Gasoil (diesel) \\
\hline
\textbf{Organe d'allumage} & Bougie & Injecteur (pas de bougie d'allumage) \\
\hline
\textbf{Mélange admis} & Mélange air + essence (préparé par le carburateur ou l'injection indirecte) & Air seul ; le gasoil est injecté ensuite dans le cylindre \\
\hline
\textbf{Mode d'allumage} & Allumage commandé (étincelle de la bougie) & Allumage par compression (auto-inflammation dans l'air chaud) \\
\hline
\textbf{Usages courants} & Motos, voitures légères (moteur léger et souple) & Camions, bus, groupes électrogènes, tracteurs (moteur robuste et économique) \\
\hline
\end{tabularx}
\end{center}

\courssub{Conséquences pratiques : qui utilise quoi, et pourquoi ?}

Ces différences techniques ont des conséquences directes sur la vie quotidienne et l'économie des transports au Cameroun.

\begin{propriete}[Conséquences pratiques des deux conceptions]
\begin{itemize}
\item Le \cle{moteur diesel}, soumis à de fortes compressions, est construit avec des pièces plus résistantes : il est donc plus \textbf{robuste}, dure plus longtemps et consomme moins de carburant sur les longues distances. Il équipe les camions, les bus, les tracteurs et les groupes électrogènes.
\item Le \cle{moteur à essence} est plus \textbf{léger}, plus compact et plus \textbf{souple} (il monte facilement en régime). Il équipe les motos et les voitures légères.
\end{itemize}
\end{propriete}

\begin{savaistu}[Le gasoil fait rouler le Cameroun]
Les camions-citernes qui approvisionnent les brasseries (SABC) et les bus de voyage qui relient Douala à Yaoundé, Bafoussam ou Garoua (Finexs, General Express, Touristique...) roulent presque tous au gasoil. Pourquoi ? Parce que sur les longues distances, le moteur diesel consomme moins de carburant et supporte mieux les charges lourdes. À l'inverse, les centaines de moto-taxis de nos villes utilisent des moteurs à essence, légers et nerveux, parfaits pour les petits trajets. Et le groupe électrogène du quartier, qui prend le relais pendant les délestages ENEO ? Le plus souvent, c'est aussi un diesel, pour sa robustesse et son économie. La SONARA (Société Nationale de Raffinage) à Limbé produit ces deux carburants à partir du pétrole brut.
\end{savaistu}

\courssub{Sécurité, entretien et environnement}

Au-delà du choix du carburant, l'utilisation d'un moteur impose des règles de sécurité, un entretien régulier et une conscience des impacts environnementaux.

\begin{attention}[Erreur de carburant à la station-service : un piège réel]
Mettre de l'\textbf{essence dans un réservoir diesel} (ou du gasoil dans un réservoir essence) \textbf{endommage gravement le moteur} : le carburant inadapté ne s'enflamme pas correctement, les organes d'injection (pompe haute pression, injecteurs) sont détériorés par manque de lubrification (essence) ou par calage (gasoil), et la réparation coûte très cher. D'où l'importance de la vigilance à la station-service : vérifie toujours la pompe avant de faire le plein, et si l'erreur est commise, \textbf{ne démarre surtout pas} le moteur ; fais vidanger le réservoir par un professionnel. Autre piège de rédaction au BEPC : ne dis jamais qu'un moteur diesel possède une bougie d'allumage — c'est faux, son allumage se fait par compression.
\end{attention}

\begin{propriete}[Entretien courant et sécurité]
\begin{itemize}
\item \textbf{Moteur à essence :} vérifier et changer la \textbf{bougie} périodiquement (usure des électrodes), contrôler le niveau d'huile moteur, changer le filtre à air et le filtre à essence.
\item \textbf{Moteur diesel :} pas de bougie d'allumage à changer, mais surveillance stricte du \textbf{filtre à gasoil} (eau et impuretés) et de la pompe d'injection. Purger le filtre à gasoil (vis de purge) régulièrement.
\item \textbf{Commun :} niveau d'huile (jauge) à froid sur terrain plat, niveau de liquide de refroidissement, état des courroies (distribution, accessoires), batterie (bornes propres, serrées). Ne jamais ouvrir le vase d'expansion du liquide de refroidissement quand le moteur est chaud (risque de brûlures graves par projection).
\end{itemize}
\end{propriete}

\begin{propriete}[Environnement : émissions et dépollution]
La combustion n'est jamais parfaite. Les gaz d'échappement contiennent :
\begin{itemize}
\item \textbf{Dioxyde de carbone (\ce{CO2}) :} gaz à effet de serre, responsable du réchauffement climatique. Le diesel en émet moins par km (meilleur rendement).
\item \textbf{Monoxyde de carbone (\ce{CO}) :} gaz toxique, mortel en espace clos (garage fermé, groupe électrogène mal ventilé). Ne jamais faire tourner un moteur dans un local non aéré.
\item \textbf{Oxydes d'azote (\ce{NOx}) et particules fines :} plus émis par les diesels anciens. Les véhicules modernes ont un \textbf{catalyseur} (essence) ou un \textbf{filtre à particules (FAP)} et une vanne EGR (diesel) pour les réduire.
\end{itemize}
L'entretien régulier (moteur bien réglé) et l'usage de carburants de qualité (normes à la pompe) limitent la pollution.
\end{propriete}

\begin{exoresolu}[Identifier un moteur et justifier]
\textbf{Énoncé.} Sur le schéma d'un moteur à quatre temps, on observe sur la culasse un organe qui produit une étincelle électrique au moment où le piston atteint le point haut. Le manuel du véhicule précise que le réservoir contient de la super.
\begin{enumerate}
\item De quel type de moteur s'agit-il ? Justifie ta réponse.
\item Comment s'appelle l'organe observé sur la culasse ? Quel est son rôle ?
\item Quel mélange est admis dans le cylindre pendant le premier temps ?
\item Cite un usage courant de ce type de moteur au Cameroun.
\end{enumerate}
\tcblower
\textbf{Solution détaillée.}
\begin{enumerate}
\item Il s'agit d'un \textbf{moteur à essence}. Justification : la combustion est déclenchée par une étincelle électrique (allumage commandé), et le carburant utilisé est l'essence (la super).
\item L'organe est la \textbf{bougie}. Son rôle est de produire, à un instant précis (fin de la compression), l'étincelle qui enflamme le mélange air--essence.
\item Pendant l'admission, le cylindre aspire un \textbf{mélange d'air et d'essence}, préparé par le carburateur ou le système d'injection.
\item Ce type de moteur équipe les \textbf{moto-taxis} et les \textbf{voitures légères}, car il est léger et souple.
\end{enumerate}
\end{exoresolu}

\begin{exoresolu}[Le groupe électrogène du quartier]
\textbf{Énoncé.} Le groupe électrogène d'un débit de boissons n'a pas de bougie. Pourtant, le gasoil qu'on y verse s'enflamme bien dans le cylindre. Un élève affirme : « Sans étincelle, aucune combustion n'est possible, ce moteur est donc anormal. »
\begin{enumerate}
\item Explique pourquoi l'affirmation de l'élève est fausse.
\item Décris en deux étapes comment le gasoil s'enflamme dans ce moteur.
\end{enumerate}
\tcblower
\textbf{Solution détaillée.}
\begin{enumerate}
\item L'affirmation est fausse car il existe un second mode d'allumage : l'\textbf{allumage par compression}. Le groupe électrogène est un moteur \textbf{diesel} : il n'a pas besoin de bougie d'allumage.
\item Étape 1 : le cylindre admet de l'\textbf{air seul}, que le piston comprime très fortement ; d'après la propriété de la compression des gaz, cet air devient \textbf{très chaud}. Étape 2 : l'\textbf{injecteur} pulvérise alors le gasoil dans cet air brûlant ; le gasoil \textbf{s'enflamme spontanément} (auto-inflammation), sans aucune étincelle.
\end{enumerate}
\end{exoresolu}

\begin{exercice}[À toi de jouer]
\begin{enumerate}
\item Recopie et complète le tableau comparatif de la leçon (carburant, organe d'allumage, mélange admis, mode d'allumage, usages) sans regarder ton cours.
\item Pourquoi le taux de compression est-il plus élevé dans un moteur diesel ? Quelle propriété des gaz cela met-il en jeu ?
\item Un chauffeur de bus verse par erreur de l'essence dans le réservoir de son bus diesel. Que doit-il faire avant toute chose, et pourquoi ?
\item Cite deux opérations d'entretien spécifiques au moteur diesel et deux spécifiques au moteur à essence.
\item Pourquoi est-il dangereux de faire tourner un groupe électrogène dans une pièce fermée ? Nomme le gaz responsable.
\end{enumerate}
\end{exercice}

\begin{corrige}[Exercice n]
\begin{enumerate}
\item Voir le tableau comparatif de la section : essence/bougie/mélange air + essence/allumage commandé/motos et voitures légères ; gasoil/injecteur/air seul puis injection/allumage par compression/camions, bus, groupes électrogènes.
\item Il faut comprimer l'air très fortement pour que sa température devienne assez élevée pour enflammer spontanément le gasoil. Propriété mise en jeu : comprimer un gaz augmente sa température.
\item Il ne doit surtout pas démarrer le moteur, puis faire vidanger le réservoir par un garagiste. Démarrer enverrait l'essence dans un moteur conçu pour le gasoil et l'endommagerait gravement (grippage, casse injection).
\item Diesel : changer filtre à gasoil, purger le circuit gasoil (eau). Essence : changer bougies, vérifier/regler l'allumage (avance).
\item Danger : intoxication mortelle par le \textbf{monoxyde de carbone (\ce{CO})}, gaz inodore, incolore, très toxique, qui se fixe sur l'hémoglobine à la place de l'oxygène. Toujours placer le groupe à l'extérieur, loin des fenêtres.
\end{enumerate}
\end{corrige}

\begin{aretenir}[L'essentiel de la section]
\begin{itemize}
\item Essence et diesel sont tous deux des moteurs à combustion interne à quatre temps.
\item Moteur à \textbf{essence} : mélange air + essence admis, \textbf{allumage commandé} par l'étincelle de la \textbf{bougie} ; moteur léger et souple (motos, voitures).
\item Moteur \textbf{diesel} : air seul admis et fortement comprimé (donc très chaud), gasoil injecté qui s'enflamme spontanément : \textbf{allumage par compression}, sans bougie d'allumage ; moteur robuste et économique (camions, bus, groupes électrogènes).
\item Le taux de compression est plus élevé dans le moteur diesel.
\item Pour identifier un moteur sur un schéma : bougie $\rightarrow$ essence ; injecteur $\rightarrow$ diesel.
\item Sécurité : ne jamais inverser les carburants ; aérer les locaux (risque \ce{CO}) ; ne pas ouvrir le circuit de refroidissement à chaud.
\item Entretien : bougies + filtres (essence) ; filtre gasoil + purge (diesel) ; huile + courroies + liquide de refroidissement (commun).
\item Environnement : \ce{CO2} (effet de serre), \ce{CO} (toxique), \ce{NOx}/particules (pollution locale) ; catalyseur / FAP / vanne EGR pour dépolluer.
\end{itemize}
\end{aretenir}

\coursec{Sécurité, entretien et environnement}

Nous venons de voir que le moteur à essence et le moteur diesel, malgré leurs différences (allumage commandé ou auto-inflammation, bougie ou injecteur), reposent tous les deux sur un même principe : la \cle{combustion} d'un carburant dans un cylindre. Or brûler un carburant, rejeter des gaz d'échappement et faire tourner une machine chaude, cela comporte des \cle{dangers} réels. Chaque année, des accidents graves se produisent au Cameroun : incendies près des stations-service, intoxications au monoxyde de carbone pendant les délestages, brûlures lors de manipulations d'essence. Cette dernière section est donc peut-être la plus importante de toute la leçon : elle peut sauver des vies. Nous allons d'abord comprendre pourquoi les carburants sont dangereux, puis étudier les gaz d'échappement, les consignes de sécurité, l'entretien du moteur et enfin l'impact environnemental.

\courssub{Le triangle du feu : comprendre le danger des carburants}

\begin{definition}[Le triangle du feu]
Pour qu'une \cle{combustion} démarre et se poursuive, trois éléments doivent être réunis en même temps :
\begin{itemize}
\item un \cle{combustible} : la matière qui brûle (essence, gasoil, bois, gaz butane) ;
\item un \cle{comburant} : le gaz qui entretient la combustion, en général le \cle{dioxygène} de l'air ;
\item une source de \cle{chaleur} (ou énergie d'activation) : flamme, étincelle, surface très chaude.
\end{itemize}
Ces trois éléments forment le \cle{triangle du feu}. Si l'un des trois manque, le feu ne peut ni démarrer ni continuer.
\end{definition}

\begin{popfigure}
\begin{center}
\begin{tikzpicture}[scale=1.05, every node/.style={font=\small}]
% Triangle
\draw[line width=2pt, poporange] (0,4.2) -- (-3.2,-1.4) -- (3.2,-1.4) -- cycle;
% Sommets
\node[circle, draw=popink, fill=poporangeL, minimum size=2.1cm, align=center, line width=1.2pt] at (0,4.2) {\textbf{CHALEUR}\\ flamme, étincelle};
\node[circle, draw=popink, fill=popblueL, minimum size=2.1cm, align=center, line width=1.2pt] at (-3.2,-1.4) {\textbf{COMBUSTIBLE}\\ essence, gasoil};
\node[circle, draw=popink, fill=popgreenL, minimum size=2.1cm, align=center, line width=1.2pt] at (3.2,-1.4) {\textbf{COMBURANT}\\ dioxygène de l'air};
% Centre
\node[align=center, font=\bfseries\large, text=popdark] at (0,0.9) {FEU};
\draw[->, line width=1.5pt, poporange] (-0.5,1.6) -- (-1.3,2.6);
\draw[->, line width=1.5pt, poporange] (0.5,1.6) -- (1.3,2.6);
\draw[->, line width=1.5pt, poporange] (-0.9,0.2) -- (-1.9,-0.6);
\draw[->, line width=1.5pt, poporange] (0.9,0.2) -- (1.9,-0.6);
% Consignes
\node[draw=popink, fill=popcream, rounded corners, align=left, text width=4.6cm, font=\footnotesize] at (-5.6,2.6) {\textbf{Supprimer la chaleur :}\\ ne pas fumer, pas de flamme, moteur coupé à la station-service.};
\node[draw=popink, fill=popcream, rounded corners, align=left, text width=4.6cm, font=\footnotesize] at (-5.6,-1.6) {\textbf{Isoler le combustible :}\\ bidons fermés, stockés loin des habitations et des flammes.};
\node[draw=popink, fill=popcream, rounded corners, align=left, text width=4.6cm, font=\footnotesize] at (5.6,-1.6) {\textbf{Étouffer le comburant :}\\ couvrir un départ de feu avec un tissu mouillé ou du sable.};
\draw[->, popmuted] (-3.4,2.6) -- (-1.2,3.9);
\draw[->, popmuted] (-3.4,-1.6) -- (-4.3,-1.4);
\draw[->, popmuted] (3.4,-1.6) -- (4.3,-1.4);
\end{tikzpicture}
\legende{Le triangle du feu : une combustion exige la réunion d'un combustible, d'un comburant et d'une source de chaleur. Toute la sécurité consiste à séparer ces trois éléments.}
\end{center}
\end{popfigure}

\begin{experience}[Pourquoi les vapeurs d'essence s'enflamment-elles à distance ?]
\textbf{Matériel :} une petite coupelle contenant quelques gouttes d'essence, une longue tige métallique portant une bougie allumée, un local très ventilé (expérience réalisée uniquement par le professeur, avec toutes les protections).\par
\textbf{Protocole :} on approche lentement la flamme de la coupelle \emph{sans la toucher}, en restant à \SI{30}{cm} puis \SI{20}{cm} du liquide.\par
\textbf{Observation :} avant même que la flamme touche le liquide, une flamme bleue apparaît brusquement à la surface de l'essence.\par
\textbf{Interprétation :} l'essence \emph{s'évapore} facilement à température ambiante. Ses \cle{vapeurs}, invisibles, se mélangent à l'air et se déplacent. C'est ce mélange vapeurs d'essence + air qui s'enflamme au contact de la flamme, puis le feu « remonte » jusqu'au liquide.\par
\textbf{Conclusion :} le danger ne vient pas seulement du liquide, mais surtout de ses \cle{vapeurs invisibles}. C'est pourquoi une flamme peut mettre le feu à de l'essence \emph{à distance}.
\end{experience}

\begin{attention}[Vapeurs d'essence : un danger invisible]
Les vapeurs d'essence sont \cle{invisibles}, plus lourdes que l'air : elles s'accumulent au ras du sol, dans les garages, les fosses et les pièces mal aérées. Une simple étincelle (interrupteur, cigarette, frottement métallique) peut les enflammer à plusieurs mètres du bidon. Ne jamais manipuler d'essence près d'une flamme, d'une cigarette ou d'un appareil électrique en fonctionnement.
\end{attention}

\courssub{Les gaz d'échappement : le danger du monoxyde de carbone}

Lorsque la combustion est \emph{complète} (assez de dioxygène), le carburant brûle selon une équation du type :
\[ \ce{carburant + O2 -> CO2 + H2O} \]
Mais dans un moteur réel, surtout mal réglé ou dans un local fermé, la combustion est souvent \emph{incomplète} : il manque du dioxygène. Il se forme alors du \cle{monoxyde de carbone} \ce{CO} (équation simplifiée pour le carbone) :
\[ \ce{2C + O2 -> 2CO} \]

\begin{definition}[Le monoxyde de carbone]
Le \cle{monoxyde de carbone}, de formule \ce{CO}, est un gaz \cle{incolore} (on ne le voit pas), \cle{inodore} (on ne le sent pas) et \cle{très toxique}. Inspiré, il prend la place du dioxygène dans le sang (il se fixe sur l'hémoglobine) et empêche l'oxygénation des organes : il provoque des maux de tête, des vertiges, la perte de conscience, puis la mort en quelques minutes à forte dose.
\end{definition}

\begin{attention}[Jamais de moteur en marche dans un local fermé]
Le monoxyde de carbone \ce{CO} est un gaz \textbf{mortel, incolore et inodore} : on ne peut pas le détecter sans appareil. Un \cle{groupe électrogène} doit \textbf{toujours} fonctionner à l'extérieur, loin des portes et des fenêtres. Ne jamais le faire tourner dans une chambre, un salon, un garage ou une véranda, même avec la porte entrouverte. De même, on ne laisse jamais tourner le moteur d'une moto ou d'une voiture dans un garage fermé.
\end{attention}

\begin{exemplebox}[Les intoxications pendant les délestages]
Au Cameroun, lors des coupures d'électricité (délestages), de nombreuses familles utilisent un groupe électrogène. Chaque année, des cas d'\cle{intoxication au monoxyde de carbone} sont signalés lorsque l'appareil fonctionne dans une pièce fermée ou juste derrière une fenêtre ouverte : toute la famille s'endort et ne se réveille pas. La bonne pratique est simple : le groupe reste \emph{dehors}, sous un abri ventilé, à plusieurs mètres des ouvertures de la maison.
\end{exemplebox}

\courssub{Les consignes de sécurité pratiques}

\begin{popfigure}
\begin{center}
\begin{tikzpicture}[scale=1, every node/.style={font=\small}]
% Pictogramme 1 : ne pas fumer
\begin{scope}[shift={(-4.6,0)}]
\draw[line width=2pt, poppink] (0,0) circle (1.35);
\draw[line width=2pt, poppink] (-0.95,0.95) -- (0.95,-0.95);
% cigarette
\draw[fill=popcream, draw=popink, line width=0.8pt] (-0.75,-0.12) rectangle (0.45,0.12);
\draw[fill=poporange, draw=popink] (0.45,-0.12) rectangle (0.75,0.12);
\draw[popmuted, decorate, decoration={snake, amplitude=0.4mm, segment length=2mm}] (0.6,0.2) -- (0.6,0.75);
\node[align=center, font=\footnotesize\bfseries, text=popdark] at (0,-1.9) {Ne pas fumer\\ près du carburant};
\end{scope}
% Pictogramme 2 : moteur coupé
\begin{scope}[shift={(0,0)}]
\draw[line width=2pt, popblue] (0,0) circle (1.35);
% clé de contact
\draw[fill=popblueL, draw=popink, line width=0.8pt] (-0.55,-0.3) rectangle (0.55,0.3);
\draw[fill=popink] (0.15,-0.12) rectangle (0.45,0.12);
\draw[line width=1.2pt, popink] (-0.3,0.3) -- (-0.3,0.75);
\node[font=\footnotesize\bfseries, text=popdark] at (-0.3,0.95) {OFF};
\node[align=center, font=\footnotesize\bfseries, text=popdark] at (0,-1.9) {Moteur coupé\\ pendant le remplissage};
\end{scope}
% Pictogramme 3 : local ventilé
\begin{scope}[shift={(4.6,0)}]
\draw[line width=2pt, popgreen] (0,0) circle (1.35);
% fenêtre + flèches d'air
\draw[draw=popink, line width=0.8pt, fill=popgreenL] (-0.5,-0.55) rectangle (0.5,0.55);
\draw[draw=popink] (-0.5,0) -- (0.5,0);
\draw[draw=popink] (0,-0.55) -- (0,0.55);
\draw[->, line width=1.2pt, popgreen!60!black] (-1.05,0.35) -- (-0.55,0.35);
\draw[->, line width=1.2pt, popgreen!60!black] (0.55,-0.35) -- (1.05,-0.35);
\node[align=center, font=\footnotesize\bfseries, text=popdark] at (0,-1.9) {Groupe électrogène\\ dehors, local ventilé};
\end{scope}
\end{tikzpicture}
\legende{Trois pictogrammes de sécurité à connaître : interdiction de fumer près d'un carburant, moteur coupé pendant le remplissage du réservoir, groupe électrogène utilisé à l'extérieur dans un espace ventilé.}
\end{center}
\end{popfigure}

\begin{aretenir}[Les consignes d'or autour d'un moteur]
\begin{itemize}
\item \textbf{Interdiction de fumer} ou de faire du feu près d'une station-service, d'un bidon d'essence ou d'un réservoir ouvert.
\item \textbf{Couper le moteur} pendant le remplissage du réservoir à la station-service.
\item \textbf{Stocker le carburant} dans des \cle{bidons homologués} (prévus pour cela, bien fermés), loin des habitations, des cuisines et des flammes, jamais dans des bouteilles de boisson.
\item \textbf{Ne jamais faire tourner} un moteur ou un groupe électrogène dans un local fermé.
\item En cas de fuite d'essence : pas de flamme, pas d'interrupteur, on aère et on éloigne.
\end{itemize}
\end{aretenir}

\begin{methode}[Que faire face à une situation à risque ?]
Face à une fuite d'essence ou à un moteur qui tourne dans un garage fermé, applique la démarche en trois étapes :
\begin{enumerate}
\item \textbf{Identifier le danger} : vapeurs d'essence inflammables (risque d'incendie) ou gaz d'échappement contenant du \ce{CO} (risque d'intoxication).
\item \textbf{Couper la source} : arrêter le moteur ou le groupe électrogène, éteindre toute flamme, ne pas toucher aux interrupteurs électriques (une étincelle suffit).
\item \textbf{Aérer et évacuer} : ouvrir portes et fenêtres, faire sortir les personnes, ne revenir que lorsque l'air est renouvelé, puis réparer la fuite.
\end{enumerate}
\end{methode}

\begin{exoresolu}[Analyse d'une situation à risque]
\textbf{Énoncé.} Pendant un délestage à Yaoundé, M. Etoundi installe son groupe électrogène dans le couloir fermé de la maison « pour éviter la pluie ». Il remplit le réservoir moteur en marche, une cigarette à la main. Relever \textbf{trois} erreurs commises et donner pour chacune le danger encouru et la bonne pratique.
\tcblower
\textbf{Solution détaillée.}
\begin{itemize}
\item \emph{Erreur 1 : groupe dans un couloir fermé.} Danger : accumulation de monoxyde de carbone \ce{CO}, gaz mortel, incolore et inodore ; risque d'intoxication de toute la famille. Bonne pratique : installer le groupe \textbf{à l'extérieur}, sous un abri ventilé, loin des portes et fenêtres.
\item \emph{Erreur 2 : remplissage moteur en marche.} Danger : le moteur chaud et les étincelles électriques peuvent enflammer les vapeurs d'essence ; risque d'incendie et de brûlures. Bonne pratique : \textbf{couper le moteur} et le laisser refroidir avant le remplissage.
\item \emph{Erreur 3 : cigarette à la main.} Danger : la flamme enflamme les vapeurs d'essence \textbf{à distance}, car elles sont invisibles et se déplacent dans l'air. Bonne pratique : \textbf{ne jamais fumer} près d'un carburant.
\end{itemize}
Conclusion : ces trois erreurs violent le principe du triangle du feu (chaleur + combustible + comburant réunis) et exposent à l'incendie et à l'intoxication au \ce{CO}.
\end{exoresolu}

\courssub{L'entretien du moteur : prolonger sa durée de vie}

Un moteur est une machine qui s'use : l'huile se salit, le filtre à air s'encrasse, la bougie s'use. Un entretien régulier et simple permet de prolonger sa durée de vie et de réduire sa consommation.

\begin{center}
\begin{tabularx}{\linewidth}{|l|X|X|}\hline
\rowcolor{popblueL}
\textbf{Opération} & \textbf{Rôle} & \textbf{Si on néglige} \\ \hline
Vidange de l'huile & L'huile lubrifie les pièces en mouvement (piston, vilebrequin) et limite l'usure par frottement. & Huile sale $\Rightarrow$ frottements $\Rightarrow$ usure rapide, surchauffe, casse du moteur. \\ \hline
Nettoyage ou remplacement du filtre à air & Le filtre retient la poussière de l'air aspiré (très chargé sur nos pistes en saison sèche). & Filtre encrassé $\Rightarrow$ moins d'air $\Rightarrow$ combustion incomplète $\Rightarrow$ surconsommation et fumées noires. \\ \hline
Vérification de la bougie (moteur à essence) & La bougie produit l'étincelle qui déclenche la combustion. & Bougie usée $\Rightarrow$ mauvais allumage $\Rightarrow$ à-coups, démarrages difficiles, consommation excessive. \\ \hline
\end{tabularx}
\end{center}

\begin{savaistu}[L'entretien : un geste économique et écologique]
Un moteur mal réglé — filtre à air encrassé, bougie usée — consomme \textbf{plus} de carburant pour la même distance et pollue \textbf{davantage}, car la combustion devient incomplète : il rejette plus de monoxyde de carbone \ce{CO} et de fumées noires (imbrûlés). Entretenir sa moto ou son groupe électrogène, c'est donc à la fois économiser l'argent du carburant et protéger la santé de tous. Les mécaniciens de quartier le savent bien : « un moteur qui fume noir mange l'essence ».
\end{savaistu}

\courssub{Impact environnemental et notion de rendement}

Les gaz d'échappement des moteurs contiennent du dioxyde de carbone \ce{CO2} (gaz à effet de serre), du monoxyde de carbone \ce{CO}, des oxydes d'azote et des particules (fumées noires). Dans les grandes villes comme \cle{Douala} et \cle{Yaoundé}, la multiplication des moto-taxis, des voitures et des groupes électrogènes dégrade la qualité de l'air et provoque des maladies respiratoires. Un parc de véhicules bien entretenus et bien réglés rejette moins de polluants : c'est un enjeu de \cle{santé publique}.

\begin{definition}[Rendement d'un moteur (complément hors strict programme)]
Le \cle{rendement} d'un moteur est le rapport entre l'énergie mécanique utile fournie et l'énergie chimique contenue dans le carburant :
\[ \eta = \frac{E_{\text{utile}}}{E_{\text{carburant}}} \]
Pour un moteur à combustion interne, ce rendement est faible : environ $25\,\%$ à $35\,\%$. Une grande partie de l'énergie du carburant est \cle{perdue en chaleur} (moteur chaud, gaz d'échappement brûlants) et en frottements.
\end{definition}

\begin{exoresolu}[Calcul de l'énergie réellement utilisée]
\textbf{Énoncé.} La combustion de \SI{1}{L} d'essence libère une énergie d'environ $E_{\text{carburant}} = \SI{34}{MJ}$. Le moteur d'une moto a un rendement $\eta = 30\,\%$. Calculer l'énergie mécanique utile réellement obtenue, puis l'énergie perdue (en chaleur et frottements).
\tcblower
\textbf{Solution détaillée.} On applique la définition du rendement :
\begin{align*}
\eta &= \frac{E_{\text{utile}}}{E_{\text{carburant}}} \\
E_{\text{utile}} &= \eta \times E_{\text{carburant}} = 0{,}30 \times 34 = \SI{10.2}{MJ}
\end{align*}
L'énergie perdue est :
\[ E_{\text{perdue}} = 34 - 10{,}2 = \SI{23.8}{MJ} \]
Conclusion : sur \SI{1}{L} d'essence, seulement \SI{10.2}{MJ} font avancer la moto ; près de $70\,\%$ de l'énergie part en chaleur dans l'atmosphère. C'est pourquoi un moteur bien réglé, qui gaspille moins, est à la fois économique et écologique.
\end{exoresolu}

\begin{attention}[Pièges fréquents au BEPC]
\begin{itemize}
\item Ne pas confondre \ce{CO} (monoxyde de carbone, \textbf{toxique}, issu d'une combustion \emph{incomplète}) et \ce{CO2} (dioxyde de carbone, produit de la combustion \emph{complète}, gaz à effet de serre mais non toxique à faible dose).
\item Dire « l'essence brûle » est imprécis : ce sont ses \textbf{vapeurs} mélangées à l'air qui s'enflamment.
\item Le rendement s'exprime sans unité ou en pourcentage ; il est toujours \textbf{inférieur à 1} (ou à $100\,\%$).
\end{itemize}
\end{attention}

\begin{exercice}[Pour s'entraîner]
\begin{enumerate}
\item Citer les trois éléments du triangle du feu et expliquer pourquoi on couvre un départ de feu d'huile avec un couvercle plutôt qu'avec de l'eau.
\item Pourquoi est-il dangereux de faire tourner un groupe électrogène dans une pièce fermée ? Nommer le gaz responsable et donner deux de ses propriétés.
\item Citer trois opérations d'entretien d'un moteur et expliquer l'intérêt de chacune.
\item Un moteur reçoit \SI{20}{MJ} d'énergie du carburant et fournit \SI{5}{MJ} d'énergie mécanique. Calculer son rendement.
\end{enumerate}
\end{exercice}

\begin{corrige}[Exercice n]
\begin{enumerate}
\item Combustible, comburant (dioxygène), chaleur. Le couvercle étouffe le feu en supprimant le comburant ; l'eau sur de l'huile brûlante provoque des projections de gouttelettes enflammées et aggrave l'incendie.
\item Le moteur rejette du monoxyde de carbone \ce{CO}, qui s'accumule dans la pièce. Propriétés : incolore, inodore (et très toxique). Il peut tuer les occupants sans qu'ils s'en aperçoivent.
\item Vidange de l'huile (lubrifier, limiter l'usure) ; nettoyage du filtre à air (garantir une bonne combustion, éviter la surconsommation) ; vérification de la bougie (assurer l'allumage, éviter les à-coups).
\item $\eta = \dfrac{E_{\text{utile}}}{E_{\text{carburant}}} = \dfrac{5}{20} = 0{,}25$, soit un rendement de $25\,\%$.
\end{enumerate}
\end{corrige}

\begin{aretenir}[L'essentiel de la section]
\begin{itemize}
\item Une combustion exige le \cle{triangle du feu} : combustible + comburant + chaleur ; la sécurité consiste à les séparer.
\item Les \cle{vapeurs d'essence} sont invisibles et s'enflamment à distance d'une flamme : pas de cigarette, moteur coupé au remplissage, bidons homologués loin des habitations.
\item Les gaz d'échappement contiennent du \cle{monoxyde de carbone} \ce{CO}, gaz mortel, incolore et inodore : jamais de moteur ou de groupe électrogène dans un local fermé.
\item L'\cle{entretien} (vidange, filtre à air, bougie) prolonge la vie du moteur, économise le carburant et réduit la pollution.
\item Le \cle{rendement} d'un moteur est faible (environ $25$ à $35\,\%$) : l'essentiel de l'énergie du carburant est perdu en chaleur.
\end{itemize}
\end{aretenir}

\newpage

\coursec{Activité d'intégration}

\coursec{Leçon 19 : Fonctionnement d'un moteur à combustion interne}

\courssub{Objectifs de la leçon}
À la fin de cette leçon, tu seras capable de :
\begin{itemize}
\item identifier sur un schéma en coupe les principaux organes d'un moteur à combustion interne et rappeler le rôle de chacun ;
\item remettre dans l'ordre les quatre temps du cycle moteur en justifiant par la position du piston et l'état des soupapes ;
\item distinguer un moteur à essence d'un moteur diesel et expliquer les conséquences d'une erreur de carburant ;
\item diagnostiquer une panne simple ou un danger lié à l'utilisation d'un moteur, puis proposer une solution et des consignes de sécurité ;
\item travailler en groupe, argumenter oralement et rédiger une fiche de conseils.
\end{itemize}

\courssub{Situation déclenchante : Au garage « Chez Papa Mballa »}
Papa Mballa tient depuis vingt ans un petit garage au quartier Nkolndongo, à Yaoundé. Chaque matin, des motos-taxis, des taxis et même des bus de voyage viennent s'y faire réparer. Ce samedi, il est débordé : trois clients arrivent en même temps, et son apprenti est absent. Il te fait confiance, toi l'élève de 3ème, pour l'aider à poser les diagnostics.

\begin{experience}Observation d'un moteur en coupe (Document 1)\end{experience}
Sur l'établi, Papa Mballa a posé la vue en coupe d'un moteur à essence à quatre temps. Six organes doivent être identifiés : le \cle{cylindre}, le \cle{piston}, les \cle{soupapes} (d'admission et d'échappement), la \cle{bougie}, la \cle{bielle} et le \cle{vilebrequin}.

\begin{popfigure}
\begin{center}
\begin{tikzpicture}[scale=0.95, every node/.style={font=\small}]
% Cylindre (murs et culasse)
\draw[thick, popink, fill=popcream] (-1.5,0) rectangle (1.5,3.5);
% Culasse (partie haute)
\draw[thick, popink, fill=popcream] (-1.8,3.5) rectangle (1.8,4.2);
% Piston
\draw[fill=popblueL, draw=popblue, thick] (-1.0,1.2) rectangle (1.0,1.8);
% Axe piston (point de liaison bielle)
\draw[fill=popdark] (0,1.2) circle (0.08);
% Bielle
\draw[very thick, poporange] (0,1.2) -- (1.0,-1.0);
% Vilebrequin : palier principal (centre) + maneton (excentrique) + bras
\draw[thick, popdark, fill=popgoldL] (0,-1.5) circle (0.6); % Contour vilebrequin simplifié
\draw[thick, popdark, fill=popdark] (0,-1.5) circle (0.15); % Palier principal
\draw[very thick, popdark] (0,-1.5) -- (1.0,-1.0); % Bras de maneton
\draw[thick, popdark, fill=popgoldL] (1.0,-1.0) circle (0.18); % Maneton (axe bielle)
% Soupape admission (gauche)
\draw[thick, popgreen] (-0.6,3.5) -- (-0.6,4.8);
\draw[thick, popgreen, fill=popgreenL] (-0.9,3.5) -- (-0.3,3.5) -- (-0.6,3.2) -- cycle; % Tête soupape
% Soupape échappement (droite)
\draw[thick, popgreen] (0.6,3.5) -- (0.6,4.8);
\draw[thick, popgreen, fill=popgreenL] (0.9,3.5) -- (0.3,3.5) -- (0.6,3.2) -- cycle; % Tête soupape
% Ressorts de soupapes (schématisés)
\draw[popmuted, decorate, decoration={coil, aspect=0.3, segment length=3pt, amplitude=1.5pt}] (-0.6,3.7) -- (-0.6,4.5);
\draw[popmuted, decorate, decoration={coil, aspect=0.3, segment length=3pt, amplitude=1.5pt}] (0.6,3.7) -- (0.6,4.5);
% Bougie (centre culasse)
\draw[thick, poppurple] (0,4.2) -- (0,5.0);
\draw[thick, poppurple] (-0.15,5.0) -- (0.15,5.0); % Écrou
\draw[poppurple, thick] (-0.1,3.45) -- (0,3.25) -- (0.1,3.45); % Électrodes
% Flèches d'étiquettes
\node[popblue, font=\bfseries\large] at (1.7,1.5) {1};
\node[poporange, font=\bfseries\large] at (1.6,-0.3) {2};
\node[popdark, font=\bfseries\large] at (-0.5,-2.4) {3};
\node[popgreen, font=\bfseries\large] at (-1.6,4.2) {4};
\node[poppurple, font=\bfseries\large] at (0.5,5.3) {5};
\node[popink, font=\bfseries\large] at (-2.3,2.0) {6};
\end{tikzpicture}
\end{center}
\legende{Vue en coupe simplifiée d'un moteur essence 4 temps. Les organes sont repérés par les numéros 1 à 6.}
\end{popfigure}

\courssub{Activité 1 : Les organes du moteur et leurs rôles}
\begin{enumerate}
\item Nomme les organes repérés de 1 à 6 sur le schéma ci-dessus.
\item Pour chaque organe, donne son rôle précis dans le fonctionnement du moteur.
\end{enumerate}

\courssub{Activité 2 : Le cycle à quatre temps (Document 2)}
L'apprenti avait dessiné les quatre temps du cycle moteur sur des fiches, mais elles se sont mélangées :
\begin{itemize}
\item \textbf{Fiche A :} Les deux soupapes sont fermées, l'étincelle jaillit à la bougie, le piston est brutalement poussé vers le bas.
\item \textbf{Fiche B :} La soupape d'admission est ouverte, le piston descend, le mélange air-essence entre dans le cylindre.
\item \textbf{Fiche C :} La soupape d'échappement est ouverte, le piston monte et chasse les gaz brûlés.
\item \textbf{Fiche D :} Les deux soupapes sont fermées, le piston monte et comprime le mélange.
\end{itemize}

\begin{enumerate}
\setcounter{enumi}{2}
\item Remets les fiches A, B, C, D dans l'ordre chronologique du cycle à quatre temps. Nomme chaque temps.
\item Justifie ton classement en utilisant la position du piston (monte/descend) et l'état des soupapes (ouvertes/fermées).
\item Combien de tours le vilebrequin effectue-t-il pendant un cycle complet de quatre temps ? À quel temps le moteur produit-il réellement de l'énergie mécanique ?
\end{enumerate}

\begin{aretenir}[Le cycle à quatre temps et les organes essentiels]
\begin{itemize}
\item Un \cle{moteur à combustion interne} transforme l'énergie chimique du carburant en énergie mécanique : la combustion a lieu \emph{à l'intérieur} du cylindre.
\item \textbf{Organes et rôles :}
\begin{itemize}
\item \cle{Cylindre} : guide le piston et abrite la chambre de combustion.
\item \cle{Piston} : coulisse dans le cylindre ; reçoit la pression des gaz et transmet la force.
\item \cle{Bielle} : relie le piston au vilebrequin ; transforme le mouvement alternatif en rotation.
\item \cle{Vilebrequin} : transforme le mouvement de va-et-vient (alternatif) du piston en mouvement de rotation continu.
\item \cle{Soupapes} : clapets commandés par l'arbre à cames. \emph{Soupape d'admission} : laisse entrer le mélange (ou l'air). \emph{Soupape d'échappement} : laisse sortir les gaz brûlés.
\item \cle{Bougie} (essence) : produit l'étincelle d'allumage au bon moment.
\end{itemize}
\item \textbf{Cycle à 4 temps (ordre immuable) :}
\begin{enumerate}
\item \textbf{Admission} : Soupape d'admission ouverte, piston descend $\rightarrow$ remplissage.
\item \textbf{Compression} : Soupapes fermées, piston monte $\rightarrow$ compression du mélange (augmentation $T$ et $P$).
\item \textbf{Combustion-Détente (Temps moteur)} : Soupapes fermées, allumage (bougie ou auto-inflammation), piston poussé violemment vers le bas $\rightarrow$ \textbf{seul temps producteur d'énergie}.
\item \textbf{Échappement} : Soupape d'échappement ouverte, piston monte $\rightarrow$ évacuation des gaz.
\end{enumerate}
\item Un cycle complet = \textbf{2 tours de vilebrequin} = 4 courses du piston (2 montées, 2 descentes).
\end{itemize}
\end{aretenir}

\begin{methode}Reconstituer le cycle sur un schéma\end{methode}
\begin{enumerate}
\item Repère l'état des soupapes : \textbf{Admission ouverte} = Admission ; \textbf{Échappement ouverte} = Échappement ; \textbf{Les deux fermées} = Compression ou Combustion.
\item Si les deux sont fermées, regarde la bougie (étincelle) ou le sens de la flèche sur le piston : \textbf{Étincelle + Piston descend} = Combustion-Détente ; \textbf{Piston monte} = Compression.
\item Vérifie l'ordre : Admission $\rightarrow$ Compression $\rightarrow$ Combustion-Détente $\rightarrow$ Échappement.
\end{enumerate}

\courssub{Activité 3 : Moteur essence, moteur diesel et sécurité (Document 3)}
Trois clients attendent :
\begin{itemize}
\item \textbf{Client 1 — Moussa (moto-taxi) :} Sa moto « cale » dès qu'il accélère. Diagnostic : le filtre à air est complètement encrassé de poussière rouge.
\item \textbf{Client 2 — Amadou (bus) :} Il a mis \SI{20}{L} d'essence dans le réservoir de son bus \textbf{diesel}. Le moteur refuse de démarrer.
\item \textbf{Client 3 — Madame Etoundi :} Pendant une coupure ENEO, le groupe électrogène a tourné dans le couloir portes fermées. La famille a eu maux de tête et nausées : intoxication au \ce{CO}.
\end{itemize}

\begin{enumerate}
\setcounter{enumi}{5}
\item \textbf{Client 1.} À quel temps du cycle le filtre à air intervient-il ? Explique pourquoi un filtre encrassé fait caler le moteur. Propose une solution.
\item \textbf{Client 2.} Rappelle la différence fondamentale d'allumage entre moteur essence et moteur diesel. Explique pourquoi l'essence ne convient pas au moteur diesel et indique la procédure à suivre.
\item \textbf{Client 3.} Quel gaz toxique est produit par la combustion incomplète ? Pourquoi est-il dangereux ? Rédige deux consignes de sécurité à afficher dans le garage.
\item Rédige collectivement une \textbf{fiche de conseils} (5 lignes max) intitulée « Bien utiliser son moteur », destinée aux clients de Papa Mballa.
\end{enumerate}

\begin{attention}[Erreurs fréquentes et pièges du BEPC]
\begin{itemize}
\item \textbf{Confusion des temps :} Ne pas inverser Compression et Combustion-Détente. Rappelle-toi : \emph{Compression = piston monte} ; \emph{Combustion = piston descend (poussé)}.
\item \textbf{Bougie vs Injecteur :} Le moteur diesel n'a \textbf{pas de bougie}. L'allumage se fait par \emph{auto-inflammation} du gasoil injecté dans l'air très chaud (forte compression). Ne pas écrire « bougie diesel ».
\item \textbf{Nombre de tours :} Un cycle 4 temps = \textbf{2 tours} de vilebrequin (pas 1, pas 4).
\item \textbf{Monoxyde de carbone \ce{CO} :} Gaz \emph{incolore, inodore, non irritant} $\rightarrow$ impossible à détecter sans détecteur. Ne jamais faire tourner un moteur (moto, groupe, voiture) dans un local fermé, même porte ouverte.
\item \textbf{Erreur carburant :} Essence dans diesel $\rightarrow$ casse pompe à injection (lubrification perdue) + cliquetis. Diesel dans essence $\rightarrow$ calage, fumée noire, bougies encrassées. \textbf{Ne jamais démarrer} : vidange immédiate.
\end{itemize}
\end{attention}

\begin{savaistu}[Contexte camerounais et histoire]
\begin{itemize}
\item \textbf{SONARA (Limbe) :} Raffine le pétrole brut pour produire l'essence (super) et le gasoil (diesel) consommés au Cameroun. La qualité du carburant influence l'encrassement du moteur (bougies, injecteurs, filtres).
\item \textbf{Nikolaus Otto (1876) :} Inventeur du cycle à 4 temps essence (allumage commandé). \textbf{Rudolf Diesel (1892) :} Inventeur du moteur à allumage par compression (diesel), plus performant (meilleur rendement) pour les gros porteurs (bus, camions, groupes électrogènes).
\item \textbf{Moto-taxis « Bend Skin » :} Majoritairement des moteurs 4 temps essence monocylindres. L'entretien du filtre à air est crucial en saison sèche (poussière rouge latéritique) et en saison des pluies (boue/eau).
\item \textbf{Groupes électrogènes :} Très utilisés lors des délestages ENEO. Danger majeur : \ce{CO}. Installer \emph{toujours} à l'extérieur, loin des fenêtres/portes, sortie d'échappement dirigée vers l'extérieur.
\end{itemize}
\end{savaistu}

\begin{exoresolu}[Corrigé détaillé : Diagnostic au garage « Chez Papa Mballa »]
Énoncé : Identifier les organes, reconstituer le cycle, diagnostiquer les trois cas clients, rédiger la fiche de conseils.
\tcblower
\textbf{Étape 1 — Légende du schéma (Question 1).}
\begin{itemize}
\item \textbf{1 : Piston.} Pièce mobile cylindrique qui coulisse dans le cylindre. Il reçoit la poussée des gaz brûlés lors de la combustion et la transmet à la bielle.
\item \textbf{2 : Bielle.} Organe de liaison articulé (tête de bielle sur piston, pied de bielle sur vilebrequin). Transforme le mouvement alternatif du piston en mouvement de rotation du vilebrequin.
\item \textbf{3 : Vilebrequin (vilebrequin).} Arbre à manetons excentrés. Reçoit le mouvement de la bielle et délivre un mouvement de rotation continu vers la boîte de vitesses / roue.
\item \textbf{4 : Soupape d'admission (gauche) / Soupape d'échappement (droite).} Clapets commandés par l'arbre à cames. Admission : ouvre le passage au mélange frais (ou air). Échappement : ouvre le passage aux gaz brûlés.
\item \textbf{5 : Bougie.} Organe d'allumage (moteur essence). Produit une étincelle haute tension entre ses électrodes pour enflammer le mélange comprimé.
\item \textbf{6 : Cylindre (corps de cylindre + culasse).} Pièce fixe en fonte ou aluminium. Sa paroi interne (chemise) guide le piston ; sa partie haute (culasse) ferme la chambre de combustion et porte soupapes et bougie.
\end{itemize}
\end{exoresolu}

\textbf{Étape 2 — Remise en ordre du cycle (Questions 2, 3, 4).}
L'ordre chronologique est : \textbf{B $\rightarrow$ D $\rightarrow$ A $\rightarrow$ C}.
\begin{center}
\begin{tabularx}{\linewidth}{|l|X|X|X|}\hline
\textbf{Temps} & \textbf{Fiche} & \textbf{Soupapes} & \textbf{Piston / Action} \\ \hline
1. Admission & B & Admission ouverte, Échappement fermée & Descend $\rightarrow$ aspire le mélange \\ \hline
2. Compression & D & \textbf{Les deux fermées} & Monte $\rightarrow$ comprime le mélange \\ \hline
3. Combustion-Détente & A & \textbf{Les deux fermées} & \textbf{Descend violemment} (poussé par l'explosion) \\ \hline
4. Échappement & C & Admission fermée, Échappement ouverte & Monte $\rightarrow$ chasse les gaz brûlés \\ \hline
\end{tabularx}
\end{center}
\begin{itemize}
\item \textbf{Justification :} L'admission (B) remplit le cylindre. La compression (D) prépare le mélange (augmentation $P, T$). La combustion (A) est le seul temps où la pression pousse le piston vers le bas : c'est le \textbf{temps moteur}. L'échappement (C) vide le cylindre pour le cycle suivant.
\item \textbf{Tours de vilebrequin :} Un cycle complet (4 temps = 4 courses) correspond à \textbf{2 tours complets} du vilebrequin.
\item \textbf{Énergie :} Seul le \textbf{3ème temps (Combustion-Détente)} fournit de l'énergie mécanique. Les 3 autres temps (admission, compression, échappement) consomment de l'énergie (fournie par l'inertie du volant moteur).
\end{itemize}

\textbf{Étape 3 — Client 1 : Filtre à air encrassé (Question 5).}
\begin{itemize}
\item \textbf{Temps concerné :} \textbf{Admission (1er temps)}. Le filtre est placé sur le conduit d'admission, avant le carburateur / l'injection.
\item \textbf{Explication :} Le filtre retient la poussière. Encrassé, il crée une perte de charge importante : le débit d'air aspiré chute. Le mélange devient \textbf{trop riche} (excès d'essence par rapport à l'air). La combustion est incomplète, la puissance chute, le moteur « suffoque » et cale à l'accélération.
\item \textbf{Solution :} Nettoyer le filtre (souffler à l'air comprimé ou laver selon type) ou le \textbf{remplacer}. Vérifier à chaque révision (tous les 5 000--10 000 km ou saison sèche).
\end{itemize}

\textbf{Étape 4 — Client 2 : Erreur carburant Essence dans Diesel (Question 6).}
\begin{itemize}
\item \textbf{Différence allumage :}
\begin{itemize}
\item \textbf{Essence (allumage commandé) :} Mélange air+essence comprimé (taux ~10:1) $\rightarrow$ étincelle \textbf{bougie}.
\item \textbf{Diesel (allumage par compression / auto-inflammation) :} Air seul comprimé très fort (taux ~18--23:1) $\rightarrow$ $T > \SI{500}{\celsius}$ $\rightarrow$ injection gasoil $\rightarrow$ \textbf{inflammation spontanée}. \textbf{Pas de bougie}.
\end{itemize}
\item \textbf{Pourquoi ça ne démarre pas :} L'essence est trop volatile et a un indice de cétane très bas (résiste mal à l'auto-inflammation). Dans le diesel, elle ne s'enflamme pas au bon moment (ou pas du tout) sous l'effet de la compression. De plus, l'essence n'a pas de pouvoir lubrifiant $\rightarrow$ \textbf{risque de grippage de la pompe à injection haute pression} (pièces métalliques en contact sans film d'huile).
\item \textbf{Procédure :} \textbf{Ne surtout pas insister au démarreur} (risque d'envoyer de l'essence dans tout le circuit d'injection). \textbf{Vidanger entièrement le réservoir}, purger le circuit (filtre, tuyaux, pompe, injecteurs), changer le filtre à gasoil, puis remplir de gasoil.
\end{itemize}

\textbf{Étape 5 — Client 3 : Intoxication au \ce{CO} (Question 7).}
\begin{itemize}
\item \textbf{Gaz :} \cle{Monoxyde de carbone} \ce{CO}. Produit par \textbf{combustion incomplète} (manque d'oxygène : $\ce{C_xH_y + O2 -> CO + H2O + C}$ au lieu de $\ce{CO2}$).
\item \textbf{Danger :} Gaz \textbf{incolore, inodore, non irritant}. Se fixe sur l'hémoglobine du sang (forme \ce{COHb}, carboxyhémoglobine) avec une affinité \textbf{200 à 250 fois supérieure} à celle de l'\ce{O2}. Bloque le transport de l'oxygène $\rightarrow$ hypoxie tissulaire. Symptômes : maux de tête, nausées, vertiges $\rightarrow$ perte de conscience $\rightarrow$ mort.
\item \textbf{Consignes de sécurité (à afficher) :}
\begin{enumerate}
\item \textbf{Interdiction formelle} de faire fonctionner un moteur thermique (moto, groupe, voiture) dans un local fermé (garage, couloir, chambre), même portes/fenêtres ouvertes.
\item Installer le groupe électrogène \textbf{uniquement à l'extérieur}, à l'air libre, à plus de \SI{2}{m} des ouvertures (portes, fenêtres, ventilations), sortie d'échappement dirigée vers l'extérieur.
\end{enumerate}
\end{itemize}

\textbf{Étape 6 — Fiche de conseils (Question 8), proposition.}
\begin{center}
\fbox{\parbox{0.9\linewidth}{\centering
\textbf{Bien utiliser son moteur} \\
1. Nettoie ou change le filtre à air régulièrement (poussière rouge !). \\
2. Vérifie le carburant \textbf{avant} le plein : Essence ou Gasoil ? \\
3. Groupe électrogène \textbf{uniquement dehors}, loin des fenêtres. \\
4. Entretiens : bougie, huile, filtre à huile, niveau liquide de refroidissement. \\
5. Malaise / odeur suspecte ? \textbf{Sors immédiatement}, appelle les secours.
}}
\end{center}

\courssub{Bilan de la leçon}
Cette leçon t'a permis de mettre en évidence que :
\begin{itemize}
\item le fonctionnement d'un moteur à combustion interne repose sur l'action coordonnée d'organes précis : piston, bielle, vilebrequin, soupapes et système d'allumage (bougie ou injection) ;
\item le cycle à quatre temps (admission, compression, combustion-détente, échappement) s'enchaîne toujours dans le même ordre ; seul le temps de combustion-détente fournit l'énergie mécanique ; un cycle = 2 tours de vilebrequin ;
\item essence et diesel ne sont pas interchangeables : modes d'allumage différents (commandé vs compression), taux de compression différents, organes différents (bougie vs injecteur/pompe HP) ;
\item la mécanique est une démarche logique : observer les symptômes $\rightarrow$ identifier le temps du cycle concerné $\rightarrow$ proposer une solution technique ;
\item l'utilisation d'un moteur impose des règles de sécurité vitales, notamment contre le danger invisible et mortel du monoxyde de carbone \ce{CO}.
\end{itemize}
Tu es maintenant capable, comme un vrai technicien, de raisonner sur un moteur : c'est exactement la démarche attendue au BEPC face à une situation-problème de technologie.

\begin{exercice}Entraînement BEPC : Moteur et Sécurité\end{exercice}
\begin{enumerate}
\item \textbf{QCM.} Lors du temps de compression d'un moteur 4 temps essence :
\begin{itemize}
\item[A.] La soupape d'admission est ouverte et le piston descend.
\item[B.] Les deux soupapes sont fermées et le piston monte.
\item[C.] La soupape d'échappement est ouverte et le piston monte.
\item[D.] Les deux soupapes sont fermées et le piston descend.
\end{itemize}
\item \textbf{Schéma.} Sur la figure ci-dessous représentant un piston au point mort haut (PMH) en fin de compression, trace la flèche indiquant le sens de rotation du vilebrequin et la flèche indiquant la force exercée sur le piston juste après l'allumage.
\begin{popfigure}
\begin{center}
\begin{tikzpicture}[scale=1.2]
% Cylindre
\draw[thick, popink] (-1,0) -- (-1,3) -- (1,3) -- (1,0);
% Piston au PMH
\draw[fill=popblueL, draw=popblue, thick] (-0.9,2.5) rectangle (0.9,2.9);
% Bielle verticale (PMH)
\draw[very thick, poporange] (0,2.5) -- (0,1.5);
% Vilebrequin (palier bas)
\draw[thick, popdark, fill=popgoldL] (0,1.0) circle (0.5);
\draw[fill=popdark] (0,1.0) circle (0.1);
% Maneton au PMH (au-dessus du palier)
\draw[fill=popgoldL] (0,1.5) circle (0.15);
% Bougie
\draw[thick, poppurple] (0,3) -- (0,3.6);
\draw[thick, poppurple] (-0.1,3.6) -- (0.1,3.6);
\draw[poppurple, thick] (-0.08,2.95) -- (0,2.8) -- (0.08,2.95);
% Soupapes fermées
\draw[thick, popgreen, fill=popgreenL] (-0.7,3) -- (-0.3,3) -- (-0.5,2.7) -- cycle;
\draw[thick, popgreen, fill=popgreenL] (0.7,3) -- (0.3,3) -- (0.5,2.7) -- cycle;
\end{tikzpicture}
\end{center}
\legende{Piston au Point Mort Haut (fin de compression).}
\end{popfigure}
\item \textbf{Calcul.} Un moteur de moto tourne à \SI{6000}{tr/min}. Combien de cycles complets (4 temps) effectue-t-il par seconde ?
\item \textbf{Rédaction.} Un voisin fait tourner son groupe électrogène dans son garage attenant à la chambre, porte fermée. Tu sens une odeur d'échappement. Explique le danger précis (gaz, mécanisme physiologique) et ce que tu dois faire immédiatement.
\end{enumerate}

\begin{corrige}Exercice d'entraînement\end{corrige}
\begin{enumerate}
\item \textbf{B.} Compression = 2ème temps. Soupapes fermées pour étanchéifier la chambre, piston monte pour comprimer le mélange.
\item \textbf{Schéma :}
\begin{itemize}
\item Force sur piston : Flèche verticale \textbf{vers le bas} (poussée des gaz brûlés).
\item Rotation vilebrequin : Le maneton est au PMH (au-dessus du palier). La poussée sur la bielle (verticale) ne crée pas de couple au PMH strict, mais immédiatement après, le piston descend, la bielle pousse le maneton vers l'avant (si rotation horaire) ou l'arrière. \textbf{Convention usuelle : rotation horaire.} Flèche circulaire horaire sur le vilebrequin.
\end{itemize}
\item \textbf{Calcul :} $N = \SI{6000}{tr/min} = \frac{6000}{60} = \SI{100}{tr/s}$.
1 cycle = 2 tours $\rightarrow$ Nombre de cycles/s $= \frac{100}{2} = \mathbf{50~cycles/s}$ (ou \SI{50}{Hz}).
\item \textbf{Danger :} Le garage fermé $\rightarrow$ manque d'\ce{O2} $\rightarrow$ combustion incomplète $\rightarrow$ production de \ce{CO} (monoxyde de carbone). Le \ce{CO} se fixe sur l'hémoglobine à la place de l'\ce{O2} (affinité 200-250x supérieure) $\rightarrow$ asphyxie cellulaire (hypoxie). Risque de mort rapide (quelques minutes à forte concentration).
\textbf{Action immédiate :} 1. \textbf{Sortir la personne} (et soi-même) \textbf{immédiatement} à l'air libre. 2. Couper le groupe (si accessible sans risque). 3. Appeler les secours (112 / 119). 4. Aérer le local longtemps avant d'y retourner.
\end{enumerate}

\newpage

\coursec{Méthodes et exercices résolus}

\coursec{Leçon 19 : Fonctionnement d'un moteur à combustion interne}

\courssub{Le moteur à combustion interne : définition et principe}

\begin{definition}[Moteur à combustion interne]
Un \cle{moteur à combustion interne} est une machine thermique qui transforme l'énergie chimique d'un carburant (essence ou gasoil) en énergie mécanique (mouvement de rotation) par la combustion de ce carburant \textbf{à l'intérieur même} du moteur, dans une chambre de combustion.
\end{definition}

\begin{savaistu}[Naissance du moteur moderne]
En 1876, l'Allemand \textbf{Nikolaus Otto} met au point le premier moteur à quatre temps fonctionnant à l'essence (cycle Otto). En 1892, \textbf{Rudolf Diesel} brevette le moteur à allumage par compression (cycle Diesel), plus robuste et plus efficace, qui équipera camions, bateaux et groupes électrogènes. Au Cameroun, la \textbf{SONARA} (Société Nationale de Raffinage) de Limbé transforme le pétrole brut en essence et en gasoil qui alimentent nos moto-taxis, voitures, camions et groupes électrogènes ENEO.
\end{savaistu}

\begin{popfigure}
\begin{tikzpicture}[scale=0.9, every node/.style={font=\small}]
% Cylindre
\draw[thick, popblue] (0,0) -- (0,5) -- (3,5) -- (3,0) -- (4,0) -- (4,-0.5) -- (3,-0.5) -- (3,0) -- (0,0);
% Piston
\draw[thick, poporange, fill=poporange!30] (0.2,2.5) rectangle (2.8,3.2);
% Bielle
\draw[thick, popdark] (1.5,2.5) -- (1.5,1.0);
% Vilebrequin (simplifié)
\draw[thick, popgreen, fill=popgreen!20] (1.5,1.0) circle (0.8);
\draw[thick, popgreen] (1.5,1.0) -- (1.5,0.2);
% Fleches mouvement
\draw[->, thick, poppink] (3.5, 2.8) -- (3.5, 1.5) node[midway, right] {Translation};
\draw[->, thick, poppink] (2.5, 0.2) arc (0:90:1) node[midway, above right] {Rotation};
% Legendes
\node[align=left, anchor=north west] at (-0.5, 5.2) {\textbf{Coupe schématique d'un cylindre}};
\node[anchor=east] at (-0.3, 2.85) {Piston};
\node[anchor=west] at (3.3, 1.0) {Bielle};
\node[anchor=north] at (1.5, -0.5) {Vilebrequin};
\node[anchor=south] at (1.5, 4.5) {Culasse};
\node[anchor=south] at (1.5, 0.5) {Carter (bassin d'huile)};
\end{tikzpicture}
\legende{Principaux organes mobiles d'un moteur à combustion interne (vue en coupe). Le piston effectue un mouvement de translation (va-et-vient) que la bielle et le vilebrequin transforment en rotation.}
\end{popfigure}

\begin{propriete}[Principe de fonctionnement]
Le moteur fonctionne par \cle{cycles} successifs. Chaque cycle comprend \textbf{4 temps} (courses du piston) correspondants à \textbf{2 tours complets du vilebrequin}. L'énergie libérée par la combustion crée une surpression qui pousse le piston : c'est le \cle{temps moteur} (seul temps qui fournit du travail).
\end{propriete}

\courssub{Les éléments d'un moteur et leurs rôles}

\begin{methode}[Identifier les organes d'un moteur sur un schéma]
Pour reconnaître et nommer les éléments d'un moteur à combustion interne sur un schéma :
\begin{enumerate}
\item Repère d'abord le \cle{cylindre} : c'est le tube vertical (le corps du moteur) dans lequel tout se passe.
\item À l'intérieur, la pièce qui coulisse de haut en bas est le \cle{piston}. Il assure l'étanchéité grâce à des \cle{segments}.
\item Sous le piston, la tige qui le relie à l'arbre coudé est la \cle{bielle} ; l'arbre coudé en bas est le \cle{vilebrequin} (il transforme le mouvement de va-et-vient en rotation).
\item En haut du cylindre, dans la \cle{culasse}, cherche les deux \cle{soupapes} (admission et échappement) qui s'ouvrent et se ferment pour laisser entrer le mélange frais et sortir les gaz brûlés.
\item Cherche enfin l'organe d'allumage : la \cle{bougie} (moteur à essence) ou l'\cle{injecteur} (moteur diesel).
\end{enumerate}
\textbf{Réflexe :} apprends la chaîne cinématique par cœur : \cle{piston} $\rightarrow$ \cle{bielle} $\rightarrow$ \cle{vilebrequin} $\rightarrow$ \cle{rotation}.
\end{methode}

\begin{aretenir}[Rôles des organes principaux]
\begin{center}
\begin{tabularx}{\linewidth}{|l|X|}\hline
\rowcolor{popblueL}
\textbf{Organe} & \textbf{Rôle} \\ \hline
Cylindre & Guide le piston ; paroi de la chambre de combustion. \\ \hline
Piston & Reçoit la pression des gaz ; assure l'étanchéité (segments). \\ \hline
Bielle & Relie piston et vilebrequin ; transforme translation en rotation. \\ \hline
Vilebrequin & Sortie de puissance : mouvement de rotation vers les roues. \\ \hline
Soupapes admission & Laissent entrer le mélange air-carburant (essence) ou l'air (diesel). \\ \hline
Soupapes échappement & Laissent sortir les gaz brûlés après combustion. \\ \hline
Bougie (essence) & Déclenche l'allumage par étincelle (allumage commandé). \\ \hline
Injecteur (diesel) & Pulvérise le gasoil dans l'air comprimé et surchauffé. \\ \hline
\end{tabularx}
\end{center}
\end{aretenir}

\begin{popfigure}
\begin{tikzpicture}[scale=0.85, every node/.style={font=\small}]
% Cylindre complet avec soupapes
\draw[thick, popblue] (0,0) -- (0,6) -- (3.5,6) -- (3.5,0) -- cycle;
% Culasse
\draw[thick, popblue, fill=popblue!10] (0,5) rectangle (3.5,6);
% Piston
\draw[thick, poporange, fill=poporange!30] (0.3,2.2) rectangle (3.2,2.8);
% Segments
\draw[thin, popdark] (0.3,2.4) -- (3.2,2.4);
\draw[thin, popdark] (0.3,2.6) -- (3.2,2.6);
% Bielle + Vilebrequin
\draw[thick, popdark] (1.75,2.2) -- (1.75,0.8);
\draw[thick, popgreen, fill=popgreen!20] (1.75,0.8) circle (0.6);
% Soupapes
% Admission (gauche)
\draw[thick, popgreen] (0.8,5) -- (0.8,4.2);
\draw[thick, popgreen, fill=popgreen!30] (0.6,4.2) rectangle (1.0,4.0);
\node[anchor=north, font=\tiny] at (0.8, 5.1) {Admission};
% Echappement (droite)
\draw[thick, poporange] (2.7,5) -- (2.7,4.2);
\draw[thick, poporange, fill=poporange!30] (2.5,4.2) rectangle (2.9,4.0);
\node[anchor=north, font=\tiny] at (2.7, 5.1) {Échappement};
% Bougie (centre haut)
\draw[thick, poppink] (1.75,5.5) -- (1.75,5.1);
\draw[thick, poppink] (1.7,5.1) -- (1.8,5.1); % electrode
\node[anchor=south, font=\tiny] at (1.75, 5.6) {Bougie};
% Fleches
\draw[->, thick, popmuted] (-0.8, 4.5) -- (-0.2, 4.5) node[midway, left] {Mélange frais};
\draw[->, thick, popmuted] (3.7, 4.5) -- (4.3, 4.5) node[midway, right] {Gaz brûlés};
\draw[<->, thick, poppink] (-0.5, 2.2) -- (-0.5, 0.2) node[midway, left] {Course $c$};
\draw[<->, thick, poppink] (0.3, -0.3) -- (3.2, -0.3) node[midway, below] {Alésage $d$};
\end{tikzpicture}
\legende{Schéma détaillé d'un cylindre 4 temps (moteur à essence) : soupapes en culasse, bougie d'allumage, piston avec segments, bielle, vilebrequin. Cotes utiles pour la cylindrée : alésage $d$ et course $c$.}
\end{popfigure}

\begin{exoresolu}[Identifier les organes d'un moteur de moto-taxi]
Le mécanicien de quartier démonte le moteur d'une moto-taxi. Sur le schéma simplifié du cylindre, on voit : une pièce qui coulisse verticalement (1), une tige articulée (2), un arbre coudé (3), deux clapets en haut du cylindre (4) et (5), et une pièce qui produit une étincelle (6).
\begin{enumerate}
\item Nomme les organes (1) à (6).
\item Quel est le rôle de l'organe (3) ?
\item S'agit-il d'un moteur à essence ou d'un moteur diesel ? Justifie.
\end{enumerate}
\tcblower
\textbf{1. Identification des organes.}

On applique la méthode en partant du cylindre :
\begin{itemize}
\item (1) La pièce qui coulisse verticalement dans le cylindre est le \cle{piston}.
\item (2) La tige articulée reliant le piston à l'arbre est la \cle{bielle}.
\item (3) L'arbre coudé situé en bas est le \cle{vilebrequin}.
\item (4) et (5) Les deux clapets situés dans la culasse sont les \cle{soupapes} : soupape d'\cle{admission} et soupape d'\cle{échappement}.
\item (6) La pièce qui produit une étincelle est la \cle{bougie} d'allumage.
\end{itemize}

\textbf{2. Rôle du vilebrequin (3).}

Le vilebrequin transforme le \cle{mouvement de translation} (va-et-vient) du piston en \cle{mouvement de rotation} utilisable pour entraîner les roues de la moto.

\textbf{3. Type de moteur.}

La présence d'une \textbf{bougie} qui produit une étincelle montre que l'allumage du mélange est commandé : c'est un \cle{moteur à essence} (allumage par étincelles). Dans un moteur diesel, il n'y a pas de bougie mais un injecteur, et l'allumage se fait par compression.

\textbf{Conclusion :} (1) piston, (2) bielle, (3) vilebrequin, (4) et (5) soupapes d'admission et d'échappement, (6) bougie ; le vilebrequin transforme la translation en rotation ; c'est un moteur à essence car il possède une bougie d'allumage.
\end{exoresolu}

\courssub{Le cycle à quatre temps}

\begin{methode}[Reconnaître et ordonner les quatre temps]
Le cycle à quatre temps comporte toujours, dans l'ordre immuable :
\begin{enumerate}
\item \cle{Admission} : la soupape d'admission est ouverte, le piston descend, le mélange air-essence (ou l'air seul en diesel) entre dans le cylindre.
\item \cle{Compression} : les deux soupapes sont fermées, le piston remonte et comprime le gaz (augmentation de la pression et de la température).
\item \cle{Combustion -- détente} (temps moteur) : la bougie produit l'étincelle (essence) ou le gasoil s'enflamme par compression (diesel), les gaz brûlés poussent violemment le piston vers le bas : c'est le seul temps qui produit du travail mécanique.
\item \cle{Échappement} : la soupape d'échappement s'ouvre, le piston remonte et chasse les gaz brûlés vers le pot d'échappement.
\end{enumerate}
\textbf{Réflexes :}
\begin{itemize}
\item Un cycle complet = 4 courses du piston = \textbf{2 tours de vilebrequin}.
\item Pour identifier un temps sur un schéma, regarde d'abord les \textbf{soupapes} (ouverte/fermée), puis le \textbf{sens de déplacement} du piston.
\item Moyen mnémotechnique : « \textbf{A}dmission, \textbf{C}ompression, \textbf{D}étente, \textbf{É}chappement ».
\end{itemize}
\end{methode}

\begin{propriete}[Bilan du cycle 4 temps]
\begin{center}
\begin{tabularx}{\linewidth}{|c|X|X|X|}\hline
\rowcolor{popblueL}
\textbf{Temps} & \textbf{Soupapes} & \textbf{Mouvement piston} & \textbf{Échange énergétique} \\ \hline
Admission & Admission ouverte & Descend (PMB $\rightarrow$ PMB) & Travail reçu (aspiration) \\ \hline
Compression & Fermées & Remonte (PMB $\rightarrow$ PMH) & Travail reçu (compression) \\ \hline
Combustion-Détente & Fermées & Descend (PMH $\rightarrow$ PMB) & \textbf{Travail fourni (moteur)} \\ \hline
Échappement & Échappement ouverte & Remonte (PMB $\rightarrow$ PMH) & Travail reçu (expulsion) \\ \hline
\end{tabularx}
\end{center}
PMH = Point Mort Haut (piston tout en haut) ; PMB = Point Mort Bas (piston tout en bas).
\end{propriete}

\begin{popfigure}
\begin{tikzpicture}[scale=0.7, every node/.style={font=\small}]
% 4 petits cylindres côte à côte
\def\cyl#1#2#3#4#5#6{ % x, y, label, piston_y, valve_adm, valve_ech
\begin{scope}[shift={(#1,#2)}]
\draw[thick] (0,0) rectangle (2,4);
\draw[thick, fill=poporange!30] (0.2,#4) rectangle (1.8,#4+0.6);
% Soupapes
\ifnum#5=1 \draw[thick, popgreen, fill=popgreen!30] (0.3,3.8) rectangle (0.7,3.4); \else \draw[thick] (0.3,3.8) rectangle (0.7,3.4); \fi
\ifnum#6=1 \draw[thick, poporange, fill=poporange!30] (1.3,3.8) rectangle (1.7,3.4); \else \draw[thick] (1.3,3.8) rectangle (1.7,3.4); \fi
% Bougie
\draw[thick, poppink] (1,4.2) -- (1,3.9);
% Fleche piston
\ifdim#4pt=2.5pt \draw[->, thick, poppink] (2.3,2.8) -- (2.3,1.5); \else \draw[->, thick, poppink] (2.3,1.5) -- (2.3,2.8); \fi
\node[anchor=north] at (1,-0.3) {\textbf{#3}};
\end{scope}
}
% Admission : piston bas (2.5), adm ouverte(1), ech fermee(0)
\cyl{0}{0}{Admission}{2.5}{1}{0}
% Compression : piston haut (0.9), fermées
\cyl{3}{0}{Compression}{0.9}{0}{0}
% Detente : piston haut (0.9), fermées, flamme
\cyl{6}{0}{Combustion-Détente}{0.9}{0}{0}
\node[poppink, font=\tiny] at (7, 2.5) {\textbf{BOUM!}};
% Echappement : piston bas (2.5), adm fermee, ech ouverte
\cyl{9}{0}{Échappement}{2.5}{0}{1}
\end{tikzpicture}
\legende{Les 4 temps du cycle : 1. Admission (piston descend, soupape adm ouverte), 2. Compression (piston monte, soupapes fermées), 3. Combustion-Détente (explosion, piston poussé vers le bas), 4. Échappement (piston monte, soupape éch ouverte).}
\end{popfigure}

\begin{exoresolu}[Ordonner les temps d'un moteur de groupe électrogène]
Le moteur d'un groupe électrogène fonctionne selon un cycle à quatre temps. Voici quatre descriptions dans le désordre :
\begin{itemize}
\item \textbf{a)} Les deux soupapes sont fermées, le piston remonte et comprime le mélange.
\item \textbf{b)} La soupape d'échappement est ouverte, le piston remonte et expulse les gaz brûlés.
\item \textbf{c)} L'étincelle de la bougie enflamme le mélange comprimé ; les gaz poussent le piston vers le bas.
\item \textbf{d)} La soupape d'admission est ouverte, le piston descend et aspire le mélange air-essence.
\end{itemize}
\begin{enumerate}
\item Remets ces temps dans l'ordre du cycle et nomme-les.
\item Pendant combien de tours de vilebrequin se déroule un cycle complet ?
\item Quel est le temps moteur ? Justifie.
\end{enumerate}
\tcblower
\textbf{1. Ordre du cycle.}

Le cycle commence toujours par l'entrée des gaz frais :
\begin{itemize}
\item \textbf{d)} = \cle{Admission} (soupape d'admission ouverte, piston descendant).
\item \textbf{a)} = \cle{Compression} (soupapes fermées, piston montant).
\item \textbf{c)} = \cle{Combustion -- détente} (étincelle, piston poussé vers le bas).
\item \textbf{b)} = \cle{Échappement} (soupape d'échappement ouverte, piston montant).
\end{itemize}
L'ordre correct est donc : \textbf{d -- a -- c -- b}.

\textbf{2. Nombre de tours de vilebrequin.}

Un cycle comporte 4 courses du piston. Or une course (montée ou descente) correspond à un demi-tour de vilebrequin. Donc :
\[ 4 \text{ courses} \times \frac{1}{2} \text{ tour} = 2 \text{ tours de vilebrequin}. \]

\textbf{3. Temps moteur.}

Le temps moteur est le \textbf{3\ieme{} temps (c) : combustion -- détente}. C'est le seul temps pendant lequel les gaz brûlés, en se détendant, \textbf{poussent} le piston et fournissent du travail mécanique. Les trois autres temps consomment de l'énergie (fournie par l'inertie du volant moteur et du vilebrequin).

\textbf{Conclusion :} l'ordre est d--a--c--b (admission, compression, combustion--détente, échappement) ; un cycle dure 2 tours de vilebrequin ; le temps moteur est la combustion--détente car c'est le seul temps qui produit du travail.
\end{exoresolu}

\courssub{Moteur à essence et moteur diesel : différences}

\begin{definition}[Allumage commandé vs Allumage par compression]
\begin{itemize}
\item \cle{Moteur à essence (cycle Otto)} : \textbf{Allumage commandé}. Le mélange air+essence est comprimé, puis une \textbf{bougie} déclenche la combustion par une étincelle au moment précis voulu.
\item \cle{Moteur diesel (cycle Diesel)} : \textbf{Allumage par compression (auto-inflammation)}. Seul de l'\textbf{air} est comprimé très fortement (rapport volumétrique $\varepsilon \approx 18$), ce qui élève sa température au-dessus de la température d'auto-inflammation du gasoil ($\approx 250\ {}^\circ\mathrm{C}$). L'\textbf{injecteur} pulvérise alors le gasoil qui s'enflamme spontanément.
\end{itemize}
\end{definition}

\begin{methode}[Comparer un moteur à essence et un moteur diesel]
Pour comparer les deux types de moteurs, raisonne toujours sur les quatre mêmes critères :
\begin{enumerate}
\item \textbf{Carburant} : essence (moteur à essence) ; gasoil (moteur diesel).
\item \textbf{Ce qui entre au cylindre à l'admission} : mélange air + essence ; air seul (diesel).
\item \textbf{Mode d'allumage} : par étincelle de la \textbf{bougie} (allumage commandé) ; par \textbf{compression} de l'air fortement chauffé, le gasoil étant injecté au moment voulu (auto-inflammation).
\item \textbf{Organe caractéristique} : bougie ; injecteur.
\end{enumerate}
\textbf{Réflexes :} « diesel » $\Rightarrow$ pas de bougie, compression plus forte, moteur plus lourd et plus robuste (camions, groupes électrogènes) ; « essence » $\Rightarrow$ bougie, moteur plus léger (motos, voitures particulières).
\end{methode}

\begin{center}
\begin{tabularx}{\linewidth}{|l|X|X|}\hline
\rowcolor{popblueL}
\textbf{Critère} & \textbf{Moteur à essence} & \textbf{Moteur diesel} \\ \hline
Carburant & Essence (plus volatile) & Gasoil (moins volatile, plus énergétique) \\ \hline
Admission & Mélange air + essence & Air seul \\ \hline
Allumage & \textbf{Commandé} par \textbf{bougie} (étincelle) & \textbf{Par compression} (auto-inflammation) via \textbf{injecteur} \\ \hline
Rapport volumétrique $\varepsilon$ & Faible ($\approx 8$ à $12$) & Élevé ($\approx 14$ à $22$) \\ \hline
Rendement & Moins bon ($\approx 25-30\%$) & Meilleur ($\approx 35-45\%$) \\ \hline
Couple / Robustesse & Moins coupleux, plus léger & Très coupleux, très robuste, plus lourd \\ \hline
Usage courant & Motos, voitures légères & Camions, bus, groupes électrogènes, bateaux \\ \hline
\end{tabularx}
\end{center}

\begin{savaistu}[La SONARA et les carburants]
La \textbf{SONARA} (Société Nationale de Raffinage) de Limbé est la seule raffinerie du Cameroun. Elle sépare le pétrole brut par distillation fractionnée : l'\textbf{essence} (fractions légères, point d'ébullition bas) sert aux moteurs à allumage commandé ; le \textbf{gasoil} (fractions plus lourdes) sert aux moteurs diesel. La qualité du carburant (indice d'octane pour l'essence, indice de cétane pour le gasoil) conditionne le bon fonctionnement et la longévité du moteur.
\end{savaistu}

\begin{exoresolu}[Essence ou diesel ?]
À la SONARA de Limbé, on produit de l'essence et du gasoil. Deux véhicules arrivent à la station : une voiture particulière dont le moteur possède des bougies d'allumage, et un camion dont le moteur, dépourvu de bougies, possède des injecteurs et comprime l'air à un rapport volumétrique élevé (environ 18).
\begin{enumerate}
\item Pour chaque véhicule, indique le type de moteur et le carburant à utiliser.
\item Explique pourquoi le moteur du camion n'a pas besoin de bougie.
\item Cite un avantage du moteur diesel pour le transport de marchandises.
\end{enumerate}
\tcblower
\textbf{1. Type de moteur et carburant.}
\begin{itemize}
\item La voiture possède des \textbf{bougies} : c'est un \cle{moteur à essence}, à allumage commandé par étincelles. Carburant : \textbf{essence}.
\item Le camion n'a pas de bougie mais des \textbf{injecteurs} et une compression élevée : c'est un \cle{moteur diesel}. Carburant : \textbf{gasoil}.
\end{itemize}

\textbf{2. Pourquoi le moteur diesel n'a pas besoin de bougie.}

Dans le moteur diesel, seul de l'\textbf{air} est admis puis comprimé très fortement (rapport volumétrique élevé, environ 18). Cette compression élève fortement la température de l'air (au-delà de la température d'auto-inflammation du gasoil, vers $600-800\ {}^\circ\mathrm{C}$). Lorsque l'injecteur pulvérise le gasoil dans cet air très chaud, le carburant \textbf{s'enflamme spontanément} : c'est l'\cle{auto-inflammation} (allumage par compression). Aucune étincelle n'est donc nécessaire.

\textbf{3. Avantage du moteur diesel.}

Le moteur diesel est plus \textbf{robuste} et plus \textbf{coupleux} (il fournit un grand effort à bas régime), et il consomme en général moins de carburant (meilleur rendement) : il est donc bien adapté aux camions qui transportent de lourdes charges sur de longues distances.

\textbf{Conclusion :} la voiture a un moteur à essence (carburant : essence) et le camion un moteur diesel (carburant : gasoil) ; le diesel s'allume par auto-inflammation grâce à la forte compression de l'air, d'où l'absence de bougie ; le diesel est avantageux pour le transport lourd grâce à sa robustesse et sa sobriété.
\end{exoresolu}

\courssub{Sécurité, entretien, environnement et calculs d'usage}

\begin{methode}[Calculs sur le cylindre : cylindrée et rapport volumétrique]
\begin{enumerate}
\item La \cle{cylindrée unitaire} $V_u$ est le volume balayé par le piston entre le point mort haut (PMH) et le point mort bas (PMB). C'est le volume d'un cylindre de révolution :
\[ V_u = \pi \times r^2 \times c \]
où $r$ est le rayon de l'alésage ($r = \dfrac{d}{2}$, $d$ = alésage) et $c$ la course du piston.
\item La \cle{cylindrée totale} : $V_{\text{tot}} = V_u \times n$, où $n$ est le nombre de cylindres.
\item Le \cle{rapport volumétrique} : $\varepsilon = \dfrac{V_u + V_c}{V_c}$, où $V_c$ est le volume de la chambre de combustion (volume restant au PMH).
\end{enumerate}
\textbf{Réflexes :} convertis toutes les longueurs dans la même unité (cm pour obtenir des cm$^3$) ; $1\,000$ cm$^3$ = 1 L ; le rayon est la \textbf{moitié} de l'alésage.
\end{methode}

\begin{exoresolu}[Cylindrée du moteur d'une moto-taxi]
Le moteur monocylindre d'une moto-taxi a un alésage $d = \SI{5,0}{cm}$ et une course $c = \SI{5,0}{cm}$. Le volume de la chambre de combustion est $V_c = \SI{25}{cm^3}$. On prendra $\pi \approx 3,14$.
\begin{enumerate}
\item Calcule la cylindrée unitaire $V_u$ de ce moteur.
\item Calcule le rapport volumétrique $\varepsilon$.
\item Exprime la cylindrée en litres.
\end{enumerate}
\tcblower
\textbf{1. Cylindrée unitaire.}

Le rayon de l'alésage est :
\[ r = \frac{d}{2} = \frac{5,0}{2} = \SI{2,5}{cm}. \]
Le volume balayé par le piston est celui d'un cylindre de rayon $r$ et de hauteur égale à la course $c$ :
\begin{align*}
V_u &= \pi \times r^2 \times c \\
&= 3,14 \times (2,5)^2 \times 5,0 \\
&= 3,14 \times 6,25 \times 5,0 \\
&= \SI{98,125}{cm^3} \approx \SI{98}{cm^3}.
\end{align*}

\textbf{2. Rapport volumétrique.}
\[ \varepsilon = \frac{V_u + V_c}{V_c} = \frac{98,125 + 25}{25} = \frac{123,125}{25} \approx 4,9. \]
Le mélange est donc comprimé environ 5 fois (typique d'un petit moteur 2 temps ou 4 temps ancien ; les moteurs modernes ont $\varepsilon \approx 10-12$).

\textbf{3. Conversion en litres.}
\[ V_u = 98,125 \text{ cm}^3 = \frac{98,125}{1000} \text{ L} \approx \SI{0,098}{L}. \]

\textbf{Conclusion :} la cylindrée unitaire vaut environ \SI{98}{cm^3} (soit \SI{0,098}{L}) et le rapport volumétrique vaut environ 4,9 : le mélange est comprimé près de 5 fois avant l'allumage.
\end{exoresolu}

\begin{methode}[Calculs de consommation et de coût]
\begin{enumerate}
\item La consommation moyenne $C$ s'exprime en litres pour 100 km : $C = \dfrac{V_{\text{carburant}}}{d} \times 100$.
\item Le volume de carburant nécessaire pour un trajet : $V = \dfrac{C \times d}{100}$.
\item Le coût du trajet : $\text{Coût} = V \times p$, où $p$ est le prix du litre de carburant.
\end{enumerate}
\textbf{Réflexes :} vérifie toujours la cohérence des unités (km, L, FCFA) ; arrondis le résultat final de façon raisonnable ; pense à la sécurité : le carburant est inflammable, on ne le manipule jamais près d'une flamme.
\end{methode}

\begin{exoresolu}[Consommation du groupe électrogène et de la moto]
Une moto-taxi consomme en moyenne 3,0 L d'essence pour 100 km. Le litre d'essence coûte 730 FCFA à la station.
\begin{enumerate}
\item Quel volume d'essence faut-il pour parcourir 250 km ?
\item Quel est le coût de ce trajet ?
\item Un groupe électrogène consomme 1,2 L d'essence par heure. Pendant une coupure ENEO de 5 h, quel volume consomme-t-il et quel est le coût correspondant ?
\end{enumerate}
\tcblower
\textbf{1. Volume pour 250 km.}
\[ V = \frac{C \times d}{100} = \frac{3,0 \times 250}{100} = \frac{750}{100} = \SI{7,5}{L}. \]

\textbf{2. Coût du trajet.}
\[ \text{Coût} = V \times p = 7,5 \times 730 = 5\,475 \text{ FCFA}. \]

\textbf{3. Groupe électrogène.}

Volume consommé en 5 h :
\[ V' = 1,2 \times 5 = \SI{6,0}{L}. \]
Coût correspondant :
\[ \text{Coût}' = 6,0 \times 730 = 4\,380 \text{ FCFA}. \]

\textbf{Conclusion :} la moto consomme 7,5 L pour 250 km, soit un coût de 5\,475 FCFA ; le groupe électrogène consomme 6,0 L en 5 h, soit 4\,380 FCFA. Rappel de sécurité : l'essence est stockée loin de toute flamme et le groupe électrogène fonctionne à l'extérieur, car les gaz d'échappement contiennent le monoxyde de carbone, gaz toxique.
\end{exoresolu}

\begin{methode}[Sécurité, entretien et environnement autour d'un moteur]
Pour répondre à une question de sécurité ou d'entretien, mobilise ces réflexes :
\begin{enumerate}
\item \textbf{Carburant} : essence et gasoil sont inflammables ; pas de flamme ni d'étincelle à proximité ; stockage dans des bidons fermés, à l'ombre, hors de portée des enfants.
\item \textbf{Gaz d'échappement} : ils contiennent le \cle{monoxyde de carbone} \ce{CO}, gaz incolore, inodore et mortel : jamais de moteur en marche dans un local fermé (groupe électrogène dehors !).
\item \textbf{Entretien courant} : vidange régulière de l'huile moteur (lubrification, refroidissement), nettoyage/remplacement du filtre à air, vérification de l'état de la bougie (calage, propreté) et du niveau de carburant.
\item \textbf{Environnement} : l'huile usagée ne se jette ni au sol ni dans les caniveaux (pollution eau/sol) ; la déposer en centre de collecte. Les gaz d'échappement polluent l'air (oxydes d'azote \ce{NOx}, \ce{CO2}, particules).
\end{enumerate}
\end{methode}

\begin{exoresolu}[Le groupe électrogène dans la boutique]
Pendant une coupure de courant, un commerçant de Douala installe son groupe électrogène à l'intérieur de sa boutique fermée pour éviter qu'on ne le vole. Au bout d'une heure, il se sent mal : maux de tête, nausées, vertiges.
\begin{enumerate}
\item Quelle est la cause probable de son malaise ? Nomme le gaz responsable et donne sa formule.
\item Pourquoi ce gaz est-il particulièrement dangereux ?
\item Quelles consignes le commerçant aurait-il dû respecter ?
\end{enumerate}
\tcblower
\textbf{1. Cause du malaise.}

Le moteur du groupe électrogène rejette des gaz d'échappement qui contiennent du \cle{monoxyde de carbone}, de formule \ce{CO}. Dans un local fermé, ce gaz s'accumule et le commerçant l'a inhalé : il souffre d'une intoxication au monoxyde de carbone (maux de tête, nausées, vertiges sont les symptômes typiques).

\textbf{2. Dangerosité du \ce{CO}.}

Le monoxyde de carbone est \textbf{incolore} et \textbf{inodore} : on ne peut ni le voir ni le sentir, donc on s'intoxique sans s'en rendre compte. De plus, il se fixe sur l'hémoglobine du sang à la place du dioxygène (affinité 200 à 250 fois plus forte) et empêche le transport de l'oxygène vers les organes : il peut provoquer la mort en quelques heures.

\textbf{3. Consignes à respecter.}
\begin{itemize}
\item Faire fonctionner le groupe électrogène \textbf{à l'extérieur}, dans un endroit bien aéré, loin des portes et fenêtres.
\item Ne jamais le mettre en marche dans une pièce fermée (boutique, garage, chambre).
\item Conserver le carburant dans un bidon fermé, loin de toute flamme.
\item En cas de malaise : éteindre le moteur, sortir la victime au grand air et appeler les secours (112 ou 119).
\end{itemize}

\textbf{Conclusion :} le malaise est dû à une intoxication au monoxyde de carbone \ce{CO}, gaz incolore et inodore qui bloque le transport de l'oxygène dans le sang ; le groupe électrogène doit toujours fonctionner à l'extérieur, dans un lieu ventilé.
\end{exoresolu}

\begin{attention}[Les 5 erreurs qui coûtent des points au BEPC]
\begin{enumerate}
\item \textbf{Inverser l'ordre des quatre temps.} L'ordre est immuable : admission $\rightarrow$ compression $\rightarrow$ combustion--détente $\rightarrow$ échappement. Écrire « échappement puis compression » est une erreur éliminatoire. Retiens : A -- C -- D -- É.
\item \textbf{Confondre bougie et injecteur.} La \textbf{bougie} caractérise le moteur à \textbf{essence} (allumage par étincelle) ; l'\textbf{injecteur} caractérise le moteur \textbf{diesel} (allumage par compression). Dire « la bougie du moteur diesel » est faux.
\item \textbf{Oublier que le rayon est la moitié de l'alésage.} Dans le calcul de la cylindrée $V_u = \pi r^2 c$, beaucoup d'élèves utilisent l'alésage $d$ à la place de $r = d/2$, ce qui donne un résultat 4 fois trop grand.
\item \textbf{Oublier les unités et les conversions.} La cylindrée se donne en cm$^3$ (ou en L : 1 L = 1\,000 cm$^3$) ; les longueurs doivent être converties dans la même unité avant le calcul. Un résultat sans unité perd des points.
\item \textbf{Négliger la consigne de sécurité dans la rédaction.} Quand l'énoncé mentionne un moteur en local fermé ou une manipulation d'essence, le correcteur attend la mention du \textbf{monoxyde de carbone \ce{CO}} (gaz toxique, incolore et inodore) ou du risque d'\textbf{incendie}. Une réponse technique sans la consigne de sécurité est incomplète.
\end{enumerate}
\end{attention}

\begin{exercice}[Entraînement BEPC - Moteur 4 temps]
\begin{enumerate}
\item Complète le tableau ci-dessous pour un moteur 4 temps essence.
\begin{center}
\begin{tabularx}{\linewidth}{|c|X|X|X|}\hline
\rowcolor{popblueL}
\textbf{Temps} & \textbf{Soupape admission} & \textbf{Soupape échappement} & \textbf{Sens piston} \\ \hline
Admission & & & \\ \hline
Compression & & & \\ \hline
Combustion-Détente & & & \\ \hline
Échappement & & & \\ \hline
\end{tabularx}
\end{center}
\item Un moteur diesel 4 cylindres a une cylindrée unitaire de \SI{500}{cm^3} et un rapport volumétrique de 18. Calcule sa cylindrée totale en litres et le volume de sa chambre de combustion $V_c$ (en cm$^3$).
\item \textbf{Sécurité :} Pourquoi est-il interdit de faire le plein d'essence moteur tournant ? Cite deux risques.
\end{enumerate}
\end{exercice}

\begin{corrige}[Exercice Entraînement BEPC]
\begin{enumerate}
\item \textbf{Tableau complété :}
\begin{center}
\begin{tabularx}{\linewidth}{|c|X|X|X|}\hline
\rowcolor{popblueL}
\textbf{Temps} & \textbf{Soupape admission} & \textbf{Soupape échappement} & \textbf{Sens piston} \\ \hline
Admission & Ouverte & Fermée & Descend (PMB $\rightarrow$ PMH) \\ \hline
Compression & Fermée & Fermée & Remonte (PMB $\rightarrow$ PMH) \\ \hline
Combustion-Détente & Fermée & Fermée & Descend (PMH $\rightarrow$ PMB) \\ \hline
Échappement & Fermée & Ouverte & Remonte (PMB $\rightarrow$ PMH) \\ \hline
\end{tabularx}
\end{center}
\item \textbf{Cylindrée totale :} $V_{\text{tot}} = 4 \times 500 = 2000 \text{ cm}^3 = \SI{2,0}{L}$.
\textbf{Volume chambre combustion :} $\varepsilon = \frac{V_u + V_c}{V_c} \Rightarrow 18 = \frac{500 + V_c}{V_c} \Rightarrow 18 V_c = 500 + V_c \Rightarrow 17 V_c = 500 \Rightarrow V_c \approx \SI{29,4}{cm^3}$.
\item \textbf{Risques :} 1) Risque d'\textbf{incendie} ou d'\textbf{explosion} : les vapeurs d'essence (invisibles, plus lourdes que l'air) peuvent être enflammées par une étincelle du système d'allumage, le pot d'échappement chaud ou l'électricité statique. 2) Risque d'\textbf{asphyxie} ou d'intoxication par les vapeurs d'essence (toxiques) ou le \ce{CO} si le moteur tourne dans un espace confiné.
\end{enumerate}
\end{corrige}

\newpage

\coursec{Exercices}

\courssub{Je vérifie mes connaissances}

\begin{exercice}[QCM : essence ou diesel ?]
Pour chaque affirmation, choisis la bonne réponse (une seule réponse correcte par question).
\begin{enumerate}
\item L'allumage du mélange air-essence est provoqué par :
\begin{enumerate}
\item une bougie d'allumage \item la compression seule \item un injecteur \item le carburateur
\end{enumerate}
\item Dans un moteur diesel, le carburant utilisé est :
\begin{enumerate}
\item l'essence \item le gasoil \item l'huile moteur \item le kérosène
\end{enumerate}
\item Dans un moteur diesel, l'inflammation du mélange se produit :
\begin{enumerate}
\item grâce à une étincelle électrique \item par compression de l'air seul \item grâce à une flamme extérieure \item par le volant moteur
\end{enumerate}
\item Le véhicule le plus souvent équipé d'un moteur diesel est :
\begin{enumerate}
\item une moto-taxi \item un camion de transport \item un vélo \item une tondeuse
\end{enumerate}
\end{enumerate}
\end{exercice}

\begin{exercice}[Vrai ou faux ?]
Réponds par VRAI ou FAUX, puis justifie ta réponse en une phrase.
\begin{enumerate}
\item Le cycle d'un moteur à quatre temps comporte : admission, compression, détente, échappement.
\item Le piston transforme le mouvement de rotation du vilebrequin en mouvement de translation.
\item Un moteur à essence possède des bougies d'allumage.
\item Le monoxyde de carbone est un gaz toxique dégagé par les moteurs à combustion interne.
\end{enumerate}
\end{exercice}

\begin{exercice}[Définitions]
Définis en une ou deux phrases chacun des termes suivants :
\begin{enumerate}
\item moteur à combustion interne ;
\item cylindrée ;
\item temps moteur ;
\item soupape.
\end{enumerate}
\end{exercice}

\begin{exercice}[Calcul direct]
Le moteur d'une moto-taxi tourne à \SI{2400}{tr/min}.
\begin{enumerate}
\item Combien de tours le vilebrequin effectue-t-il en une minute ?
\item Sachant qu'un cycle complet à quatre temps correspond à 2 tours de vilebrequin, combien de cycles complets chaque cylindre réalise-t-il en une minute ?
\item Combien de temps moteur (temps de détente) se produisent par minute dans ce cylindre ?
\end{enumerate}
\end{exercice}

\courssub{Je m'entraîne}

\begin{exercice}[Remettre le cycle dans l'ordre]
Voici les descriptions mélangées des quatre temps d'un moteur à essence :
\begin{itemize}
\item \textbf{A} : le mélange brûlé est chassé hors du cylindre par le piston qui remonte, la soupape d'échappement étant ouverte ;
\item \textbf{B} : le mélange air-essence est aspiré dans le cylindre pendant que le piston descend, la soupape d'admission étant ouverte ;
\item \textbf{C} : l'étincelle de la bougie enflamme le mélange comprimé ; les gaz chauds poussent violemment le piston vers le bas ;
\item \textbf{D} : les deux soupapes étant fermées, le piston remonte et comprime le mélange air-essence.
\end{itemize}
\begin{enumerate}
\item Remets ces quatre temps dans l'ordre chronologique du cycle.
\item Nomme chaque temps.
\item Indique lequel est le temps moteur et justifie ta réponse.
\end{enumerate}
\end{exercice}

\begin{exercice}[Légender un schéma de moteur]
On donne le schéma simplifié d'un cylindre de moteur à combustion interne portant les repères 1 à 6 : piston, cylindre, bougie, soupape d'admission, soupape d'échappement, bielle.
\begin{enumerate}
\item Reproduis le schéma et place chaque repère au bon endroit.
\item Donne le rôle de quatre organes au choix parmi les six.
\item Indique sur ton schéma le sens de déplacement du piston pendant le temps de compression.
\end{enumerate}
\end{exercice}

\begin{exercice}[Chaîne énergétique d'un moteur]
\begin{enumerate}
\item Recopie et complète la chaîne de transformation d'énergie d'un moteur à combustion interne :
\[
\text{énergie chimique} \longrightarrow \ldots \longrightarrow \ldots
\]
\item Explique, à partir de cette chaîne, pourquoi un moteur chauffe pendant son fonctionnement.
\item Quel organe permet d'évacuer une partie de cette chaleur ? Cite un exemple de système de refroidissement.
\end{enumerate}
\end{exercice}

\begin{exercice}[Moteur à quatre cylindres]
Le moteur d'une voiture possède 4 cylindres et tourne à \SI{3000}{tr/min}.
\begin{enumerate}
\item Calcule le nombre de cycles complets réalisés par chaque cylindre en une minute.
\item Calcule le nombre total de temps moteur (détentes) produits par l'ensemble des 4 cylindres en une minute.
\item Explique pourquoi un moteur à 4 cylindres tourne de façon plus régulière qu'un moteur à 1 seul cylindre.
\end{enumerate}
\end{exercice}

\courssub{Je me prépare au Bac}

\begin{exercice}[Le groupe électrogène du quartier]
Pendant les délestages, M. Etoundi utilise un groupe électrogène à essence pour alimenter sa boutique. Une nuit, il place le groupe à l'intérieur de la boutique, portes et fenêtres fermées, « pour éviter qu'on ne le vole ».
\begin{enumerate}
\item Définis un moteur à combustion interne et cite les deux types étudiés en classe.
\item Le groupe électrogène de M. Etoundi possède une bougie d'allumage. S'agit-il d'un moteur essence ou diesel ? Justifie.
\item Nomme les quatre temps du cycle de ce moteur dans l'ordre.
\item Quel gaz toxique, invisible et inodore, s'accumule dans la boutique ? Pourquoi est-il dangereux ?
\item Propose deux consignes de sécurité que M. Etoundi aurait dû respecter.
\item Pourquoi ne faut-il jamais faire le plein d'essence pendant que le moteur tourne ?
\end{enumerate}
\end{exercice}

\begin{exercice}[La moto-taxi de Ngaoundéré]
La moto-taxi de Ibrahim est équipée d'un moteur à essence monocylindre (1 seul cylindre) de cylindrée \SI{125}{cm^3}, tournant à \SI{4000}{tr/min} en vitesse de croisière.
\begin{enumerate}
\item Explique ce que signifie « cylindrée de \SI{125}{cm^3} ».
\item Calcule le nombre de tours de vilebrequin effectués en une seconde.
\item Sachant qu'un cycle complet correspond à 2 tours de vilebrequin, calcule le nombre de temps moteur par seconde.
\item Pendant le temps d'admission, quelle soupape est ouverte ? Quel mélange entre dans le cylindre ?
\item Ibrahim remarque que son moteur chauffe beaucoup. Donne l'origine de cette chaleur et cite une conséquence possible d'une surchauffe.
\item Complète la chaîne énergétique : énergie chimique $\rightarrow$ \ldots $\rightarrow$ \ldots
\end{enumerate}
\end{exercice}

\begin{exercice}[Essence ou gasoil à la SONARA ?]
Un transporteur de la ville de Douala hésite entre deux camions : le camion A est équipé d'un moteur à essence, le camion B d'un moteur diesel. Les deux carburants sont produits par raffinage du pétrole.
\begin{enumerate}
\item Recopie et complète le tableau comparatif suivant :
\begin{center}
\begin{tabularx}{\linewidth}{|l|X|X|}
\hline
\rowcolor{popblueL} \textbf{Critère} & \textbf{Moteur à essence} & \textbf{Moteur diesel} \\
\hline
Carburant utilisé & & \\
\hline
Organe d'allumage & & \\
\hline
Mode d'inflammation du mélange & & \\
\hline
\end{tabularx}
\end{center}
\item Dans le moteur diesel, quel fluide est comprimé seul avant l'injection du carburant ?
\item Pourquoi le taux de compression d'un moteur diesel est-il plus élevé que celui d'un moteur à essence ?
\item Le transporteur parcourt de longues distances avec de lourdes charges. Quel camion lui conseilles-tu ? Justifie ton choix par deux arguments.
\item Cite deux mesures d'entretien régulier qui prolongent la durée de vie d'un moteur.
\item Les gaz d'échappement des camions contribuent à la pollution de l'air à Douala. Cite un gaz polluant rejeté et propose une mesure pour réduire cette pollution.
\end{enumerate}
\end{exercice}

\newpage

\coursec{Corrigés des exercices}

\coursec{Corrigés des exercices — Fonctionnement d'un moteur à combustion interne}

\begin{corrige}[Exercice 1 — QCM : essence ou diesel ?]
\begin{enumerate}
\item \textbf{Réponse 1 : une bougie d'allumage.}
      \par Dans un moteur à essence, le mélange air-essence est enflammé par une étincelle électrique produite par la bougie d'allumage (allumage commandé).
\item \textbf{Réponse 2 : le gasoil.}
      \par Le moteur diesel fonctionne au gasoil (aussi appelé gazole), un carburant moins volatil que l'essence, obtenu par raffinage du pétrole.
\item \textbf{Réponse 3 : par compression de l'air seul.}
      \par Dans un moteur diesel, seul l'air est comprimé fortement, ce qui élève sa température au-dessus de la température d'auto-inflammation du gasoil (environ 400--500 $^\circ\mathrm{C}$). Le carburant est injecté à la fin de la compression et s'enflamme spontanément au contact de l'air brûlant (allumage par compression).
\item \textbf{Réponse 4 : un camion de transport.}
      \par Les camions, bus et groupes électrogènes de forte puissance sont quasi exclusivement équipés de moteurs diesel pour leur couple élevé à bas régime et leur meilleur rendement énergétique. Les motos-taxis et tondeuses utilisent généralement des moteurs à essence (plus légers, moins chers pour de faibles puissances).
\end{enumerate}
\end{corrige}

\begin{corrige}[Exercice 2 — Vrai ou faux ?]
\begin{enumerate}
\item \textbf{VRAI.} Le cycle théorique à quatre temps se déroule dans l'ordre : admission (remplissage), compression (compression du mélange), détente ou explosion (combustion et travail), échappement (vidange des gaz brûlés).
\item \textbf{FAUX.} C'est le mécanisme bielle-vilebrequin qui transforme le mouvement alternatif (de translation) du piston en mouvement de rotation du vilebrequin. Le vilebrequin, par sa rotation, ramène le piston vers le haut (translation) lors des temps résistants.
\item \textbf{VRAI.} Le moteur à essence est un moteur à allumage commandé : une bougie d'allumage par cylindre délivre l'étincelle nécessaire à l'inflammation du mélange air-essence comprimé.
\item \textbf{VRAI.} La combustion incomplète de l'essence ou du gasoil (manque d'oxygène) produit du monoxyde de carbone (\ce{CO}), gaz incolore, inodore, très toxique (se fixe sur l'hémoglobine à la place de l'\ce{O2}) et mortel en espace confiné.
\end{enumerate}
\end{corrige}

\begin{corrige}[Exercice 3 — Définitions]
\begin{enumerate}
\item \textbf{Moteur à combustion interne :} Moteur thermique dans lequel la combustion du carburant (essence ou gasoil) avec le comburant (air) a lieu directement à l'intérieur de la chambre de combustion (le cylindre), produisant une forte pression qui pousse le piston.
\item \textbf{Cylindrée :} Volume balayé par le piston lors de son déplacement entre le point mort bas (PMB) et le point mort haut (PMH). Elle correspond au volume utile d'un cylindre. Pour un moteur multicylindres, la cylindrée totale est la somme des cylindrées unitaires. Unité : \si{cm^3} ou \si{L} (1 L = 1000 \si{cm^3}).
\item \textbf{Temps moteur (ou temps de détente) :} Le troisième temps du cycle à quatre temps. C'est le seul temps où le moteur fournit de l'énergie mécanique : les gaz brûlés à haute pression, issus de la combustion, se détendent et poussent violemment le piston du PMH vers le PMB, entraînant la rotation du vilebrequin.
\item \textbf{Soupape :} Organe mobile (clapet) commandé par l'arbre à cames, qui assure l'étanchéité de la chambre de combustion et en commande l'ouverture/fermeture. La soupape d'admission laisse entrer le mélange (ou l'air), la soupape d'échappement laisse sortir les gaz brûlés.
\end{enumerate}
\end{corrige}

\begin{corrige}[Exercice 4 — Calcul direct]
Données : régime moteur $N = \SI{2400}{tr/min}$.
\begin{enumerate}
\item Le vilebrequin effectue \textbf{2400 tours} en une minute (définition du régime en tours par minute).
\item Un cycle complet à quatre temps correspond à 2 tours de vilebrequin.
      \[
      \text{Nombre de cycles par minute} = \frac{N}{2} = \frac{2400}{2} = \SI{1200}{cycles/min}
      \]
\item Chaque cycle complet ne comporte qu'un seul temps moteur (le temps de détente).
      \[
      \text{Nombre de temps moteur par minute} = \text{Nombre de cycles par minute} = \SI{1200}{temps moteur/min}
      \]
\end{enumerate}
\end{corrige}

\begin{corrige}[Exercice 5 — Remettre le cycle dans l'ordre]
\begin{enumerate}
\item \textbf{Ordre chronologique :} \textbf{B} $\rightarrow$ \textbf{D} $\rightarrow$ \textbf{C} $\rightarrow$ \textbf{A}.
\item \textbf{Noms des temps :}
      \begin{itemize}
      \item \textbf{B} : Temps d'\textbf{admission} (ou aspiration).
      \item \textbf{D} : Temps de \textbf{compression}.
      \item \textbf{C} : Temps de \textbf{détente} (ou explosion, ou temps moteur).
      \item \textbf{A} : Temps d'\textbf{échappement}.
      \end{itemize}
\item \textbf{Temps moteur :} C'est le temps \textbf{C} (détente).
      \par \textbf{Justification :} C'est le seul temps où la pression des gaz dans le cylindre effectue un travail moteur positif sur le piston (force $\times$ déplacement dans le même sens), ce qui entraîne la rotation du vilebrequin. Les trois autres temps (admission, compression, échappement) sont des temps « résistants » ou « passifs » qui consomment de l'énergie (fournie par l'inertie du volant moteur).
\end{enumerate}
\end{corrige}

\begin{corrige}[Exercice 6 — Légender un schéma de moteur]
\begin{enumerate}
\item \textbf{Schéma légendé :}
\begin{popfigure}
\begin{tikzpicture}[scale=0.9, every node/.style={font=\small}]
% Cylindre
\draw[thick, popdark] (0,0) -- (0,4) -- (3,4) -- (3,0);
% Culasse
\draw[thick, popdark] (-0.3,4) -- (3.3,4) -- (3.3,4.5) -- (-0.3,4.5) -- cycle;
% Piston
\draw[thick, popblue, fill=popblueL] (0.2,1.5) rectangle (2.8,2.0);
% Axe piston
\draw[thick, popdark] (1.5,1.75) -- (1.5,1.0);
% Bielle
\draw[thick, poporange, fill=poporangeL] (1.5,1.0) -- (1.5,0.2) -- (2.0,-0.3) -- (2.0,0.2) -- cycle;
% Vilebrequin (axe)
\fill[popdark] (1.5,-0.3) circle (0.15);
% Soupape Admission (gauche)
\draw[thick, popgreen] (0.7,4.5) -- (0.7,5.5);
\draw[thick, fill=popgreenL] (0.5,5.5) rectangle (0.9,5.7);
% Soupape Echappement (droite)
\draw[thick, poppink] (2.3,4.5) -- (2.3,5.5);
\draw[thick, fill=poppinkL] (2.1,5.5) rectangle (2.5,5.7);
% Bougie (milieu haut)
\draw[thick, poppurple] (1.5,4.5) -- (1.5,5.2);
\draw[thick, fill=poppurpleL] (1.3,5.2) rectangle (1.7,5.4);
% Etiquettes repères
\node[popblue, anchor=east] at (-0.3,1.75) {\textbf{1. Piston}};
\node[popdark, anchor=east] at (-0.3,3.0) {\textbf{2. Cylindre}};
\node[poppurple, anchor=west] at (3.3,5.3) {\textbf{3. Bougie}};
\node[popgreen, anchor=east] at (0.4,5.6) {\textbf{4. Soupape d'admission}};
\node[poppink, anchor=west] at (2.6,5.6) {\textbf{5. Soupape d'échappement}};
\node[poporange, anchor=north] at (1.5,0.0) {\textbf{6. Bielle}};
% Flèche compression (piston monte)
\draw[->, thick, popgold] (3.5,1.75) -- (3.5,2.75) node[midway, right] {Compression};
\end{tikzpicture}
\legende{Schéma simplifié d'un cylindre de moteur à essence (4 temps) — Coupe longitudinale. Le piston (1) est au milieu de sa course.}
\end{popfigure}

\item \textbf{Rôles de quatre organes (au choix) :}
      \begin{itemize}
      \item \textbf{Piston (1) :} Reçoit la pression des gaz, transforme le mouvement alternatif en rotation via la bielle, assure l'étanchéité (segments).
      \item \textbf{Cylindre (2) :} Guide le piston, délimite la chambre de combustion, évacue la chaleur vers le système de refroidissement.
      \item \textbf{Bougie (3) :} Produit l'étincelle électrique qui enflamme le mélange air-essence compressé (allumage commandé).
      \item \textbf{Soupape d'admission (4) :} S'ouvre pendant le temps d'admission pour laisser entrer le mélange air-carburant (ou l'air seul en diesel).
      \item \textbf{Soupape d'échappement (5) :} S'ouvre pendant le temps d'échappement pour laisser sortir les gaz brûlés.
      \item \textbf{Bielle (6) :} Relie le piston au vilebrequin, transforme le mouvement de translation du piston en mouvement de rotation du vilebrequin.
      \end{itemize}
\item \textbf{Sens du piston pendant la compression :} Vers le haut (du PMB vers le PMH). Sur le schéma, une flèche montante est indiquée à droite du cylindre.
\end{enumerate}
\end{corrige}

\begin{corrige}[Exercice 7 — Chaîne énergétique d'un moteur]
\begin{enumerate}
\item \[
      \text{Énergie chimique (carburant + air)} \longrightarrow
      \text{Énergie thermique (combustion, gaz chauds)} \longrightarrow
      \text{Énergie mécanique (rotation vilebrequin)}
      \]
\item Le rendement énergétique d'un moteur thermique est inférieur à 1 (souvent 25--35 \% pour l'essence, 35--45 \% pour le diesel). Une grande partie de l'énergie thermique issue de la combustion n'est pas convertie en travail mécanique : elle est dissipée sous forme de chaleur dans les parois du cylindre, la culasse, le piston, et évacuée par les gaz d'échappement et le système de refroidissement. C'est pour cela que le moteur chauffe.
\item \textbf{Organe :} Le système de refroidissement.
      \par \textbf{Exemple :} Refroidissement par liquide (circuit eau + radiateur + ventilateur + pompe à eau + calorstat) ou refroidissement par air (ailettes sur le cylindre + ventilateur forcé, fréquent sur motos et petits groupes).
\end{enumerate}
\end{corrige}

\begin{corrige}[Exercice 8 — Moteur à quatre cylindres]
Données : $N = \SI{3000}{tr/min}$, 4 cylindres.
\begin{enumerate}
\item Cycles par cylindre par minute :
      \[
      \frac{N}{2} = \frac{3000}{2} = \SI{1500}{cycles/min \text{ par cylindre}}
      \]
\item Chaque cycle produit 1 temps moteur par cylindre. Pour 4 cylindres :
      \[
      \text{Total temps moteur/min} = 1500 \times 4 = \SI{6000}{temps moteur/min}
      \]
      (Soit 100 temps moteur par seconde).
\item \textbf{Explication :} Dans un moteur 4 cylindres, les temps moteurs des différents cylindres sont décalés dans le temps (généralement de 180° de vilebrequin pour un 4 cylindres en ligne). Alors qu'un monocylindre ne produit un effort utile que pendant 1 temps sur 4 (1 tour sur 2), le 4 cylindres produit un effort utile en permanence (un cylindre en détente à tout moment). Cela lisse le couple moteur, supprime les à-coups, réduit les vibrations et assure un fonctionnement beaucoup plus régulier et souple.
\end{enumerate}
\end{corrige}

\begin{corrige}[Exercice 9 — Le groupe électrogène du quartier]
\begin{enumerate}
\item \textbf{Définition :} Un moteur à combustion interne est un moteur thermique où la combustion du carburant se produit à l'intérieur même de la chambre de combustion (cylindre), créant une pression qui actionne directement un piston.
      \par \textbf{Deux types étudiés :} Moteur à essence (allumage commandé) et moteur diesel (allumage par compression).
\item \textbf{Moteur à essence.}
      \par \textbf{Justification :} La présence d'une bougie d'allumage est la signature de l'allumage commandé, propre aux moteurs à essence. Le moteur diesel n'a pas de bougie (sauf bougie de préchauffage pour le démarrage à froid, mais pas d'allumage en fonctionnement).
\item \textbf{Quatre temps dans l'ordre :} Admission $\rightarrow$ Compression $\rightarrow$ Détente (Explosion) $\rightarrow$ Échappement.
\item \textbf{Gaz :} Le \textbf{monoxyde de carbone (\ce{CO})}.
      \par \textbf{Danger :} C'est un gaz incolore, inodore, non irritant, donc indétectable par les sens. Il se fixe sur l'hémoglobine du sang avec une affinité 200 à 250 fois supérieure à celle de l'oxygène, formant de la carboxyhémoglobine. Cela empêche le transport de l'\ce{O2} vers les organes (cerveau, cœur) $\rightarrow$ maux de tête, nausées, perte de connaissance, mort par asphyxie hypoxique.
\item \textbf{Consignes de sécurité (au moins deux) :}
      \begin{itemize}
      \item Placer le groupe électrogène \textbf{à l'extérieur} du bâtiment, dans un endroit bien ventilé, loin des portes/fenêtres.
      \item Ne jamais faire fonctionner un moteur thermique dans un local clos ou mal ventilé (risque d'accumulation de \ce{CO}).
      \item Installer un détecteur de monoxyde de carbone dans la boutique.
      \item Respecter les distances de sécurité par rapport aux matériaux inflammables.
      \end{itemize}
\item \textbf{Pourquoi pas de plein en marche :} L'essence est très volatile et inflammable. Les vapeurs d'essence peuvent s'enflammer au contact d'une étincelle (bougie, démarreur), d'une pièce chaude (échappement, culasse) ou d'électricité statique. Risque d'incendie ou d'explosion grave. Il faut couper le moteur et attendre son refroidissement.
\end{enumerate}

\textbf{Barème indicatif (sur 10 pts) :}
\begin{itemize}
\item Q1 : 1,5 pt (def 1 pt + 2 types 0,5 pt)
\item Q2 : 1,5 pt (réponse 0,5 pt + justif 1 pt)
\item Q3 : 1 pt (ordre exact)
\item Q4 : 2 pts (gaz 1 pt + danger mécanisme 1 pt)
\item Q5 : 2 pts (2 consignes correctes 1 pt chacune)
\item Q6 : 2 pts (risque incendie/explosion + explication vapeurs/chaleur)
\end{itemize}

\textbf{Points de vigilance :}
\begin{itemize}
\item Ne pas confondre bougie d'allumage (essence) et bougie de préchauffage (diesel).
\item Préciser que le \ce{CO} est indétectable (inodore/incolore) : c'est ce qui le rend si dangereux.
\item Pour le plein : citer explicitement « vapeurs d'essence » + « source de chaleur/étincelle ».
\end{itemize}
\end{corrige}

\begin{corrige}[Exercice 10 — La moto-taxi de Ngaoundéré]
\begin{enumerate}
\item \textbf{Cylindrée de \SI{125}{cm^3} :} Cela signifie que le volume balayé par le piston entre son point mort bas et son point mort haut (volume du cylindre utile) est de \SI{125}{cm^3}. C'est la capacité unitaire du moteur (monocylindre).
\item \textbf{Tours par seconde :}
      \[
      N = \SI{4000}{tr/min} = \frac{4000}{60} \approx \SI{66,7}{tr/s} \quad \left(\text{ou } \frac{200}{3} \text{ tr/s}\right)
      \]
\item \textbf{Temps moteur par seconde :}
      \begin{align*}
      \text{Cycles par seconde} &= \frac{N/2}{60} = \frac{2000}{60} \approx \SI{33,3}{cycles/s} \\
      \text{Temps moteur par seconde} &= \text{Cycles par seconde} \approx \SI{33,3}{temps moteur/s}
      \end{align*}
      (Car 1 cycle = 1 temps moteur pour un monocylindre).
\item \textbf{Admission :} La \textbf{soupape d'admission} est ouverte. Le mélange qui entre est le \textbf{mélange air + essence} (préparé par le carburateur ou l'injection indirecte).
\item \textbf{Origine de la chaleur :} La combustion du mélange (transformation énergie chimique $\rightarrow$ thermique) et les frottements internes (piston/cylindre, vilebrequin, distribution). Une partie seulement est convertie en travail mécanique (rendement $\approx$ 25--30 \%), le reste chauffe le moteur.
      \par \textbf{Conséquence d'une surchauffe :} Dilatation excessive du piston $\rightarrow$ grippage (blocage) dans le cylindre ; déformation de la culasse (perte d'étanchéité joint de culasse) ; dégradation de l'huile (perte lubrification) $\rightarrow$ casse moteur.
\item \textbf{Chaîne énergétique :}
      \[
      \text{Énergie chimique} \longrightarrow \text{Énergie thermique} \longrightarrow \text{Énergie mécanique}
      \]
\end{enumerate}

\textbf{Barème indicatif (sur 10 pts) :}
\begin{itemize}
\item Q1 : 1 pt
\item Q2 : 1 pt (conversion min $\rightarrow$ s)
\item Q3 : 2 pts (calcul cycles + lien 1 cycle = 1 temps moteur)
\item Q4 : 1,5 pt (soupape 0,5 pt + mélange 1 pt)
\item Q5 : 2,5 pts (origine 1 pt + conséquence 1,5 pt)
\item Q6 : 1 pt
\item Présentation/Unités : 1 pt
\end{itemize}

\textbf{Points de vigilance :}
\begin{itemize}
\item Attention à la conversion tr/min $\rightarrow$ tr/s (division par 60).
\item Ne pas oublier que 1 cycle = 2 tours = 1 temps moteur.
\item Préciser « mélange air-essence » pour l'admission (pas « essence seule » ni « air seul »).
\end{itemize}
\end{corrige}

\begin{corrige}[Exercice 11 — Essence ou gasoil à la SONARA ?]
\begin{enumerate}
\item \textbf{Tableau comparatif complété :}
\begin{center}
\begin{tabularx}{\linewidth}{|l|X|X|}
\hline
\rowcolor{popblueL} \textbf{Critère} & \textbf{Moteur à essence} & \textbf{Moteur diesel} \\
\hline
Carburant utilisé & Essence (super, sans plomb) & Gasoil (gazole) \\
\hline
Organe d'allumage & Bougie d'allumage & Aucun (injecteur seulement) \\
\hline
Mode d'inflammation du mélange & Allumage commandé par étincelle électrique & Allumage par compression (auto-inflammation) \\
\hline
\end{tabularx}
\end{center}

\item \textbf{L'air seul} est comprimé dans le cylindre pendant le temps de compression. Le gasoil n'est injecté qu'à la fin de ce temps (au PMH).
\item Le taux de compression (rapport volumes PMB/PMH) est plus élevé en diesel (14:1 à 22:1) qu'en essence (8:1 à 12:1) pour deux raisons :
      \begin{enumerate}
      \item Atteindre une température de l'air comprimé supérieure à la température d'auto-inflammation du gasoil ($\approx \SI{400}{\celsius}$) sans source d'allumage externe.
      \item Améliorer le rendement thermodynamique (cycle Diesel plus efficace à haut taux de compression). En essence, un taux trop élevé provoquerait le cliquetis (auto-inflammation anarchique avant l'étincelle).
      \end{enumerate}
\item \textbf{Conseil : Camion B (Diesel).}
      \par \textbf{Arguments :}
      \begin{enumerate}
      \item \textbf{Couple et puissance à bas régime :} Le diesel développe un couple maximal à bas régime, idéal pour démarrer et tracter de lourdes charges en côte sans rétrograder sans cesse.
      \item \textbf{Consommation et coût :} Rendement supérieur (environ 40 \% vs 30 \%) $\rightarrow$ consommation plus faible (L/100 km). Le gasoil est souvent moins cher à la pompe (SONARA) que l'essence. Autonomie plus grande pour le même volume de réservoir.
      \end{enumerate}
      (Autres arguments acceptés : longévité mécanique supérieure, pas de système d'allumage à entretenir).
\item \textbf{Deux mesures d'entretien régulier :}
      \begin{itemize}
      \item Vidange de l'huile moteur + remplacement du filtre à huile (tous les 5000--10000 km ou selon constructeur).
      \item Contrôle/remplacement du filtre à air (propreté admission).
      \item Contrôle niveau liquide de refroidissement + concentration antigel.
      \item Vérification état/courroie de distribution (risque casse moteur).
      \item Pour essence : changement bougies. Pour diesel : purge filtre gasoil / contrôle injecteurs.
      \end{itemize}
\item \textbf{Gaz polluant :} \textbf{Dioxyde de carbone (\ce{CO2})} (gaz à effet de serre, réchauffement climatique) \textbf{OU} \textbf{Oxydes d'azote (\ce{NOx})} / \textbf{Particules fines (diesel)} / \textbf{Monoxyde de carbone (\ce{CO})} / \textbf{Hydrocarbures imbrûlés}.
      \par \textbf{Mesure de réduction :} Pot catalytique (3 voies essence, oxydation diesel) ; Filtre à particules (FAP) pour diesel ; Normes Euro successives (injection common rail, EGR, AdBlue/SCR) ; Entretien rigoureux ; Éco-conduite ; Renouvellement parc vers véhicules moins polluants (électrique, gaz).
\end{enumerate}

\textbf{Barème indicatif (sur 12 pts) :}
\begin{itemize}
\item Q1 : 3 pts (3 lignes $\times$ 1 pt)
\item Q2 : 1 pt
\item Q3 : 2 pts (T° auto-inflammation + rendement/cliquetis)
\item Q4 : 3 pts (Choix 0,5 pt + 2 arguments justifiés 1,25 pt chacun)
\item Q5 : 1,5 pts (2 mesures distinctes)
\item Q6 : 1,5 pts (Gaz 0,5 pt + Mesure 1 pt)
\end{itemize}

\textbf{Points de vigilance :}
\begin{itemize}
\item Ne pas écrire « injection » comme organe d'allumage pour le diesel (l'injecteur injecte, il n'allume pas).
\item Justifier le taux de compression par la physique (température) et non par « c'est plus costaud ».
\item Pour le choix du camion, l'argument « gasoil moins cher » seul ne suffit pas, il faut le coupler au rendement/consommation.
\item Citer un gaz précis (\ce{CO2}, \ce{NOx}, particules) et une mesure technique ou organisationnelle concrète.
\end{itemize}
\end{corrige}

\newpage

\coursec{Fiche bilan}

\begin{popfigure}
\centering
\begin{tikzpicture}[mindmap, grow cyclic, every node/.style={concept, text=white, font=\small\bfseries},
root concept/.append style={concept color=popdark, font=\bfseries, minimum size=3.2cm},
level 1/.append style={level distance=3.6cm, sibling angle=60},
level 1 concept/.append style={minimum size=2.5cm}]
\node[root concept]{Moteur à combustion interne}
child[concept color=popblue]{ node {Définition et principe} }
child[concept color=poporange]{ node {Éléments et rôles} }
child[concept color=popgreen]{ node {Cycle à 4 temps} }
child[concept color=poppurple]{ node {Essence ou Diesel ?} }
child[concept color=popgold]{ node {Sécurité et entretien} }
child[concept color=poppink]{ node {Environnement} };
\end{tikzpicture}
\legende{Carte mentale de la leçon : les six grandes idées à retenir sur le moteur à combustion interne.}
\end{popfigure}

\begin{bilan}
\textbf{Définitions clés}
\begin{itemize}
\item Un \cle{moteur à combustion interne} est une machine thermique qui transforme l'énergie chimique d'un carburant (essence, gasoil) en \cle{énergie mécanique}, la combustion ayant lieu \emph{à l'intérieur} du moteur (dans le cylindre).
\item Le \cle{cylindre} est le tube dans lequel coulisse le \cle{piston}. Le \cle{cylindrée} est le volume total balayé par les pistons.
\item La \cle{bougie} produit l'étincelle qui enflamme le mélange (moteur à essence). Les \cle{soupapes} (d'admission et d'échappement) laissent entrer les gaz frais et sortir les gaz brûlés.
\item Le système \cle{bielle-manivelle} (bielle + vilebrequin) transforme le mouvement alternatif du piston en \cle{mouvement de rotation}.
\item Un \cle{temps} est une course complète du piston (montée ou descente). Un cycle à \cle{quatre temps} correspond à 4 courses du piston, soit 2 tours de vilebrequin.
\end{itemize}

\textbf{Les éléments du moteur et leurs rôles}
\begin{center}
\begin{tabularx}{\linewidth}{|l|X|}
\hline
\rowcolor{popblueL} \textbf{Élément} & \textbf{Rôle} \\
\hline
Cylindre & Chambre où se déroule la combustion ; guide le piston. \\
\hline
Piston & Reçoit la poussée des gaz brûlés et la transmet à la bielle. \\
\hline
Bielle & Relie le piston au vilebrequin. \\
\hline
Vilebrequin (manivelle) & Transforme la translation du piston en rotation. \\
\hline
Soupapes & Admission : entrée des gaz frais ; échappement : sortie des gaz brûlés. \\
\hline
Bougie (essence) & Produit l'étincelle d'allumage. \\
\hline
Injecteur (diesel) & Pulvérise le gasoil dans l'air comprimé chaud. \\
\hline
\end{tabularx}
\end{center}

\textbf{Le cycle à quatre temps (à connaître par cœur)}
\begin{enumerate}
\item \cle{Admission} : le piston descend, la soupape d'admission est ouverte, le mélange (ou l'air) entre.
\item \cle{Compression} : soupapes fermées, le piston remonte et comprime le mélange.
\item \cle{Combustion -- détente} (temps moteur) : la bougie enflamme le mélange (ou le gasoil s'enflamme spontanément) ; les gaz poussent le piston vers le bas.
\item \cle{Échappement} : la soupape d'échappement s'ouvre, le piston remonte et chasse les gaz brûlés.
\end{enumerate}

\textbf{Moteur à essence / moteur diesel : les différences}
\begin{center}
\begin{tabularx}{\linewidth}{|l|X|X|}
\hline
\rowcolor{popblueL} \textbf{Critère} & \textbf{Moteur à essence} & \textbf{Moteur diesel} \\
\hline
Carburant & Essence & Gasoil \\
\hline
Gaz admis & Mélange air + essence & Air seul \\
\hline
Allumage & Par étincelle (bougie) & Par compression (auto-inflammation) \\
\hline
Pièce d'allumage & Bougie & Injecteur \\
\hline
Exemples & Moto-taxi, groupe électrogène & Camions, gros groupes \\
\hline
\end{tabularx}
\end{center}

\textbf{Méthode express pour reconnaître un temps moteur sur un schéma}
\begin{enumerate}
\item Regarder le sens de déplacement du piston (flèche).
\item Vérifier l'état des soupapes (ouverte ou fermée).
\item Conclure : descente + admission ouverte = admission ; montée + tout fermé = compression ; descente + tout fermé (+ étincelle) = détente ; montée + échappement ouvert = échappement.
\end{enumerate}

\textbf{Sécurité et entretien}
\begin{itemize}
\item Ne jamais faire tourner un moteur dans un local fermé : le \cle{monoxyde de carbone} \ce{CO} est un gaz toxique et mortel.
\item Ne jamais approcher une flamme d'un réservoir d'essence ; couper le moteur pour faire le plein.
\item Entretien : vidange régulière de l'huile, nettoyage du filtre à air, vérification des bougies.
\item Les gaz d'échappement (\ce{CO2}, \ce{CO}, oxydes d'azote) polluent l'air : un moteur bien réglé pollue moins.
\end{itemize}
\end{bilan}

\begin{aretenir}[Checklist avant le BEPC]
$\square$ Je sais définir un moteur à combustion interne.\\
$\square$ Je connais la chaîne de conversion d'énergie : chimique $\rightarrow$ thermique $\rightarrow$ mécanique.\\
$\square$ Je sais citer les éléments du moteur : cylindre, piston, bielle, vilebrequin, soupapes, bougie.\\
$\square$ Je sais donner le rôle de chaque élément.\\
$\square$ Je connais les quatre temps dans l'ordre : admission, compression, combustion-détente, échappement.\\
$\square$ Je sais que le temps moteur est le 3\textsuperscript{e} temps (combustion-détente).\\
$\square$ Je sais qu'un cycle à 4 temps = 4 courses de piston = 2 tours de vilebrequin.\\
$\square$ Je sais identifier un temps moteur sur un schéma (position du piston, état des soupapes).\\
$\square$ Je sais distinguer moteur à essence (bougie, mélange air-essence) et moteur diesel (injecteur, air seul, allumage par compression).\\
$\square$ Je sais expliquer pourquoi le monoxyde de carbone \ce{CO} est dangereux.\\
$\square$ Je connais les consignes de sécurité liées au carburant et à l'échappement.\\
$\square$ Je sais citer deux gestes d'entretien qui réduisent la pollution.
\end{aretenir}

\findecours

\end{document}
